% GENERATED FILE -- DO NOT EDIT.
%
% Produced by harness/assemble_contract.py from module.yaml and claims/*.tex.
% Edit the claim files and regenerate; `--check` fails if this file is stale.
\documentclass[11pt,a4paper]{article}
\pdftrailerid{7be292a621b11dfc217b4c36955e9d54}
%% --- 1. Mathematics (order matters: physics after braket wins \ket, \tr) ---
\usepackage{amsmath}
\usepackage{amssymb}
\usepackage{amsfonts}
\usepackage{amsthm}
\usepackage{mathtools}
\usepackage{dsfont}
\usepackage{braket}
\usepackage{physics}
%% --- 2. Page furniture ------------------------------------------------------
\usepackage{xcolor}
\usepackage{array}    % >{...} column prefixes in the wide tables of some claims
\usepackage{enumitem}
\usepackage{graphicx}
\usepackage{booktabs}
\usepackage{longtable}
\usepackage{xltabular}              % front- and back-matter tables across pages
\usepackage{needspace}
\usepackage{xurl}
\usepackage{tcolorbox}
\tcbuselibrary{skins,breakable}
%% --- 3. Links and navigation (hyperref late; bookmark, refcount after it) ---
\usepackage{xr-hyper}   % before hyperref: the contract's labels
\usepackage{hyperref}
\usepackage{bookmark}
\usepackage{refcount}
%% --- 4. Colours: every distinction also has a shape (rule, caps, size, weight)
\definecolor{nnGrey}{gray}{0.40}      % non-normative text
\definecolor{metaGrey}{gray}{0.30}    % the generated meta line
\definecolor{techGrey}{gray}{0.45}    % heading of a technical claim
\definecolor{mainRule}{RGB}{0,110,140}
\definecolor{reqAccent}{RGB}{165,70,0}
\definecolor{linkBlue}{RGB}{0,62,150}
\definecolor{extLink}{RGB}{120,40,140}  % links into another PDF
\hypersetup{colorlinks=true, linkcolor=linkBlue, citecolor=linkBlue,
  urlcolor=linkBlue, filecolor=extLink, bookmarksnumbered=true,
  bookmarksopen=true, bookmarksopenlevel=1, pdfdisplaydoctitle=true}
% A link into another PDF (\href to a .pdf, \hyperref{file}..., xr-hyper) is a
% relative URI, <file>#nameddest=<destination>, in the file-link colour:
% browser viewers (Edge, Chrome) ignore hyperref's default GoToR action.
\makeatletter
\renewcommand*{\hyper@linkfile}[3]{% anchor text, file name, destination
  \begingroup
    \ifx\\#3\\\def\CQ@target{#2}\else\def\CQ@target{#2\hyper@hash nameddest=#3}\fi
    \let\@urlcolor\@filecolor
    \hyper@linkurl{#1}{\CQ@target}%
  \endgroup}
\makeatother
%% --- 5. The mapped Unicode set (harness/contract_navigation.py, UNICODE_TEX)
\DeclareUnicodeCharacter{03B1}{\ensuremath{\alpha}}
\DeclareUnicodeCharacter{03B2}{\ensuremath{\beta}}
\DeclareUnicodeCharacter{03B3}{\ensuremath{\gamma}}
\DeclareUnicodeCharacter{03B4}{\ensuremath{\delta}}
\DeclareUnicodeCharacter{03B5}{\ensuremath{\varepsilon}}
\DeclareUnicodeCharacter{03F5}{\ensuremath{\epsilon}}
\DeclareUnicodeCharacter{03B6}{\ensuremath{\zeta}}
\DeclareUnicodeCharacter{03B7}{\ensuremath{\eta}}
\DeclareUnicodeCharacter{03B8}{\ensuremath{\theta}}
\DeclareUnicodeCharacter{03D1}{\ensuremath{\vartheta}}
\DeclareUnicodeCharacter{03B9}{\ensuremath{\iota}}
\DeclareUnicodeCharacter{03BA}{\ensuremath{\kappa}}
\DeclareUnicodeCharacter{03BB}{\ensuremath{\lambda}}
\DeclareUnicodeCharacter{03BC}{\ensuremath{\mu}}
\DeclareUnicodeCharacter{03BD}{\ensuremath{\nu}}
\DeclareUnicodeCharacter{03BE}{\ensuremath{\xi}}
\DeclareUnicodeCharacter{03C0}{\ensuremath{\pi}}
\DeclareUnicodeCharacter{03C1}{\ensuremath{\rho}}
\DeclareUnicodeCharacter{03F1}{\ensuremath{\varrho}}
\DeclareUnicodeCharacter{03C3}{\ensuremath{\sigma}}
\DeclareUnicodeCharacter{03C2}{\ensuremath{\varsigma}}
\DeclareUnicodeCharacter{03C4}{\ensuremath{\tau}}
\DeclareUnicodeCharacter{03C5}{\ensuremath{\upsilon}}
\DeclareUnicodeCharacter{03C6}{\ensuremath{\varphi}}
\DeclareUnicodeCharacter{03D5}{\ensuremath{\phi}}
\DeclareUnicodeCharacter{03C7}{\ensuremath{\chi}}
\DeclareUnicodeCharacter{03C8}{\ensuremath{\psi}}
\DeclareUnicodeCharacter{03C9}{\ensuremath{\omega}}
\DeclareUnicodeCharacter{0393}{\ensuremath{\Gamma}}
\DeclareUnicodeCharacter{0394}{\ensuremath{\Delta}}
\DeclareUnicodeCharacter{0398}{\ensuremath{\Theta}}
\DeclareUnicodeCharacter{039B}{\ensuremath{\Lambda}}
\DeclareUnicodeCharacter{039E}{\ensuremath{\Xi}}
\DeclareUnicodeCharacter{03A0}{\ensuremath{\Pi}}
\DeclareUnicodeCharacter{03A3}{\ensuremath{\Sigma}}
\DeclareUnicodeCharacter{03A5}{\ensuremath{\Upsilon}}
\DeclareUnicodeCharacter{03A6}{\ensuremath{\Phi}}
\DeclareUnicodeCharacter{03A8}{\ensuremath{\Psi}}
\DeclareUnicodeCharacter{03A9}{\ensuremath{\Omega}}
\DeclareUnicodeCharacter{00E1}{\'a}
\DeclareUnicodeCharacter{00E0}{\`a}
\DeclareUnicodeCharacter{00E2}{\^a}
\DeclareUnicodeCharacter{00E4}{\"a}
\DeclareUnicodeCharacter{00E5}{\aa{}}
\DeclareUnicodeCharacter{00E9}{\'e}
\DeclareUnicodeCharacter{00E8}{\`e}
\DeclareUnicodeCharacter{00EA}{\^e}
\DeclareUnicodeCharacter{00EB}{\"e}
\DeclareUnicodeCharacter{00ED}{\'{\i}}
\DeclareUnicodeCharacter{00EF}{\"{\i}}
\DeclareUnicodeCharacter{00F3}{\'o}
\DeclareUnicodeCharacter{00F2}{\`o}
\DeclareUnicodeCharacter{00F4}{\^o}
\DeclareUnicodeCharacter{00F6}{\"o}
\DeclareUnicodeCharacter{00F8}{\o{}}
\DeclareUnicodeCharacter{00FA}{\'u}
\DeclareUnicodeCharacter{00F9}{\`u}
\DeclareUnicodeCharacter{00FB}{\^u}
\DeclareUnicodeCharacter{00FC}{\"u}
\DeclareUnicodeCharacter{00F1}{\~n}
\DeclareUnicodeCharacter{00E7}{\c{c}}
\DeclareUnicodeCharacter{00DF}{\ss{}}
\DeclareUnicodeCharacter{00C4}{\"A}
\DeclareUnicodeCharacter{00C5}{\AA{}}
\DeclareUnicodeCharacter{00C9}{\'E}
\DeclareUnicodeCharacter{00D6}{\"O}
\DeclareUnicodeCharacter{00DC}{\"U}
\DeclareUnicodeCharacter{2013}{\textendash{}}
\DeclareUnicodeCharacter{2014}{\textemdash{}}
\DeclareUnicodeCharacter{2018}{`}
\DeclareUnicodeCharacter{2019}{'}
\DeclareUnicodeCharacter{201C}{``}
\DeclareUnicodeCharacter{201D}{''}
\DeclareUnicodeCharacter{2026}{\ldots{}}
\DeclareUnicodeCharacter{00A0}{~}
\DeclareUnicodeCharacter{00A7}{\S{}}
\DeclareUnicodeCharacter{2264}{\ensuremath{\le}}
\DeclareUnicodeCharacter{2265}{\ensuremath{\ge}}
\DeclareUnicodeCharacter{2260}{\ensuremath{\neq}}
\DeclareUnicodeCharacter{2248}{\ensuremath{\approx}}
\DeclareUnicodeCharacter{2261}{\ensuremath{\equiv}}
\DeclareUnicodeCharacter{223C}{\ensuremath{\sim}}
\DeclareUnicodeCharacter{2243}{\ensuremath{\simeq}}
\DeclareUnicodeCharacter{2245}{\ensuremath{\cong}}
\DeclareUnicodeCharacter{221D}{\ensuremath{\propto}}
\DeclareUnicodeCharacter{226A}{\ensuremath{\ll}}
\DeclareUnicodeCharacter{226B}{\ensuremath{\gg}}
\DeclareUnicodeCharacter{2192}{\ensuremath{\to}}
\DeclareUnicodeCharacter{2190}{\ensuremath{\leftarrow}}
\DeclareUnicodeCharacter{2194}{\ensuremath{\leftrightarrow}}
\DeclareUnicodeCharacter{21D2}{\ensuremath{\Rightarrow}}
\DeclareUnicodeCharacter{21D0}{\ensuremath{\Leftarrow}}
\DeclareUnicodeCharacter{21D4}{\ensuremath{\Leftrightarrow}}
\DeclareUnicodeCharacter{21A6}{\ensuremath{\mapsto}}
\DeclareUnicodeCharacter{2208}{\ensuremath{\in}}
\DeclareUnicodeCharacter{2209}{\ensuremath{\notin}}
\DeclareUnicodeCharacter{2286}{\ensuremath{\subseteq}}
\DeclareUnicodeCharacter{2282}{\ensuremath{\subset}}
\DeclareUnicodeCharacter{2287}{\ensuremath{\supseteq}}
\DeclareUnicodeCharacter{2283}{\ensuremath{\supset}}
\DeclareUnicodeCharacter{2229}{\ensuremath{\cap}}
\DeclareUnicodeCharacter{222A}{\ensuremath{\cup}}
\DeclareUnicodeCharacter{2216}{\ensuremath{\setminus}}
\DeclareUnicodeCharacter{2205}{\ensuremath{\emptyset}}
\DeclareUnicodeCharacter{2200}{\ensuremath{\forall}}
\DeclareUnicodeCharacter{2203}{\ensuremath{\exists}}
\DeclareUnicodeCharacter{2115}{\ensuremath{\mathbb{N}}}
\DeclareUnicodeCharacter{2124}{\ensuremath{\mathbb{Z}}}
\DeclareUnicodeCharacter{211A}{\ensuremath{\mathbb{Q}}}
\DeclareUnicodeCharacter{211D}{\ensuremath{\mathbb{R}}}
\DeclareUnicodeCharacter{2102}{\ensuremath{\mathbb{C}}}
\DeclareUnicodeCharacter{00B7}{\ensuremath{\cdot}}
\DeclareUnicodeCharacter{00D7}{\ensuremath{\times}}
\DeclareUnicodeCharacter{2212}{\ensuremath{-}}
\DeclareUnicodeCharacter{00B1}{\ensuremath{\pm}}
\DeclareUnicodeCharacter{2213}{\ensuremath{\mp}}
\DeclareUnicodeCharacter{221A}{\ensuremath{\surd}}
\DeclareUnicodeCharacter{2211}{\ensuremath{\sum}}
\DeclareUnicodeCharacter{220F}{\ensuremath{\prod}}
\DeclareUnicodeCharacter{222B}{\ensuremath{\int}}
\DeclareUnicodeCharacter{221E}{\ensuremath{\infty}}
\DeclareUnicodeCharacter{2202}{\ensuremath{\partial}}
\DeclareUnicodeCharacter{2207}{\ensuremath{\nabla}}
\DeclareUnicodeCharacter{2297}{\ensuremath{\otimes}}
\DeclareUnicodeCharacter{2295}{\ensuremath{\oplus}}
\DeclareUnicodeCharacter{2218}{\ensuremath{\circ}}
\DeclareUnicodeCharacter{2308}{\ensuremath{\lceil}}
\DeclareUnicodeCharacter{2309}{\ensuremath{\rceil}}
\DeclareUnicodeCharacter{230A}{\ensuremath{\lfloor}}
\DeclareUnicodeCharacter{230B}{\ensuremath{\rfloor}}
\DeclareUnicodeCharacter{27E8}{\ensuremath{\langle}}
\DeclareUnicodeCharacter{27E9}{\ensuremath{\rangle}}
\DeclareUnicodeCharacter{2016}{\ensuremath{\|}}
\DeclareUnicodeCharacter{2113}{\ensuremath{\ell}}
\DeclareUnicodeCharacter{00B0}{\ensuremath{^{\circ}}}
\DeclareUnicodeCharacter{00B2}{\ensuremath{^{2}}}
\DeclareUnicodeCharacter{00B3}{\ensuremath{^{3}}}
\DeclareUnicodeCharacter{00B9}{\ensuremath{^{1}}}
\DeclareUnicodeCharacter{2080}{\ensuremath{_{0}}}
\DeclareUnicodeCharacter{2081}{\ensuremath{_{1}}}
\DeclareUnicodeCharacter{2082}{\ensuremath{_{2}}}
\DeclareUnicodeCharacter{207F}{\ensuremath{^{n}}}
\DeclareUnicodeCharacter{2096}{\ensuremath{_{k}}}
%% --- 6. Numbering: one counter per section -----------------------
\numberwithin{equation}{section}
\setcounter{tocdepth}{1}
%% --- 7. The project macro set, then the claim environments -----------------
% spec/macros.tex, read relative to this module's directory (latexmk -cd).
% =============================================================================
%  spec/macros.tex -- the one macro set every generated contract inputs.
%
%  Two parts: part A, the macros spec/notation.md catalogues, and part B,
%  the macros local to one module, grouped by module.
%  harness/assemble_contract.py emits % =============================================================================
%  spec/macros.tex -- the one macro set every generated contract inputs.
%
%  Two parts: part A, the macros spec/notation.md catalogues, and part B,
%  the macros local to one module, grouped by module.
%  harness/assemble_contract.py emits % =============================================================================
%  spec/macros.tex -- the one macro set every generated contract inputs.
%
%  Two parts: part A, the macros spec/notation.md catalogues, and part B,
%  the macros local to one module, grouped by module.
%  harness/assemble_contract.py emits \input{../../macros} for it, so the
%  file is read from spec/modules/<id>/.
%
%  RULES
%    * A macro that spec/notation.md catalogues is a verbatim copy of it
%      (part A). Changing one is a notation change: bump spec/notation.md.
%    * Part B holds the module-local macros. A macro defined by two modules
%      with the same expansion appears once, under the first module (in the
%      order below) that defines it. Two macros with the same name and
%      DIFFERENT expansions are both kept under a module prefix, and the
%      claims write the prefixed names:
%        \Chain  -> \redChain (reductions, \mathfrak{C}: the reduction chain)
%                   \pcChain  (protocol-compiler, \mathfrak{R}: the two-segment
%                             chain; \Chsrc and \Chrec are built on it)
%        \Nmsm   -> \redNmsm  (reductions, \mathcal{N}: a POVM measurement
%                             channel)
%                   \pcNmsm   (\mathcal{N}_{\mathsf A}: Alice's measurement
%                             channel, written by
%                             transmitter/thm:certificate-sufficiency,
%                             transmitter/rem:xreg and
%                             reductions/lem:source-transfer)
%    * Deliberately omitted: a \provenance{tag}{details} macro. Provenance
%      is a header field of every claim file (spec/claim-format.md section
%      3); a claim body that calls such a macro fails to compile, which is
%      the intended signal.
%
%  REQUIRED PACKAGES, loaded by the assembler's preamble before this file:
%    amsmath, amssymb, amsfonts, mathtools, dsfont, braket, physics, xcolor
%    (braket before physics, so that physics wins for \ket, \bra, \tr, ...).
%
%  This file contains only definitions: no \usepackage, no document body.
% =============================================================================


% =============================================================================
%  PART A -- notation macros, the ones spec/notation.md catalogues, each a
%  verbatim copy of its catalogue row; grouped under numbered headings.
% =============================================================================

% -----------------------------------------------------------------------------
% 1. State spaces and operator sets (a "system label" -> spaces of operators)
% -----------------------------------------------------------------------------
\newcommand{\Hs}[1]{\mathcal{H}(#1)}        % Hilbert space of a system label
\newcommand{\Lin}[1]{\mathcal{L}(#1)}       % bounded linear operators on H(#1)
\newcommand{\Linn}[2]{\mathcal{L}(#1,#2)}   % linear maps H(#1) -> H(#2)
\newcommand{\Pos}[1]{\mathrm{Pos}(#1)}      % positive semidefinite operators on H(#1)
\newcommand{\Density}{\mathrm{D}}           % density operators: \Density(A) = D(A)
\newcommand{\Seq}[1]{S_{=}(#1)}             % normalized states (trace = 1)
\newcommand{\Sleq}[1]{S_{\leq}(#1)}         % subnormalized states (0 <= trace <= 1)
\newcommand{\id}{\mathds{1}}                % identity operator
\newcommand{\idon}[1]{\mathds{1}_{#1}}      % identity on a named system, e.g. \idon{C}

% DEPRECATED: applies \mathcal to both arguments itself, so it takes bare
% letters, not system labels. Prefer \Linn; do not use in new contracts.
\providecommand{\Lop}[2]{\mathcal{L}(\mathcal{#1}, \mathcal{#2})}


% -----------------------------------------------------------------------------
% 2. Composite systems
% -----------------------------------------------------------------------------
\newcommand{\dsum}{\oplus}                  % direct sum of system labels
% Tensor product: use \otimes directly.
% n-fold direct sum of an operator space: \Lin{A}^{\oplus n}


% -----------------------------------------------------------------------------
% 3. Events, conditioning and classical registers (for STATES)
%    Convention:  \wedge  -> subnormalized "and-event" component, e.g. \rho_{A\wedge x}
%                 |        -> normalized conditional state,        e.g. \rho_{A|x}
% (Write these directly in subscripts; no extra macros needed.)
% -----------------------------------------------------------------------------


% -----------------------------------------------------------------------------
% 4. Photon-number subspaces / direct-sum blocks
%    Convention:  operators carry the restriction "on top", in brackets;
%                 states carry it as a \wedge subscript (see Sec. 3).
% -----------------------------------------------------------------------------
\newcommand{\Pnum}[1]{P_{#1}}               % photon-number projector (system from context)
\newcommand{\Pnumon}[2]{P_{#1}^{#2}}        % with explicit system superscript, e.g. P_n^C
\newcommand{\blk}[2]{#1^{(#2)}}             % operator block, e.g. \blk{M}{\leqN} = M^{(<=N)}
\newcommand{\leqN}{\leq N}                  % use as \blk{M}{\leqN}
\newcommand{\gtN}{> N}                      % use as \blk{M}{\gtN}
\newcommand{\eqn}{= n}                      % use as \blk{M}{\eqn}


% -----------------------------------------------------------------------------
% 5. System labels used in the protocol
% -----------------------------------------------------------------------------
\newcommand{\Chan}{C}                       % quantum-channel (Eve) input to Bob
\newcommand{\Prb}{P}                        % probe input
\newcommand{\Flag}{{\mathrm{Flg}}}          % flag system (EXISTS ONLY AFTER squashing)
\newcommand{\Loss}{{\mathrm{Lss}}}          % loss system
\newcommand{\Vc}{{\mathrm{Vcm}}}            % vacuum system
\newcommand{\Nout}{N_{\mathrm{out}}}
\newcommand{\Nin}{N_{\mathrm{in}}}


% -----------------------------------------------------------------------------
% 6. Alphabets
% -----------------------------------------------------------------------------
\newcommand{\Yset}{\mathcal{Y}}             % Bob's (coarse-grained) detection alphabet
\newcommand{\Yraw}{\mathcal{Y}_{\mathrm{raw}}}% raw click patterns {0,1}^4
\newcommand{\Yscset}{\mathcal{Y}_{\singlec}}% single-click outcomes
\newcommand{\Ymcset}{\mathcal{Y}_{\mc}}     % multi-click outcomes
\newcommand{\Tset}{\mathcal{T}_{\mathrm{set}}}% private setting set (the probe settings in the optical instance)
\newcommand{\Xset}{\mathcal{X}}             % Alice's alphabet
\newcommand{\Dset}{\{H,V,+,-\}}             % polarization / detector labels
\newcommand{\pass}{\top}                    % pass sifting


% -----------------------------------------------------------------------------
% 7. Mode / outcome abbreviations
% -----------------------------------------------------------------------------
\newcommand{\on}{{\mathrm{on}}}
\newcommand{\off}{{\mathrm{off}}}
\newcommand{\dc}{{\mathrm{dc}}}             % dark count
\newcommand{\sq}{\mathrm{sq}}               % squashed
\newcommand{\mc}{\mathrm{mc}}               % multi-click
\newcommand{\nc}{\mathrm{nc}}               % no-click
\newcommand{\singlec}{\mathrm{sc}}          % single-click
\newcommand{\vac}{\mathrm{vac}}             % vacuum
\newcommand{\kout}{{k_\mathrm{out}}}
\newcommand{\kin}{{k_\mathrm{in}}}


% -----------------------------------------------------------------------------
% 8. Maps, channels and miscellaneous operators
% -----------------------------------------------------------------------------
\newcommand{\Prob}{\mathrm{Prob}}           % probability (use \Prob[\cdot] or \Prob(\cdot))
\newcommand{\lmin}{\lambda_{\min}}          % minimum eigenvalue
\newcommand{\Wflag}{W_{\Flag}}              % flag-space weight
\newcommand{\transpose}{\mathsf{T}}
\newcommand{\innerp}[2]{\left\langle #1 \vert #2 \right\rangle}
\renewcommand{\vec}[1]{\boldsymbol{#1}}     % NB: overrides the LaTeX arrow accent


% -----------------------------------------------------------------------------
% 9. Key-rate numerics
% -----------------------------------------------------------------------------
\newcommand{\cG}{\mathcal{G}}               % protocol map
\newcommand{\cZ}{\mathcal{Z}}               % pinching map
\newcommand{\Amap}{\mathcal{A}_{\mathrm{ann}}} % announcement + measurement map
\newcommand{\Amapnb}{\mathcal{A}_{\mathrm{ann}|\mathrm{nb}}}
\newcommand{\simpl}{}                       % empty suffix for simplified block maps
\newcommand{\acSet}{\mathcal{S}_{\mathrm{acc}}} % acceptance set
\newcommand{\testE}{\mathcal{E}_{\mathrm{test}}} % test-statistics map
\newcommand{\ecLeak}{\mathrm{leak}^{\mathrm{EC}}_{\mathrm{obs}}}
\newcommand{\pgen}{p_{\mathrm{gen}}}
\newcommand{\askeyr}{R_{\infty}}             % asymptotic key rate
\newcommand{\ppass}{p_{\mathrm{pass}}}       % sifting pass probability


% -----------------------------------------------------------------------------
% 10. Finite-size analysis / Renyi entropies
% -----------------------------------------------------------------------------

% --- Round types: test vs. key generation (register T, alphabet \Fset) --------
\newcommand{\tst}{\mathrm{test}}            % test round
\newcommand{\gen}{\mathrm{gen}}             % key-generation round
\newcommand{\Fset}{\mathcal{T}}             % round-type alphabet {test, gen} (not \Tset)
\newcommand{\Chat}{\widehat{C}}             % acceptance-test register
\newcommand{\Chatset}{\widehat{\mathcal{C}}}% its alphabet

% --- Basis / key extraction functions ----------------------------------------
\newcommand{\bA}{b_{A}}                     % Alice's basis-extraction function
\newcommand{\bB}{b_{B}}                     % Bob's basis-extraction function
\newcommand{\kB}{k_{B}}                     % Bob's key-extraction function

% --- Maps --------------------------------------------------------------------
\newcommand{\Mprot}{\mathcal{M}}            % single-round protocol map
\newcommand{\Phimap}{\boldsymbol{\Phi}}     % test-statistics map
\newcommand{\Pdeph}{\mathcal{P}}            % dephasing (pinching) map on \Blreg

% --- Isometries --------------------------------------------------------------
\newcommand{\Vmsm}{V_{\mathrm{msm}}}        % measurement isometry
\newcommand{\Vann}{V_{\mathrm{ann}}}        % announcement isometry
\newcommand{\VI}{V_{I}}                     % Eve's copy of the announcements
\newcommand{\VChat}{V_{\Chat}}              % acceptance-data isometry
\newcommand{\VR}{V_{R}}                     % sifting + key map -> raw key R
\newcommand{\VZ}{V_{Z}}                     % raw key -> key register Z
\newcommand{\Wprot}{W}                      % full protocol isometry

% --- Renyi entropies ---------------------------------------------------------
% SECURITY-CRITICAL: see spec/notation.md §4. Never change an arrow to
% make an inequality work out; never convert silently between the sandwiched
% (tilde) and Petz (overline) families.
\newcommand{\Hsandup}{\widetilde{H}^{\uparrow}_{\alpha}}   % sandwiched, up
\newcommand{\Hsanddown}{\widetilde{H}^{\downarrow}_{\alpha}}% sandwiched, down
\newcommand{\Hpetzup}{\overline{H}^{\uparrow}_{\alpha}}    % Petz, up
\newcommand{\Hpetzdown}{\overline{H}^{\downarrow}_{\alpha}}% Petz, down
\newcommand{\hup}{h^{\uparrow}_{\alpha}}                   % single-round rate
\newcommand{\hqkd}{h^{\uparrow,\mathrm{QKD}}_{\alpha}}     % single-round QKD rate
\newcommand{\hlow}{\underline{h}_{\alpha}}                 % lower bound on \hqkd
% --- spec/notation.md 3.5 and 4.5 -------------------------------------------
% \hlb  is the first-order linearization lower bound on \hlow: the headline
%       key rate (the `first-order` tier of foundations/def:tiers). NOT \hlow
%       itself.
% \tnum is the instance's one declared numerical tolerance. Its
%       near-collision with the declared marginal \tau_A is disclosed in
%       spec/notation.md 3.5; bare \tau is reserved and denotes neither.
\newcommand{\hlb}{\underline{h}^{\mathrm{FO}}_{\alpha}}    % first-order lower bound on \hlow
\newcommand{\tnum}{\tau_{\mathrm{num}}}                    % declared numerical tolerance


% =============================================================================
%  PART B -- module-local macros, grouped by module. Deprecated aliases,
%  defined so that a text written with them compiles, used in no new contract: \Lop,
%  \testE, \Pdeph, \pass, \ecLeak (spec/notation.md section 3.9).
% =============================================================================

% --- meat-core -----------------------------------------------------------------
\newcommand{\SecSet}{\mathbb{S}}            % sector index set of Q
\newcommand{\Qsys}{Q}                       % declared sector-carrying measured system
\newcommand{\Vset}{\mathcal{V}}             % declared joint outcome alphabet of the Q-POVM
\newcommand{\CAset}{\mathcal{C}_{A}}        % Alice announcement alphabet
\newcommand{\CBset}{\mathcal{C}_{B}}        % Bob announcement alphabet
\newcommand{\Party}{\mathfrak{P}}           % declared party partition
\newcommand{\Oset}[1]{\mathcal{O}_{#1}}     % readout alphabet of party #1
\newcommand{\rd}[1]{\pi_{#1}}               % readout map of party #1
\newcommand{\Sadm}{\mathcal{S}_{A}}         % admissible Alice subspace
\newcommand{\nomarg}{\mathsf{none}}         % "no trusted marginal declared"
\newcommand{\Zset}{\mathcal{Z}_{\perp}}     % raw-key alphabet incl. \perp
\newcommand{\Zval}{\mathcal{Z}}             % key alphabet without \perp
\newcommand{\Etil}{\widetilde{E}}           % single-round purifying register
\newcommand{\Ehat}{\widehat{E}}             % Eve's n-round register
\newcommand{\OmAT}{\Omega_{\mathrm{AT}}}    % acceptance event
\newcommand{\epa}{\varepsilon_{\mathrm{PA}}}
\newcommand{\eev}{\varepsilon_{\mathrm{EV}}}
\newcommand{\lEC}{\lambda_{\mathrm{EC}}}
\newcommand{\Decl}{\mathfrak{D}}            % protocol declaration
\newcommand{\CList}{\mathfrak{L}}           % declared constraint list
\newcommand{\PinS}{\mathcal{P}_{\SecSet}}   % sector pinching on Q
\newcommand{\Feas}{\mathcal{W}}             % feasible set of block states
\newcommand{\Chkey}{\Chatset_{\mathrm{key}}}% key-live acceptance values
\newcommand{\Kann}{\mathcal{K}}             % announcement-class alphabet
\newcommand{\hold}{\mathsf{h}}              % declared key holder (A or B)
\newcommand{\sfD}{D}                        % classical KL divergence (base 2)
\newcommand{\Dtwo}{D}                       % quantum relative entropy (base 2)
\newcommand{\Dsand}{\widetilde{D}_{\alpha}} % sandwiched Renyi divergence
\newcommand{\idmap}{\mathrm{id}}            % identity channel
\newcommand{\Aev}{\mathcal{A}}              % Eve's per-round attack channel
\newcommand{\AttSet}{\mathfrak{A}}          % licensed attack-channel set
\newcommand{\Chacc}{\Chatset_{\mathrm{acc}}}% acceptance support of \acSet
\newcommand{\epsstar}{\varepsilon^{\star}}      % distance to an exact POVM
\newcommand{\dfH}{\Delta_{\mathrm{herm}}}       % Hermiticity defect
\newcommand{\dfP}{\Delta_{\mathrm{pos}}}        % negative part
\newcommand{\dfC}{\Delta_{\mathrm{cpl}}}        % completeness defect
\newcommand{\lmoq}{g}                       % acceptance-side gradient block
\newcommand{\lmoJ}{G}                       % feasible-side gradient block
\newcommand{\bacc}{\beta_{\mathrm{acc}}}    % certified acceptance-side minimum
\newcommand{\Itset}{\mathfrak{I}}           % finite index set of iterates (C22)

% --- reductions ----------------------------------------------------------------
\newcommand{\Bsys}{B}                        % physical (pre-squashing) system
\newcommand{\Tcl}[1]{\mathcal{T}(#1)}        % trace-class operators on H(#1)
\newcommand{\Ins}{\mathcal{I}}               % generalized-instrument branch map
\newcommand{\redNmsm}{\mathcal{N}}              % POVM measurement channel
\newcommand{\Eprot}{\mathcal{E}}             % n-round protocol channel
\newcommand{\Grnd}{\mathsf{G}}               % one round map of a protocol
\newcommand{\Kfin}{\mathsf{K}}               % end-of-protocol processing
\newcommand{\Crit}{\mathsf{sec}}             % value of a security criterion
\newcommand{\Pflag}{\Pi_{\Flag}}             % flag-sector projector on Q
\newcommand{\Pleq}{\Pi_{\leqN}}              % preserved-block projector
\newcommand{\Pgt}{\Pi_{\gtN}}                % above-cutoff projector
\newcommand{\Sph}{S^{\mathrm{ph}}}           % family of physical inputs
\newcommand{\Srch}{S^{\mathrm{rch}}}         % reachable squashed states
\newcommand{\Ssq}{S^{\sq}}                   % squashed admissible set
\newcommand{\Ffss}{F^{\mathrm{FSS}}}         % flag-state-squasher target
\newcommand{\Fwp}{F^{\mathrm{WP}}}           % multi-bound WPFSS target
\newcommand{\Mwp}{M^{\mathrm{WP}}}           % renormalized above-cutoff POVM
\newcommand{\Mdec}{{M'}^{\mathrm{ph}}}       % declared physical family on B
\newcommand{\Ppvm}{P^{\Flag}}                % flag PVM
\newcommand{\Lmp}{{\Lambda_{\mathrm{mp}}}}   % measure-and-prepare channel
\newcommand{\Lwp}{{\Lambda^{\mathrm{WP}}}}   % multi-bound WPFSS channel
\newcommand{\Lfss}{{\Lambda^{\mathrm{FSS}}}} % flag-state-squasher channel
\newcommand{\lmax}{\lambda_{\max}}           % maximum eigenvalue
\newcommand{\Lnoise}{\widetilde{\Lambda}}    % flag-noise channel
\newcommand{\cg}{\varpi}                     % declared coarse-graining map
\newcommand{\MCK}{\mathsf{MC}}               % multiclick raw-pattern set
\newcommand{\redChain}{\mathfrak{C}}            % reduction chain
\newcommand{\Lnk}{\ell}                      % a single typed link
\newcommand{\dsq}{\preceq}                      % the squashing relation
\newcommand{\opn}[1]{\left\lVert #1\right\rVert}           % operator (spectral) norm
\newcommand{\mxn}[1]{\left\lVert #1\right\rVert_{\max}}    % entrywise max modulus
\newcommand{\Fld}{\mathbb{F}}                   % exact ordered substrate field
\newcommand{\Fldi}{\mathbb{F}_{i}}              % its Gaussian extension F + iF
\newcommand{\lat}[1]{\big\lfloor #1\big\rfloor_{\kappa}} % round DOWN onto the lattice

% --- protocol-compiler ---------------------------------------------------------
\newcommand{\Spec}{\mathfrak{S}}             % protocol specification (compiler input)
\newcommand{\pcChain}{\mathfrak{R}}            % reduction chain (two typed segments)
\newcommand{\Chsrc}{\pcChain_{\mathrm{src}}}   % its Alice-side source-map segment
\newcommand{\Chrec}{\pcChain_{\mathrm{rec}}}   % its receiver-side squashing segment
\newcommand{\rvirt}{\rho^{\mathrm{v}}}       % virtual (PRE-source-map) signal state
\newcommand{\Fph}{F^{\mathrm{ph}}}           % physical (pre-reduction) POVM element
\newcommand{\Vph}{\mathcal{V}^{\mathrm{ph}}} % physical outcome alphabet
\newcommand{\Tpub}{\mathcal{T}_{\mathrm{pub}}}% public-choice alphabet
\newcommand{\Uset}{\mathcal{U}}              % residual outcome alphabet
\newcommand{\idxm}{\mathsf{idx}}             % label-to-index map
\newcommand{\Asig}{A'}                       % transmitted (signal) system -- the REAL one
\newcommand{\Avirt}{A''}                     % virtual transmitted system (source descriptor of Xi_1)
\newcommand{\Abar}{\bar{A}}                  % Alice's measured register
\newcommand{\Ahat}{\hat{A}}                  % Alice's shield register
\newcommand{\pcNmsm}{\mathcal{N}_{\mathsf A}}  % Alice's measurement channel (see the RULES block)
\newcommand{\Phip}{\Phi^{+}}                 % maximally entangled vector label
\newcommand{\Smap}{\mathfrak{s}}             % the emitted sector map (fraktur: a compile-time object)

% --- receiver ------------------------------------------------------------------
\newcommand{\Modes}{\mathbb{M}}              % declared ordered mode set
\newcommand{\Inp}{\mathbb{I}}                % declared input modes
\newcommand{\InpC}{\mathbb{I}_{\Chan}}       % channel input modes
\newcommand{\InpS}{\mathbb{I}_{\Prb}}        % probe (source) input modes
\newcommand{\Env}{\mathbb{E}}                % dilation (environment) modes
\newcommand{\Port}[1]{\mathbb{P}_{#1}}       % mode set of detector port #1
\newcommand{\Emd}[1]{E_{#1}}                 % coordinate projector on modes #1
\newcommand{\Gmat}{G}                        % instrument (transfer) matrix
\newcommand{\Fock}{\mathfrak{F}}             % Fock / second-quantization functor
\newcommand{\Symn}[1]{\mathrm{Sym}^{#1}}     % symmetric tensor power
\newcommand{\normord}[1]{{:}\mspace{2mu}#1\mspace{2mu}{:}}   % normal ordering
\newcommand{\rk}{\operatorname{rank}}        % rank (physics.sty owns \rank)
\newcommand{\SCK}{\mathsf{SC}}               % single-click pattern set
\newcommand{\Recv}{\mathfrak{R}}             % receiver instrument tuple
\newcommand{\Cat}{\mathfrak{C}}              % component catalogue
\newcommand{\Grph}{\mathfrak{E}}             % declared element graph
\newcommand{\defic}{\mathbb{L}}              % loss deficiency operator
\newcommand{\Pvac}[1]{\Pi_{\vac}(#1)}        % joint no-click element on ports #1
\newcommand{\pmc}{p_{\mc}}                   % uniform high-photon mc weight
\newcommand{\Bnd}{\beta}                     % the scalar bound function beta(n)
\newcommand{\Rot}{R}                         % planar rotation matrix
\newcommand{\Wplate}{W}                      % retarder (wave-plate) matrix
\newcommand{\Bsp}{B}                         % beamsplitter core matrix

% --- statistics ----------------------------------------------------------------
\newcommand{\honest}{p^{\mathrm{hon}}}          % honest per-round distribution
\newcommand{\ecomp}{\varepsilon_{\mathrm{comp}}}% completeness budget
\newcommand{\simplex}[1]{\Delta(#1)}            % probability simplex
\newcommand{\Chatcon}{\Chatset_{\mathrm{con}}}  % constrained outcomes
\newcommand{\Ireg}[3]{\mathcal{I}_{#1}(#2,#3)}  % regularized incomplete beta
\newcommand{\Ncount}[1]{\mathsf{N}_{#1}}        % per-outcome count
\newcommand{\qlo}{\mathbf{q}^{-}}               % box lower bounds
\newcommand{\qhi}{\mathbf{q}^{+}}               % box upper bounds
 for it, so the
%  file is read from spec/modules/<id>/.
%
%  RULES
%    * A macro that spec/notation.md catalogues is a verbatim copy of it
%      (part A). Changing one is a notation change: bump spec/notation.md.
%    * Part B holds the module-local macros. A macro defined by two modules
%      with the same expansion appears once, under the first module (in the
%      order below) that defines it. Two macros with the same name and
%      DIFFERENT expansions are both kept under a module prefix, and the
%      claims write the prefixed names:
%        \Chain  -> \redChain (reductions, \mathfrak{C}: the reduction chain)
%                   \pcChain  (protocol-compiler, \mathfrak{R}: the two-segment
%                             chain; \Chsrc and \Chrec are built on it)
%        \Nmsm   -> \redNmsm  (reductions, \mathcal{N}: a POVM measurement
%                             channel)
%                   \pcNmsm   (\mathcal{N}_{\mathsf A}: Alice's measurement
%                             channel, written by
%                             transmitter/thm:certificate-sufficiency,
%                             transmitter/rem:xreg and
%                             reductions/lem:source-transfer)
%    * Deliberately omitted: a \provenance{tag}{details} macro. Provenance
%      is a header field of every claim file (spec/claim-format.md section
%      3); a claim body that calls such a macro fails to compile, which is
%      the intended signal.
%
%  REQUIRED PACKAGES, loaded by the assembler's preamble before this file:
%    amsmath, amssymb, amsfonts, mathtools, dsfont, braket, physics, xcolor
%    (braket before physics, so that physics wins for \ket, \bra, \tr, ...).
%
%  This file contains only definitions: no \usepackage, no document body.
% =============================================================================


% =============================================================================
%  PART A -- notation macros, the ones spec/notation.md catalogues, each a
%  verbatim copy of its catalogue row; grouped under numbered headings.
% =============================================================================

% -----------------------------------------------------------------------------
% 1. State spaces and operator sets (a "system label" -> spaces of operators)
% -----------------------------------------------------------------------------
\newcommand{\Hs}[1]{\mathcal{H}(#1)}        % Hilbert space of a system label
\newcommand{\Lin}[1]{\mathcal{L}(#1)}       % bounded linear operators on H(#1)
\newcommand{\Linn}[2]{\mathcal{L}(#1,#2)}   % linear maps H(#1) -> H(#2)
\newcommand{\Pos}[1]{\mathrm{Pos}(#1)}      % positive semidefinite operators on H(#1)
\newcommand{\Density}{\mathrm{D}}           % density operators: \Density(A) = D(A)
\newcommand{\Seq}[1]{S_{=}(#1)}             % normalized states (trace = 1)
\newcommand{\Sleq}[1]{S_{\leq}(#1)}         % subnormalized states (0 <= trace <= 1)
\newcommand{\id}{\mathds{1}}                % identity operator
\newcommand{\idon}[1]{\mathds{1}_{#1}}      % identity on a named system, e.g. \idon{C}

% DEPRECATED: applies \mathcal to both arguments itself, so it takes bare
% letters, not system labels. Prefer \Linn; do not use in new contracts.
\providecommand{\Lop}[2]{\mathcal{L}(\mathcal{#1}, \mathcal{#2})}


% -----------------------------------------------------------------------------
% 2. Composite systems
% -----------------------------------------------------------------------------
\newcommand{\dsum}{\oplus}                  % direct sum of system labels
% Tensor product: use \otimes directly.
% n-fold direct sum of an operator space: \Lin{A}^{\oplus n}


% -----------------------------------------------------------------------------
% 3. Events, conditioning and classical registers (for STATES)
%    Convention:  \wedge  -> subnormalized "and-event" component, e.g. \rho_{A\wedge x}
%                 |        -> normalized conditional state,        e.g. \rho_{A|x}
% (Write these directly in subscripts; no extra macros needed.)
% -----------------------------------------------------------------------------


% -----------------------------------------------------------------------------
% 4. Photon-number subspaces / direct-sum blocks
%    Convention:  operators carry the restriction "on top", in brackets;
%                 states carry it as a \wedge subscript (see Sec. 3).
% -----------------------------------------------------------------------------
\newcommand{\Pnum}[1]{P_{#1}}               % photon-number projector (system from context)
\newcommand{\Pnumon}[2]{P_{#1}^{#2}}        % with explicit system superscript, e.g. P_n^C
\newcommand{\blk}[2]{#1^{(#2)}}             % operator block, e.g. \blk{M}{\leqN} = M^{(<=N)}
\newcommand{\leqN}{\leq N}                  % use as \blk{M}{\leqN}
\newcommand{\gtN}{> N}                      % use as \blk{M}{\gtN}
\newcommand{\eqn}{= n}                      % use as \blk{M}{\eqn}


% -----------------------------------------------------------------------------
% 5. System labels used in the protocol
% -----------------------------------------------------------------------------
\newcommand{\Chan}{C}                       % quantum-channel (Eve) input to Bob
\newcommand{\Prb}{P}                        % probe input
\newcommand{\Flag}{{\mathrm{Flg}}}          % flag system (EXISTS ONLY AFTER squashing)
\newcommand{\Loss}{{\mathrm{Lss}}}          % loss system
\newcommand{\Vc}{{\mathrm{Vcm}}}            % vacuum system
\newcommand{\Nout}{N_{\mathrm{out}}}
\newcommand{\Nin}{N_{\mathrm{in}}}


% -----------------------------------------------------------------------------
% 6. Alphabets
% -----------------------------------------------------------------------------
\newcommand{\Yset}{\mathcal{Y}}             % Bob's (coarse-grained) detection alphabet
\newcommand{\Yraw}{\mathcal{Y}_{\mathrm{raw}}}% raw click patterns {0,1}^4
\newcommand{\Yscset}{\mathcal{Y}_{\singlec}}% single-click outcomes
\newcommand{\Ymcset}{\mathcal{Y}_{\mc}}     % multi-click outcomes
\newcommand{\Tset}{\mathcal{T}_{\mathrm{set}}}% private setting set (the probe settings in the optical instance)
\newcommand{\Xset}{\mathcal{X}}             % Alice's alphabet
\newcommand{\Dset}{\{H,V,+,-\}}             % polarization / detector labels
\newcommand{\pass}{\top}                    % pass sifting


% -----------------------------------------------------------------------------
% 7. Mode / outcome abbreviations
% -----------------------------------------------------------------------------
\newcommand{\on}{{\mathrm{on}}}
\newcommand{\off}{{\mathrm{off}}}
\newcommand{\dc}{{\mathrm{dc}}}             % dark count
\newcommand{\sq}{\mathrm{sq}}               % squashed
\newcommand{\mc}{\mathrm{mc}}               % multi-click
\newcommand{\nc}{\mathrm{nc}}               % no-click
\newcommand{\singlec}{\mathrm{sc}}          % single-click
\newcommand{\vac}{\mathrm{vac}}             % vacuum
\newcommand{\kout}{{k_\mathrm{out}}}
\newcommand{\kin}{{k_\mathrm{in}}}


% -----------------------------------------------------------------------------
% 8. Maps, channels and miscellaneous operators
% -----------------------------------------------------------------------------
\newcommand{\Prob}{\mathrm{Prob}}           % probability (use \Prob[\cdot] or \Prob(\cdot))
\newcommand{\lmin}{\lambda_{\min}}          % minimum eigenvalue
\newcommand{\Wflag}{W_{\Flag}}              % flag-space weight
\newcommand{\transpose}{\mathsf{T}}
\newcommand{\innerp}[2]{\left\langle #1 \vert #2 \right\rangle}
\renewcommand{\vec}[1]{\boldsymbol{#1}}     % NB: overrides the LaTeX arrow accent


% -----------------------------------------------------------------------------
% 9. Key-rate numerics
% -----------------------------------------------------------------------------
\newcommand{\cG}{\mathcal{G}}               % protocol map
\newcommand{\cZ}{\mathcal{Z}}               % pinching map
\newcommand{\Amap}{\mathcal{A}_{\mathrm{ann}}} % announcement + measurement map
\newcommand{\Amapnb}{\mathcal{A}_{\mathrm{ann}|\mathrm{nb}}}
\newcommand{\simpl}{}                       % empty suffix for simplified block maps
\newcommand{\acSet}{\mathcal{S}_{\mathrm{acc}}} % acceptance set
\newcommand{\testE}{\mathcal{E}_{\mathrm{test}}} % test-statistics map
\newcommand{\ecLeak}{\mathrm{leak}^{\mathrm{EC}}_{\mathrm{obs}}}
\newcommand{\pgen}{p_{\mathrm{gen}}}
\newcommand{\askeyr}{R_{\infty}}             % asymptotic key rate
\newcommand{\ppass}{p_{\mathrm{pass}}}       % sifting pass probability


% -----------------------------------------------------------------------------
% 10. Finite-size analysis / Renyi entropies
% -----------------------------------------------------------------------------

% --- Round types: test vs. key generation (register T, alphabet \Fset) --------
\newcommand{\tst}{\mathrm{test}}            % test round
\newcommand{\gen}{\mathrm{gen}}             % key-generation round
\newcommand{\Fset}{\mathcal{T}}             % round-type alphabet {test, gen} (not \Tset)
\newcommand{\Chat}{\widehat{C}}             % acceptance-test register
\newcommand{\Chatset}{\widehat{\mathcal{C}}}% its alphabet

% --- Basis / key extraction functions ----------------------------------------
\newcommand{\bA}{b_{A}}                     % Alice's basis-extraction function
\newcommand{\bB}{b_{B}}                     % Bob's basis-extraction function
\newcommand{\kB}{k_{B}}                     % Bob's key-extraction function

% --- Maps --------------------------------------------------------------------
\newcommand{\Mprot}{\mathcal{M}}            % single-round protocol map
\newcommand{\Phimap}{\boldsymbol{\Phi}}     % test-statistics map
\newcommand{\Pdeph}{\mathcal{P}}            % dephasing (pinching) map on \Blreg

% --- Isometries --------------------------------------------------------------
\newcommand{\Vmsm}{V_{\mathrm{msm}}}        % measurement isometry
\newcommand{\Vann}{V_{\mathrm{ann}}}        % announcement isometry
\newcommand{\VI}{V_{I}}                     % Eve's copy of the announcements
\newcommand{\VChat}{V_{\Chat}}              % acceptance-data isometry
\newcommand{\VR}{V_{R}}                     % sifting + key map -> raw key R
\newcommand{\VZ}{V_{Z}}                     % raw key -> key register Z
\newcommand{\Wprot}{W}                      % full protocol isometry

% --- Renyi entropies ---------------------------------------------------------
% SECURITY-CRITICAL: see spec/notation.md §4. Never change an arrow to
% make an inequality work out; never convert silently between the sandwiched
% (tilde) and Petz (overline) families.
\newcommand{\Hsandup}{\widetilde{H}^{\uparrow}_{\alpha}}   % sandwiched, up
\newcommand{\Hsanddown}{\widetilde{H}^{\downarrow}_{\alpha}}% sandwiched, down
\newcommand{\Hpetzup}{\overline{H}^{\uparrow}_{\alpha}}    % Petz, up
\newcommand{\Hpetzdown}{\overline{H}^{\downarrow}_{\alpha}}% Petz, down
\newcommand{\hup}{h^{\uparrow}_{\alpha}}                   % single-round rate
\newcommand{\hqkd}{h^{\uparrow,\mathrm{QKD}}_{\alpha}}     % single-round QKD rate
\newcommand{\hlow}{\underline{h}_{\alpha}}                 % lower bound on \hqkd
% --- spec/notation.md 3.5 and 4.5 -------------------------------------------
% \hlb  is the first-order linearization lower bound on \hlow: the headline
%       key rate (the `first-order` tier of foundations/def:tiers). NOT \hlow
%       itself.
% \tnum is the instance's one declared numerical tolerance. Its
%       near-collision with the declared marginal \tau_A is disclosed in
%       spec/notation.md 3.5; bare \tau is reserved and denotes neither.
\newcommand{\hlb}{\underline{h}^{\mathrm{FO}}_{\alpha}}    % first-order lower bound on \hlow
\newcommand{\tnum}{\tau_{\mathrm{num}}}                    % declared numerical tolerance


% =============================================================================
%  PART B -- module-local macros, grouped by module. Deprecated aliases,
%  defined so that a text written with them compiles, used in no new contract: \Lop,
%  \testE, \Pdeph, \pass, \ecLeak (spec/notation.md section 3.9).
% =============================================================================

% --- meat-core -----------------------------------------------------------------
\newcommand{\SecSet}{\mathbb{S}}            % sector index set of Q
\newcommand{\Qsys}{Q}                       % declared sector-carrying measured system
\newcommand{\Vset}{\mathcal{V}}             % declared joint outcome alphabet of the Q-POVM
\newcommand{\CAset}{\mathcal{C}_{A}}        % Alice announcement alphabet
\newcommand{\CBset}{\mathcal{C}_{B}}        % Bob announcement alphabet
\newcommand{\Party}{\mathfrak{P}}           % declared party partition
\newcommand{\Oset}[1]{\mathcal{O}_{#1}}     % readout alphabet of party #1
\newcommand{\rd}[1]{\pi_{#1}}               % readout map of party #1
\newcommand{\Sadm}{\mathcal{S}_{A}}         % admissible Alice subspace
\newcommand{\nomarg}{\mathsf{none}}         % "no trusted marginal declared"
\newcommand{\Zset}{\mathcal{Z}_{\perp}}     % raw-key alphabet incl. \perp
\newcommand{\Zval}{\mathcal{Z}}             % key alphabet without \perp
\newcommand{\Etil}{\widetilde{E}}           % single-round purifying register
\newcommand{\Ehat}{\widehat{E}}             % Eve's n-round register
\newcommand{\OmAT}{\Omega_{\mathrm{AT}}}    % acceptance event
\newcommand{\epa}{\varepsilon_{\mathrm{PA}}}
\newcommand{\eev}{\varepsilon_{\mathrm{EV}}}
\newcommand{\lEC}{\lambda_{\mathrm{EC}}}
\newcommand{\Decl}{\mathfrak{D}}            % protocol declaration
\newcommand{\CList}{\mathfrak{L}}           % declared constraint list
\newcommand{\PinS}{\mathcal{P}_{\SecSet}}   % sector pinching on Q
\newcommand{\Feas}{\mathcal{W}}             % feasible set of block states
\newcommand{\Chkey}{\Chatset_{\mathrm{key}}}% key-live acceptance values
\newcommand{\Kann}{\mathcal{K}}             % announcement-class alphabet
\newcommand{\hold}{\mathsf{h}}              % declared key holder (A or B)
\newcommand{\sfD}{D}                        % classical KL divergence (base 2)
\newcommand{\Dtwo}{D}                       % quantum relative entropy (base 2)
\newcommand{\Dsand}{\widetilde{D}_{\alpha}} % sandwiched Renyi divergence
\newcommand{\idmap}{\mathrm{id}}            % identity channel
\newcommand{\Aev}{\mathcal{A}}              % Eve's per-round attack channel
\newcommand{\AttSet}{\mathfrak{A}}          % licensed attack-channel set
\newcommand{\Chacc}{\Chatset_{\mathrm{acc}}}% acceptance support of \acSet
\newcommand{\epsstar}{\varepsilon^{\star}}      % distance to an exact POVM
\newcommand{\dfH}{\Delta_{\mathrm{herm}}}       % Hermiticity defect
\newcommand{\dfP}{\Delta_{\mathrm{pos}}}        % negative part
\newcommand{\dfC}{\Delta_{\mathrm{cpl}}}        % completeness defect
\newcommand{\lmoq}{g}                       % acceptance-side gradient block
\newcommand{\lmoJ}{G}                       % feasible-side gradient block
\newcommand{\bacc}{\beta_{\mathrm{acc}}}    % certified acceptance-side minimum
\newcommand{\Itset}{\mathfrak{I}}           % finite index set of iterates (C22)

% --- reductions ----------------------------------------------------------------
\newcommand{\Bsys}{B}                        % physical (pre-squashing) system
\newcommand{\Tcl}[1]{\mathcal{T}(#1)}        % trace-class operators on H(#1)
\newcommand{\Ins}{\mathcal{I}}               % generalized-instrument branch map
\newcommand{\redNmsm}{\mathcal{N}}              % POVM measurement channel
\newcommand{\Eprot}{\mathcal{E}}             % n-round protocol channel
\newcommand{\Grnd}{\mathsf{G}}               % one round map of a protocol
\newcommand{\Kfin}{\mathsf{K}}               % end-of-protocol processing
\newcommand{\Crit}{\mathsf{sec}}             % value of a security criterion
\newcommand{\Pflag}{\Pi_{\Flag}}             % flag-sector projector on Q
\newcommand{\Pleq}{\Pi_{\leqN}}              % preserved-block projector
\newcommand{\Pgt}{\Pi_{\gtN}}                % above-cutoff projector
\newcommand{\Sph}{S^{\mathrm{ph}}}           % family of physical inputs
\newcommand{\Srch}{S^{\mathrm{rch}}}         % reachable squashed states
\newcommand{\Ssq}{S^{\sq}}                   % squashed admissible set
\newcommand{\Ffss}{F^{\mathrm{FSS}}}         % flag-state-squasher target
\newcommand{\Fwp}{F^{\mathrm{WP}}}           % multi-bound WPFSS target
\newcommand{\Mwp}{M^{\mathrm{WP}}}           % renormalized above-cutoff POVM
\newcommand{\Mdec}{{M'}^{\mathrm{ph}}}       % declared physical family on B
\newcommand{\Ppvm}{P^{\Flag}}                % flag PVM
\newcommand{\Lmp}{{\Lambda_{\mathrm{mp}}}}   % measure-and-prepare channel
\newcommand{\Lwp}{{\Lambda^{\mathrm{WP}}}}   % multi-bound WPFSS channel
\newcommand{\Lfss}{{\Lambda^{\mathrm{FSS}}}} % flag-state-squasher channel
\newcommand{\lmax}{\lambda_{\max}}           % maximum eigenvalue
\newcommand{\Lnoise}{\widetilde{\Lambda}}    % flag-noise channel
\newcommand{\cg}{\varpi}                     % declared coarse-graining map
\newcommand{\MCK}{\mathsf{MC}}               % multiclick raw-pattern set
\newcommand{\redChain}{\mathfrak{C}}            % reduction chain
\newcommand{\Lnk}{\ell}                      % a single typed link
\newcommand{\dsq}{\preceq}                      % the squashing relation
\newcommand{\opn}[1]{\left\lVert #1\right\rVert}           % operator (spectral) norm
\newcommand{\mxn}[1]{\left\lVert #1\right\rVert_{\max}}    % entrywise max modulus
\newcommand{\Fld}{\mathbb{F}}                   % exact ordered substrate field
\newcommand{\Fldi}{\mathbb{F}_{i}}              % its Gaussian extension F + iF
\newcommand{\lat}[1]{\big\lfloor #1\big\rfloor_{\kappa}} % round DOWN onto the lattice

% --- protocol-compiler ---------------------------------------------------------
\newcommand{\Spec}{\mathfrak{S}}             % protocol specification (compiler input)
\newcommand{\pcChain}{\mathfrak{R}}            % reduction chain (two typed segments)
\newcommand{\Chsrc}{\pcChain_{\mathrm{src}}}   % its Alice-side source-map segment
\newcommand{\Chrec}{\pcChain_{\mathrm{rec}}}   % its receiver-side squashing segment
\newcommand{\rvirt}{\rho^{\mathrm{v}}}       % virtual (PRE-source-map) signal state
\newcommand{\Fph}{F^{\mathrm{ph}}}           % physical (pre-reduction) POVM element
\newcommand{\Vph}{\mathcal{V}^{\mathrm{ph}}} % physical outcome alphabet
\newcommand{\Tpub}{\mathcal{T}_{\mathrm{pub}}}% public-choice alphabet
\newcommand{\Uset}{\mathcal{U}}              % residual outcome alphabet
\newcommand{\idxm}{\mathsf{idx}}             % label-to-index map
\newcommand{\Asig}{A'}                       % transmitted (signal) system -- the REAL one
\newcommand{\Avirt}{A''}                     % virtual transmitted system (source descriptor of Xi_1)
\newcommand{\Abar}{\bar{A}}                  % Alice's measured register
\newcommand{\Ahat}{\hat{A}}                  % Alice's shield register
\newcommand{\pcNmsm}{\mathcal{N}_{\mathsf A}}  % Alice's measurement channel (see the RULES block)
\newcommand{\Phip}{\Phi^{+}}                 % maximally entangled vector label
\newcommand{\Smap}{\mathfrak{s}}             % the emitted sector map (fraktur: a compile-time object)

% --- receiver ------------------------------------------------------------------
\newcommand{\Modes}{\mathbb{M}}              % declared ordered mode set
\newcommand{\Inp}{\mathbb{I}}                % declared input modes
\newcommand{\InpC}{\mathbb{I}_{\Chan}}       % channel input modes
\newcommand{\InpS}{\mathbb{I}_{\Prb}}        % probe (source) input modes
\newcommand{\Env}{\mathbb{E}}                % dilation (environment) modes
\newcommand{\Port}[1]{\mathbb{P}_{#1}}       % mode set of detector port #1
\newcommand{\Emd}[1]{E_{#1}}                 % coordinate projector on modes #1
\newcommand{\Gmat}{G}                        % instrument (transfer) matrix
\newcommand{\Fock}{\mathfrak{F}}             % Fock / second-quantization functor
\newcommand{\Symn}[1]{\mathrm{Sym}^{#1}}     % symmetric tensor power
\newcommand{\normord}[1]{{:}\mspace{2mu}#1\mspace{2mu}{:}}   % normal ordering
\newcommand{\rk}{\operatorname{rank}}        % rank (physics.sty owns \rank)
\newcommand{\SCK}{\mathsf{SC}}               % single-click pattern set
\newcommand{\Recv}{\mathfrak{R}}             % receiver instrument tuple
\newcommand{\Cat}{\mathfrak{C}}              % component catalogue
\newcommand{\Grph}{\mathfrak{E}}             % declared element graph
\newcommand{\defic}{\mathbb{L}}              % loss deficiency operator
\newcommand{\Pvac}[1]{\Pi_{\vac}(#1)}        % joint no-click element on ports #1
\newcommand{\pmc}{p_{\mc}}                   % uniform high-photon mc weight
\newcommand{\Bnd}{\beta}                     % the scalar bound function beta(n)
\newcommand{\Rot}{R}                         % planar rotation matrix
\newcommand{\Wplate}{W}                      % retarder (wave-plate) matrix
\newcommand{\Bsp}{B}                         % beamsplitter core matrix

% --- statistics ----------------------------------------------------------------
\newcommand{\honest}{p^{\mathrm{hon}}}          % honest per-round distribution
\newcommand{\ecomp}{\varepsilon_{\mathrm{comp}}}% completeness budget
\newcommand{\simplex}[1]{\Delta(#1)}            % probability simplex
\newcommand{\Chatcon}{\Chatset_{\mathrm{con}}}  % constrained outcomes
\newcommand{\Ireg}[3]{\mathcal{I}_{#1}(#2,#3)}  % regularized incomplete beta
\newcommand{\Ncount}[1]{\mathsf{N}_{#1}}        % per-outcome count
\newcommand{\qlo}{\mathbf{q}^{-}}               % box lower bounds
\newcommand{\qhi}{\mathbf{q}^{+}}               % box upper bounds
 for it, so the
%  file is read from spec/modules/<id>/.
%
%  RULES
%    * A macro that spec/notation.md catalogues is a verbatim copy of it
%      (part A). Changing one is a notation change: bump spec/notation.md.
%    * Part B holds the module-local macros. A macro defined by two modules
%      with the same expansion appears once, under the first module (in the
%      order below) that defines it. Two macros with the same name and
%      DIFFERENT expansions are both kept under a module prefix, and the
%      claims write the prefixed names:
%        \Chain  -> \redChain (reductions, \mathfrak{C}: the reduction chain)
%                   \pcChain  (protocol-compiler, \mathfrak{R}: the two-segment
%                             chain; \Chsrc and \Chrec are built on it)
%        \Nmsm   -> \redNmsm  (reductions, \mathcal{N}: a POVM measurement
%                             channel)
%                   \pcNmsm   (\mathcal{N}_{\mathsf A}: Alice's measurement
%                             channel, written by
%                             transmitter/thm:certificate-sufficiency,
%                             transmitter/rem:xreg and
%                             reductions/lem:source-transfer)
%    * Deliberately omitted: a \provenance{tag}{details} macro. Provenance
%      is a header field of every claim file (spec/claim-format.md section
%      3); a claim body that calls such a macro fails to compile, which is
%      the intended signal.
%
%  REQUIRED PACKAGES, loaded by the assembler's preamble before this file:
%    amsmath, amssymb, amsfonts, mathtools, dsfont, braket, physics, xcolor
%    (braket before physics, so that physics wins for \ket, \bra, \tr, ...).
%
%  This file contains only definitions: no \usepackage, no document body.
% =============================================================================


% =============================================================================
%  PART A -- notation macros, the ones spec/notation.md catalogues, each a
%  verbatim copy of its catalogue row; grouped under numbered headings.
% =============================================================================

% -----------------------------------------------------------------------------
% 1. State spaces and operator sets (a "system label" -> spaces of operators)
% -----------------------------------------------------------------------------
\newcommand{\Hs}[1]{\mathcal{H}(#1)}        % Hilbert space of a system label
\newcommand{\Lin}[1]{\mathcal{L}(#1)}       % bounded linear operators on H(#1)
\newcommand{\Linn}[2]{\mathcal{L}(#1,#2)}   % linear maps H(#1) -> H(#2)
\newcommand{\Pos}[1]{\mathrm{Pos}(#1)}      % positive semidefinite operators on H(#1)
\newcommand{\Density}{\mathrm{D}}           % density operators: \Density(A) = D(A)
\newcommand{\Seq}[1]{S_{=}(#1)}             % normalized states (trace = 1)
\newcommand{\Sleq}[1]{S_{\leq}(#1)}         % subnormalized states (0 <= trace <= 1)
\newcommand{\id}{\mathds{1}}                % identity operator
\newcommand{\idon}[1]{\mathds{1}_{#1}}      % identity on a named system, e.g. \idon{C}

% DEPRECATED: applies \mathcal to both arguments itself, so it takes bare
% letters, not system labels. Prefer \Linn; do not use in new contracts.
\providecommand{\Lop}[2]{\mathcal{L}(\mathcal{#1}, \mathcal{#2})}


% -----------------------------------------------------------------------------
% 2. Composite systems
% -----------------------------------------------------------------------------
\newcommand{\dsum}{\oplus}                  % direct sum of system labels
% Tensor product: use \otimes directly.
% n-fold direct sum of an operator space: \Lin{A}^{\oplus n}


% -----------------------------------------------------------------------------
% 3. Events, conditioning and classical registers (for STATES)
%    Convention:  \wedge  -> subnormalized "and-event" component, e.g. \rho_{A\wedge x}
%                 |        -> normalized conditional state,        e.g. \rho_{A|x}
% (Write these directly in subscripts; no extra macros needed.)
% -----------------------------------------------------------------------------


% -----------------------------------------------------------------------------
% 4. Photon-number subspaces / direct-sum blocks
%    Convention:  operators carry the restriction "on top", in brackets;
%                 states carry it as a \wedge subscript (see Sec. 3).
% -----------------------------------------------------------------------------
\newcommand{\Pnum}[1]{P_{#1}}               % photon-number projector (system from context)
\newcommand{\Pnumon}[2]{P_{#1}^{#2}}        % with explicit system superscript, e.g. P_n^C
\newcommand{\blk}[2]{#1^{(#2)}}             % operator block, e.g. \blk{M}{\leqN} = M^{(<=N)}
\newcommand{\leqN}{\leq N}                  % use as \blk{M}{\leqN}
\newcommand{\gtN}{> N}                      % use as \blk{M}{\gtN}
\newcommand{\eqn}{= n}                      % use as \blk{M}{\eqn}


% -----------------------------------------------------------------------------
% 5. System labels used in the protocol
% -----------------------------------------------------------------------------
\newcommand{\Chan}{C}                       % quantum-channel (Eve) input to Bob
\newcommand{\Prb}{P}                        % probe input
\newcommand{\Flag}{{\mathrm{Flg}}}          % flag system (EXISTS ONLY AFTER squashing)
\newcommand{\Loss}{{\mathrm{Lss}}}          % loss system
\newcommand{\Vc}{{\mathrm{Vcm}}}            % vacuum system
\newcommand{\Nout}{N_{\mathrm{out}}}
\newcommand{\Nin}{N_{\mathrm{in}}}


% -----------------------------------------------------------------------------
% 6. Alphabets
% -----------------------------------------------------------------------------
\newcommand{\Yset}{\mathcal{Y}}             % Bob's (coarse-grained) detection alphabet
\newcommand{\Yraw}{\mathcal{Y}_{\mathrm{raw}}}% raw click patterns {0,1}^4
\newcommand{\Yscset}{\mathcal{Y}_{\singlec}}% single-click outcomes
\newcommand{\Ymcset}{\mathcal{Y}_{\mc}}     % multi-click outcomes
\newcommand{\Tset}{\mathcal{T}_{\mathrm{set}}}% private setting set (the probe settings in the optical instance)
\newcommand{\Xset}{\mathcal{X}}             % Alice's alphabet
\newcommand{\Dset}{\{H,V,+,-\}}             % polarization / detector labels
\newcommand{\pass}{\top}                    % pass sifting


% -----------------------------------------------------------------------------
% 7. Mode / outcome abbreviations
% -----------------------------------------------------------------------------
\newcommand{\on}{{\mathrm{on}}}
\newcommand{\off}{{\mathrm{off}}}
\newcommand{\dc}{{\mathrm{dc}}}             % dark count
\newcommand{\sq}{\mathrm{sq}}               % squashed
\newcommand{\mc}{\mathrm{mc}}               % multi-click
\newcommand{\nc}{\mathrm{nc}}               % no-click
\newcommand{\singlec}{\mathrm{sc}}          % single-click
\newcommand{\vac}{\mathrm{vac}}             % vacuum
\newcommand{\kout}{{k_\mathrm{out}}}
\newcommand{\kin}{{k_\mathrm{in}}}


% -----------------------------------------------------------------------------
% 8. Maps, channels and miscellaneous operators
% -----------------------------------------------------------------------------
\newcommand{\Prob}{\mathrm{Prob}}           % probability (use \Prob[\cdot] or \Prob(\cdot))
\newcommand{\lmin}{\lambda_{\min}}          % minimum eigenvalue
\newcommand{\Wflag}{W_{\Flag}}              % flag-space weight
\newcommand{\transpose}{\mathsf{T}}
\newcommand{\innerp}[2]{\left\langle #1 \vert #2 \right\rangle}
\renewcommand{\vec}[1]{\boldsymbol{#1}}     % NB: overrides the LaTeX arrow accent


% -----------------------------------------------------------------------------
% 9. Key-rate numerics
% -----------------------------------------------------------------------------
\newcommand{\cG}{\mathcal{G}}               % protocol map
\newcommand{\cZ}{\mathcal{Z}}               % pinching map
\newcommand{\Amap}{\mathcal{A}_{\mathrm{ann}}} % announcement + measurement map
\newcommand{\Amapnb}{\mathcal{A}_{\mathrm{ann}|\mathrm{nb}}}
\newcommand{\simpl}{}                       % empty suffix for simplified block maps
\newcommand{\acSet}{\mathcal{S}_{\mathrm{acc}}} % acceptance set
\newcommand{\testE}{\mathcal{E}_{\mathrm{test}}} % test-statistics map
\newcommand{\ecLeak}{\mathrm{leak}^{\mathrm{EC}}_{\mathrm{obs}}}
\newcommand{\pgen}{p_{\mathrm{gen}}}
\newcommand{\askeyr}{R_{\infty}}             % asymptotic key rate
\newcommand{\ppass}{p_{\mathrm{pass}}}       % sifting pass probability


% -----------------------------------------------------------------------------
% 10. Finite-size analysis / Renyi entropies
% -----------------------------------------------------------------------------

% --- Round types: test vs. key generation (register T, alphabet \Fset) --------
\newcommand{\tst}{\mathrm{test}}            % test round
\newcommand{\gen}{\mathrm{gen}}             % key-generation round
\newcommand{\Fset}{\mathcal{T}}             % round-type alphabet {test, gen} (not \Tset)
\newcommand{\Chat}{\widehat{C}}             % acceptance-test register
\newcommand{\Chatset}{\widehat{\mathcal{C}}}% its alphabet

% --- Basis / key extraction functions ----------------------------------------
\newcommand{\bA}{b_{A}}                     % Alice's basis-extraction function
\newcommand{\bB}{b_{B}}                     % Bob's basis-extraction function
\newcommand{\kB}{k_{B}}                     % Bob's key-extraction function

% --- Maps --------------------------------------------------------------------
\newcommand{\Mprot}{\mathcal{M}}            % single-round protocol map
\newcommand{\Phimap}{\boldsymbol{\Phi}}     % test-statistics map
\newcommand{\Pdeph}{\mathcal{P}}            % dephasing (pinching) map on \Blreg

% --- Isometries --------------------------------------------------------------
\newcommand{\Vmsm}{V_{\mathrm{msm}}}        % measurement isometry
\newcommand{\Vann}{V_{\mathrm{ann}}}        % announcement isometry
\newcommand{\VI}{V_{I}}                     % Eve's copy of the announcements
\newcommand{\VChat}{V_{\Chat}}              % acceptance-data isometry
\newcommand{\VR}{V_{R}}                     % sifting + key map -> raw key R
\newcommand{\VZ}{V_{Z}}                     % raw key -> key register Z
\newcommand{\Wprot}{W}                      % full protocol isometry

% --- Renyi entropies ---------------------------------------------------------
% SECURITY-CRITICAL: see spec/notation.md §4. Never change an arrow to
% make an inequality work out; never convert silently between the sandwiched
% (tilde) and Petz (overline) families.
\newcommand{\Hsandup}{\widetilde{H}^{\uparrow}_{\alpha}}   % sandwiched, up
\newcommand{\Hsanddown}{\widetilde{H}^{\downarrow}_{\alpha}}% sandwiched, down
\newcommand{\Hpetzup}{\overline{H}^{\uparrow}_{\alpha}}    % Petz, up
\newcommand{\Hpetzdown}{\overline{H}^{\downarrow}_{\alpha}}% Petz, down
\newcommand{\hup}{h^{\uparrow}_{\alpha}}                   % single-round rate
\newcommand{\hqkd}{h^{\uparrow,\mathrm{QKD}}_{\alpha}}     % single-round QKD rate
\newcommand{\hlow}{\underline{h}_{\alpha}}                 % lower bound on \hqkd
% --- spec/notation.md 3.5 and 4.5 -------------------------------------------
% \hlb  is the first-order linearization lower bound on \hlow: the headline
%       key rate (the `first-order` tier of foundations/def:tiers). NOT \hlow
%       itself.
% \tnum is the instance's one declared numerical tolerance. Its
%       near-collision with the declared marginal \tau_A is disclosed in
%       spec/notation.md 3.5; bare \tau is reserved and denotes neither.
\newcommand{\hlb}{\underline{h}^{\mathrm{FO}}_{\alpha}}    % first-order lower bound on \hlow
\newcommand{\tnum}{\tau_{\mathrm{num}}}                    % declared numerical tolerance


% =============================================================================
%  PART B -- module-local macros, grouped by module. Deprecated aliases,
%  defined so that a text written with them compiles, used in no new contract: \Lop,
%  \testE, \Pdeph, \pass, \ecLeak (spec/notation.md section 3.9).
% =============================================================================

% --- meat-core -----------------------------------------------------------------
\newcommand{\SecSet}{\mathbb{S}}            % sector index set of Q
\newcommand{\Qsys}{Q}                       % declared sector-carrying measured system
\newcommand{\Vset}{\mathcal{V}}             % declared joint outcome alphabet of the Q-POVM
\newcommand{\CAset}{\mathcal{C}_{A}}        % Alice announcement alphabet
\newcommand{\CBset}{\mathcal{C}_{B}}        % Bob announcement alphabet
\newcommand{\Party}{\mathfrak{P}}           % declared party partition
\newcommand{\Oset}[1]{\mathcal{O}_{#1}}     % readout alphabet of party #1
\newcommand{\rd}[1]{\pi_{#1}}               % readout map of party #1
\newcommand{\Sadm}{\mathcal{S}_{A}}         % admissible Alice subspace
\newcommand{\nomarg}{\mathsf{none}}         % "no trusted marginal declared"
\newcommand{\Zset}{\mathcal{Z}_{\perp}}     % raw-key alphabet incl. \perp
\newcommand{\Zval}{\mathcal{Z}}             % key alphabet without \perp
\newcommand{\Etil}{\widetilde{E}}           % single-round purifying register
\newcommand{\Ehat}{\widehat{E}}             % Eve's n-round register
\newcommand{\OmAT}{\Omega_{\mathrm{AT}}}    % acceptance event
\newcommand{\epa}{\varepsilon_{\mathrm{PA}}}
\newcommand{\eev}{\varepsilon_{\mathrm{EV}}}
\newcommand{\lEC}{\lambda_{\mathrm{EC}}}
\newcommand{\Decl}{\mathfrak{D}}            % protocol declaration
\newcommand{\CList}{\mathfrak{L}}           % declared constraint list
\newcommand{\PinS}{\mathcal{P}_{\SecSet}}   % sector pinching on Q
\newcommand{\Feas}{\mathcal{W}}             % feasible set of block states
\newcommand{\Chkey}{\Chatset_{\mathrm{key}}}% key-live acceptance values
\newcommand{\Kann}{\mathcal{K}}             % announcement-class alphabet
\newcommand{\hold}{\mathsf{h}}              % declared key holder (A or B)
\newcommand{\sfD}{D}                        % classical KL divergence (base 2)
\newcommand{\Dtwo}{D}                       % quantum relative entropy (base 2)
\newcommand{\Dsand}{\widetilde{D}_{\alpha}} % sandwiched Renyi divergence
\newcommand{\idmap}{\mathrm{id}}            % identity channel
\newcommand{\Aev}{\mathcal{A}}              % Eve's per-round attack channel
\newcommand{\AttSet}{\mathfrak{A}}          % licensed attack-channel set
\newcommand{\Chacc}{\Chatset_{\mathrm{acc}}}% acceptance support of \acSet
\newcommand{\epsstar}{\varepsilon^{\star}}      % distance to an exact POVM
\newcommand{\dfH}{\Delta_{\mathrm{herm}}}       % Hermiticity defect
\newcommand{\dfP}{\Delta_{\mathrm{pos}}}        % negative part
\newcommand{\dfC}{\Delta_{\mathrm{cpl}}}        % completeness defect
\newcommand{\lmoq}{g}                       % acceptance-side gradient block
\newcommand{\lmoJ}{G}                       % feasible-side gradient block
\newcommand{\bacc}{\beta_{\mathrm{acc}}}    % certified acceptance-side minimum
\newcommand{\Itset}{\mathfrak{I}}           % finite index set of iterates (C22)

% --- reductions ----------------------------------------------------------------
\newcommand{\Bsys}{B}                        % physical (pre-squashing) system
\newcommand{\Tcl}[1]{\mathcal{T}(#1)}        % trace-class operators on H(#1)
\newcommand{\Ins}{\mathcal{I}}               % generalized-instrument branch map
\newcommand{\redNmsm}{\mathcal{N}}              % POVM measurement channel
\newcommand{\Eprot}{\mathcal{E}}             % n-round protocol channel
\newcommand{\Grnd}{\mathsf{G}}               % one round map of a protocol
\newcommand{\Kfin}{\mathsf{K}}               % end-of-protocol processing
\newcommand{\Crit}{\mathsf{sec}}             % value of a security criterion
\newcommand{\Pflag}{\Pi_{\Flag}}             % flag-sector projector on Q
\newcommand{\Pleq}{\Pi_{\leqN}}              % preserved-block projector
\newcommand{\Pgt}{\Pi_{\gtN}}                % above-cutoff projector
\newcommand{\Sph}{S^{\mathrm{ph}}}           % family of physical inputs
\newcommand{\Srch}{S^{\mathrm{rch}}}         % reachable squashed states
\newcommand{\Ssq}{S^{\sq}}                   % squashed admissible set
\newcommand{\Ffss}{F^{\mathrm{FSS}}}         % flag-state-squasher target
\newcommand{\Fwp}{F^{\mathrm{WP}}}           % multi-bound WPFSS target
\newcommand{\Mwp}{M^{\mathrm{WP}}}           % renormalized above-cutoff POVM
\newcommand{\Mdec}{{M'}^{\mathrm{ph}}}       % declared physical family on B
\newcommand{\Ppvm}{P^{\Flag}}                % flag PVM
\newcommand{\Lmp}{{\Lambda_{\mathrm{mp}}}}   % measure-and-prepare channel
\newcommand{\Lwp}{{\Lambda^{\mathrm{WP}}}}   % multi-bound WPFSS channel
\newcommand{\Lfss}{{\Lambda^{\mathrm{FSS}}}} % flag-state-squasher channel
\newcommand{\lmax}{\lambda_{\max}}           % maximum eigenvalue
\newcommand{\Lnoise}{\widetilde{\Lambda}}    % flag-noise channel
\newcommand{\cg}{\varpi}                     % declared coarse-graining map
\newcommand{\MCK}{\mathsf{MC}}               % multiclick raw-pattern set
\newcommand{\redChain}{\mathfrak{C}}            % reduction chain
\newcommand{\Lnk}{\ell}                      % a single typed link
\newcommand{\dsq}{\preceq}                      % the squashing relation
\newcommand{\opn}[1]{\left\lVert #1\right\rVert}           % operator (spectral) norm
\newcommand{\mxn}[1]{\left\lVert #1\right\rVert_{\max}}    % entrywise max modulus
\newcommand{\Fld}{\mathbb{F}}                   % exact ordered substrate field
\newcommand{\Fldi}{\mathbb{F}_{i}}              % its Gaussian extension F + iF
\newcommand{\lat}[1]{\big\lfloor #1\big\rfloor_{\kappa}} % round DOWN onto the lattice

% --- protocol-compiler ---------------------------------------------------------
\newcommand{\Spec}{\mathfrak{S}}             % protocol specification (compiler input)
\newcommand{\pcChain}{\mathfrak{R}}            % reduction chain (two typed segments)
\newcommand{\Chsrc}{\pcChain_{\mathrm{src}}}   % its Alice-side source-map segment
\newcommand{\Chrec}{\pcChain_{\mathrm{rec}}}   % its receiver-side squashing segment
\newcommand{\rvirt}{\rho^{\mathrm{v}}}       % virtual (PRE-source-map) signal state
\newcommand{\Fph}{F^{\mathrm{ph}}}           % physical (pre-reduction) POVM element
\newcommand{\Vph}{\mathcal{V}^{\mathrm{ph}}} % physical outcome alphabet
\newcommand{\Tpub}{\mathcal{T}_{\mathrm{pub}}}% public-choice alphabet
\newcommand{\Uset}{\mathcal{U}}              % residual outcome alphabet
\newcommand{\idxm}{\mathsf{idx}}             % label-to-index map
\newcommand{\Asig}{A'}                       % transmitted (signal) system -- the REAL one
\newcommand{\Avirt}{A''}                     % virtual transmitted system (source descriptor of Xi_1)
\newcommand{\Abar}{\bar{A}}                  % Alice's measured register
\newcommand{\Ahat}{\hat{A}}                  % Alice's shield register
\newcommand{\pcNmsm}{\mathcal{N}_{\mathsf A}}  % Alice's measurement channel (see the RULES block)
\newcommand{\Phip}{\Phi^{+}}                 % maximally entangled vector label
\newcommand{\Smap}{\mathfrak{s}}             % the emitted sector map (fraktur: a compile-time object)

% --- receiver ------------------------------------------------------------------
\newcommand{\Modes}{\mathbb{M}}              % declared ordered mode set
\newcommand{\Inp}{\mathbb{I}}                % declared input modes
\newcommand{\InpC}{\mathbb{I}_{\Chan}}       % channel input modes
\newcommand{\InpS}{\mathbb{I}_{\Prb}}        % probe (source) input modes
\newcommand{\Env}{\mathbb{E}}                % dilation (environment) modes
\newcommand{\Port}[1]{\mathbb{P}_{#1}}       % mode set of detector port #1
\newcommand{\Emd}[1]{E_{#1}}                 % coordinate projector on modes #1
\newcommand{\Gmat}{G}                        % instrument (transfer) matrix
\newcommand{\Fock}{\mathfrak{F}}             % Fock / second-quantization functor
\newcommand{\Symn}[1]{\mathrm{Sym}^{#1}}     % symmetric tensor power
\newcommand{\normord}[1]{{:}\mspace{2mu}#1\mspace{2mu}{:}}   % normal ordering
\newcommand{\rk}{\operatorname{rank}}        % rank (physics.sty owns \rank)
\newcommand{\SCK}{\mathsf{SC}}               % single-click pattern set
\newcommand{\Recv}{\mathfrak{R}}             % receiver instrument tuple
\newcommand{\Cat}{\mathfrak{C}}              % component catalogue
\newcommand{\Grph}{\mathfrak{E}}             % declared element graph
\newcommand{\defic}{\mathbb{L}}              % loss deficiency operator
\newcommand{\Pvac}[1]{\Pi_{\vac}(#1)}        % joint no-click element on ports #1
\newcommand{\pmc}{p_{\mc}}                   % uniform high-photon mc weight
\newcommand{\Bnd}{\beta}                     % the scalar bound function beta(n)
\newcommand{\Rot}{R}                         % planar rotation matrix
\newcommand{\Wplate}{W}                      % retarder (wave-plate) matrix
\newcommand{\Bsp}{B}                         % beamsplitter core matrix

% --- statistics ----------------------------------------------------------------
\newcommand{\honest}{p^{\mathrm{hon}}}          % honest per-round distribution
\newcommand{\ecomp}{\varepsilon_{\mathrm{comp}}}% completeness budget
\newcommand{\simplex}[1]{\Delta(#1)}            % probability simplex
\newcommand{\Chatcon}{\Chatset_{\mathrm{con}}}  % constrained outcomes
\newcommand{\Ireg}[3]{\mathcal{I}_{#1}(#2,#3)}  % regularized incomplete beta
\newcommand{\Ncount}[1]{\mathsf{N}_{#1}}        % per-outcome count
\newcommand{\qlo}{\mathbf{q}^{-}}               % box lower bounds
\newcommand{\qhi}{\mathbf{q}^{+}}               % box upper bounds

\makeatletter
\newcommand*{\ClaimHeadFont}{\bfseries}   % switched by prominence
\newcommand*{\ClaimNoteFont}{\mdseries}
\newcommand*{\ClaimShortId}{}
\newcommand*{\SetShortId}[1]{\renewcommand*{\ClaimShortId}{#1}}
\newcommand*{\RestatedName}{}
% Upright bodies throughout: italic is reserved for non-normative text.
% \@endparpenalty\@M keeps the meta line on the statement's page.
\newtheoremstyle{claim}{\topsep}{\topsep}{\normalfont\@endparpenalty\@M}{}%
  {\ClaimHeadFont}{.}{.5em}{\thmname{#1}\thmnumber{ #2}\thmnote{ {\ClaimNoteFont(#3)}}}
\newtheoremstyle{claimass}{\topsep}{\topsep}{\normalfont\@endparpenalty\@M}{}%
  {\ClaimHeadFont}{.}{.5em}{\thmname{#1}\thmnumber{ #2}%
  \ifx\ClaimShortId\@empty\else\ [\textsf{\ClaimShortId}]\fi%
  \thmnote{ {\ClaimNoteFont(#3)}}}
\newtheoremstyle{claimremark}{.5\topsep}{.5\topsep}{\small\@endparpenalty\@M}{}%
  {\itshape}{.}{.5em}{\thmname{#1}\thmnumber{ #2}\thmnote{ {\upshape(#3)}}}
\newtheoremstyle{claimreq}{\topsep}{\topsep}{\normalfont\@endparpenalty\@M}{}%
  {\color{reqAccent}\scshape}{.}{.5em}%
  {\thmname{#1}\thmnumber{ #2}\thmnote{ {\upshape(#3)}}}
\makeatother
\theoremstyle{claim}
\newtheorem{theorem}{Theorem}[section]
\newtheorem{lemma}[theorem]{Lemma}
\newtheorem{proposition}[theorem]{Proposition}
\newtheorem{corollary}[theorem]{Corollary}
\newtheorem{fact}[theorem]{Fact}
\newtheorem{definition}[theorem]{Definition}
\newtheorem*{restatedclaim}{\RestatedName}
\theoremstyle{claimass}
\newtheorem{assumption}[theorem]{Assumption}
\newtheorem*{restatedassumption}{\RestatedName}
\theoremstyle{claimremark}
\newtheorem{remark}[theorem]{Remark}
\newtheorem{example}[theorem]{Example}
\theoremstyle{claimreq}
\newtheorem{requirement}[theorem]{Requirement}
\newtheorem*{restatedrequirement}{\RestatedName}
%% --- 8. Prominence: main results a margin rule and a bold heading ----------
\newtcolorbox{mainclaim}{enhanced jigsaw, breakable, blanker,
  borderline west={1.5pt}{-8pt}{mainRule}, left=0pt, right=0pt, top=0pt, bottom=0pt,
  boxsep=0pt, before skip=\topsep, after skip=\topsep, parbox=false,
  before upper={\renewcommand*{\ClaimNoteFont}{\bfseries}}}
\newenvironment{technicalclaim}{\renewcommand*{\ClaimHeadFont}{\mdseries\color{techGrey}}%
  \renewcommand*{\ClaimNoteFont}{}}{}
%% --- 9. Non-normative text, the meta line, navigation ------------------------
\makeatletter
% Tie what follows to this paragraph: no page break, no list penalty before
% the next theorem head (the \@nobreak route a section heading uses).
\newcommand{\keepwithnext}{\par\nobreak\vspace{3pt}\@nobreaktrue
  \everypar{\if@nobreak\@nobreakfalse\clubpenalty\@M\else\clubpenalty\@clubpenalty\everypar{}\fi}}
\newcommand{\claimsummary}[2]{% #1 claim id (named destination), #2 summary
  \par\addvspace{\bigskipamount}\Needspace{6\baselineskip}%
  {\small\itshape\color{nnGrey}\noindent\hypertarget{claim-#1}{}#2\par}\keepwithnext}
\newcommand{\claimanchor}[1]{% a claim without a summary: the destination alone
  \par\addvspace{\bigskipamount}\Needspace{6\baselineskip}%
  \phantomsection\hypertarget{claim-#1}{}\keepwithnext}
\newcommand{\claimleadin}[1]{{\color{nnGrey}\noindent#1\par}\keepwithnext}
\newcommand{\claimmeta}[1]{\par\nobreak\vspace{2pt}%
  {\footnotesize\color{metaGrey}\raggedright\interlinepenalty\@M\noindent#1\par}%
  \addvspace{\medskipamount}}
% After an assumption: the sentence generated from its verification field; after
% any claim with one: its plain-text reading (user_statement). Both
% non-normative, both grey.
\newcommand{\claimverification}[1]{\par\nobreak\vspace{2pt}%
  {\footnotesize\color{metaGrey}\noindent#1\par}\nobreak}
\newcommand{\claimreading}[1]{\par\nobreak\vspace{2pt}%
  {\footnotesize\color{nnGrey}\noindent\emph{Reading:}\enspace#1\par}\nobreak}
\makeatother
\newcommand{\claimbookmark}[4]{% #1 id, #2 kind word, #3 label, #4 plain title
  \bookmark[dest=claim-#1,level=2]{#2 \getrefnumber{#3} #4}}
\newcommand{\metasep}{\unskip\enspace\textperiodcentered\enspace\ignorespaces}
\newcommand{\metalabel}[1]{\emph{#1}\enspace}
\newcommand{\sectionintro}[1]{{\color{nnGrey}\noindent#1\par}\smallskip}
\newcommand{\sectionimports}[1]{{\footnotesize\color{metaGrey}\raggedright\noindent\metalabel{Imports}#1\par}\smallskip}
\newenvironment{overview}{\par\color{nnGrey}\itshape\noindent\ignorespaces}{\par}
% References: the whole "Def. 2.1" is the link; ids are the fallback text of
% every link into another PDF (a relative URI to a named destination).
\newcommand{\mref}[2]{\hyperref[#2]{#1~\ref*{#2}}}
\newcommand{\aref}[2]{\hyperref[#2]{\textsf{#1}}}
\DeclareUrlCommand{\claimid}{\urlstyle{sf}\def\UrlBreaks{\do\/\do\:\do\-\do\.\do\_}%
  \def\UrlBigBreaks{}\def\UrlNoBreaks{}}
\newcommand{\extref}[2]{\hyperref{#1}{}{claim-#2}{\claimid{#2}}}
\newcommand{\extaref}[3]{\hyperref{#1}{}{claim-#2}{\textsf{#3}}}
\newcommand{\coderef}[2]{\href{#1}{\claimid{#2}}}
% Claim-body commands (spec/claim-format.md section 4): \xref{<full id>} prints
% the kind, title and id of a claim of another contract, linked to its stable
% destination in that contract's PDF (one \DeclareXref per cited claim below);
% \term[gloss]{phrase} marks a defined phrase for the index (the gloss feeds
% the Terms table); \tref[phrase]{text} is a use of a defined term. The
% assembler rewrites both in the body into \termdef (the definition's named
% destination term-<slug>) and \termref / \exttermref (a link to it, in this
% or another contract's PDF); the fallbacks below serve text it copies
% elsewhere, such as the digest. \iface is an interface key.
\makeatletter
\newcommand{\DeclareXref}[4]{% id, kind word, title, contract PDF
  \@namedef{xref@#1}{\hyperref{#4}{}{claim-#1}{#2 ``#3'' ({\small\claimid{#1}})}}}
\newcommand{\xref}[1]{\@ifundefined{xref@#1}%
  {\textbf{[unresolved \nolinkurl{#1}]}\@latex@warning{Unresolved xref #1}}%
  {\@nameuse{xref@#1}}}
\makeatother
\DeclareUrlCommand{\iface}{\urlstyle{tt}}
\newcommand{\term}[2][]{\emph{#2}\index{#2}}
\newcommand{\tref}[2][]{#2}
\newcommand{\termdef}[2]{\hypertarget{#1}{}\emph{#2}\index{#2}}
\newcommand{\termref}[2]{\hyperlink{#1}{#2}}
\newcommand{\exttermref}[3]{\hyperref{#1}{}{#2}{#3}}
\hypersetup{pdftitle={Foundations: quantum-information preliminaries, number model, exact admissibility and emission -- statements digest 0.12.0}}

% The digest: the contract's labels under the prefix C-, its numbers checked.
\externaldocument[C-]{contract}
\theoremstyle{claim}\newtheorem*{digestclaim}{\RestatedName}
\theoremstyle{claimass}\newtheorem*{digestassumption}{\RestatedName}
\theoremstyle{claimreq}\newtheorem*{digestrequirement}{\RestatedName}
\makeatletter
% #1 label, #2 the number the assembler computed: a local label carrying it, so
% that a reference between two digest items stays in the digest; the build
% stops when the contract's own number differs.
\newcommand{\digestlabel}[2]{\phantomsection
  \edef\digest@theirs{\getrefnumber{C-#1}}\def\digest@ours{#2}%
  \ifx\digest@theirs\digest@ours\else
    \PackageError{digest}{#1 is #2 here but \digest@theirs\space in the contract}{}\fi
  \def\@currentlabel{#2}\label{#1}}
\makeatother
% A display's local label (amsmath records the \tag, i.e. the contract's number).
\newcommand{\digesteqlabel}[1]{\label{#1}}

% Claims of other contracts cited here (\xref, [[...]]).
\DeclareXref{meat-core/ass:fixed}{Assumption}{Protocol parameters fixed before the run}{../meat-core/contract.pdf}
\DeclareXref{meat-core/frank-wolfe/lem:certified-bracket}{Lemma}{The certified bracket from enclosures of the gradient}{../meat-core/contract.pdf}
\DeclareXref{meat-core/frank-wolfe/req:bracket-on-reference}{Requirement}{The first-order bound evaluated on the reference companion}{../meat-core/contract.pdf}
\DeclareXref{meat-core/frank-wolfe/req:certified-evaluation}{Requirement}{The certified evaluation of the first-order bound}{../meat-core/contract.pdf}
\DeclareXref{meat-core/frank-wolfe/thm:firstorder}{Theorem}{The estimate and the first-order lower bound}{../meat-core/contract.pdf}
\DeclareXref{meat-core/lem:dk}{Lemma}{Divided-difference derivative of a matrix function}{../meat-core/contract.pdf}
\DeclareXref{meat-core/lem:estimate}{Lemma}{Upper estimate of the value at any feasible point}{../meat-core/contract.pdf}
\DeclareXref{meat-core/lem:memory-lift}{Lemma}{Approximating an attack by all-at-once attacks with finite memory}{../meat-core/contract.pdf}
\DeclareXref{meat-core/prop:optol-nobound}{Proposition}{Instability of the rate under measurement perturbations}{../meat-core/contract.pdf}
\DeclareXref{meat-core/thm:gradient}{Theorem}{Gradient of the pinched trace power on a face}{../meat-core/contract.pdf}
\DeclareXref{statistics/binomial-box/lem:widening}{Lemma}{Widening the box keeps its guarantees}{../statistics/contract.pdf}

\begin{document}
\begin{center}
{\LARGE\bfseries Foundations: quantum-information preliminaries, number model, exact admissibility and emission}\\[4pt]
{\large Statements digest, version 0.12.0}
\end{center}
\noindent{\small Every definition, fact, assumption, result and requirement of the \href{contract.pdf}{contract}, with the contract's numbers; remarks, examples, technical steps and all proofs are left out. Each heading opens its statement in the contract.}

\part*{Mathematics}\pdfbookmark[0]{Mathematics}{digest-part-1}

\section*{1\quad Quantum-information preliminaries}\pdfbookmark[1]{1 Quantum-information preliminaries}{digest-group-foundations/def:pinching}

% claim foundations/def:pinching
\digestlabel{def:pinching}{1.1}\pdfbookmark[2]{Definition 1.1 Operator conventions and pinching}{digest-foundations/def:pinching}
\renewcommand*{\RestatedName}{\hyperref{contract.pdf}{}{claim-foundations/def:pinching}{Definition~1.1}}%
\begin{digestclaim}[Operator conventions and pinching]
All Hilbert spaces are finite-dimensional. For a system label $A$,
$\Lin{A}$ is the space of linear operators on $\Hs{A}$, $\Pos{A}$ its
positive semidefinite operators, and $\Seq{A}$ and $\Sleq{A}$ the
normalized and the subnormalized states, of trace equal to one and at most
one. A channel is a CPTP map; an instrument branch is a CP
trace-non-increasing (CPTNI) map. For $X\in\Pos{A}$ and $p>0$, the power
$X^{p}$ is defined by functional calculus on the support of $X$ and by $0^p=0$
on its kernel; in particular $X^{p}$ is block-diagonal for any orthogonal
decomposition for which $X$ is.

For a family of orthogonal projectors $\{\Pi_j\}_{j}$ on $\Hs{Q}$ with
$\sum_j\Pi_j=\idon{Q}$, the associated \emph{pinching} is the channel
$\tau\mapsto\sum_j\Pi_j\tau\Pi_j$. A pinching is CPTP, self-adjoint as a
superoperator ($\tr[Y\,\mathcal{E}(W)]=\tr[\mathcal{E}(Y)\,W]$), idempotent,
and unital. An operator $X$ is a fixed point iff it is block-diagonal,
$X=\sum_j\Pi_jX\Pi_j$; functional calculus of a block-diagonal PSD operator
acts blockwise, so a pinching commutes with taking powers of its fixed
points. Each of these properties is a one-line computation from the form
$\sum_j\Pi_j(\cdot)\Pi_j$ and orthogonality of the $\Pi_j$;
self-adjointness is
$\tr[Y\sum_j\Pi_jW\Pi_j]=\tr[\sum_j\Pi_jY\Pi_j\,W]$ by cyclicity.

A pinching is a uniform mixture of unitary conjugations. Enumerate the
projectors as $\Pi_1,\dots,\Pi_m$, put $\zeta\coloneqq e^{2\pi i/m}$ and
$U_k\coloneqq\sum_{j=1}^{m}\zeta^{jk}\Pi_j$ for $k\in\{0,\dots,m-1\}$; each
$U_k$ is unitary, $U_kU_k^{\dagger}=\sum_j|\zeta^{jk}|^2\Pi_j=\idon{Q}$ by
orthogonality, and
\begin{equation*}\tag{\getrefnumber{C-eq:pinch-unitary-mixture}}\digesteqlabel{eq:pinch-unitary-mixture}
  \sum_{j}\Pi_j X\Pi_j=\frac1m\sum_{k=0}^{m-1}U_kXU_k^{\dagger}
  \qquad\text{for every }X\in\Lin{Q},
\end{equation*}
since the right-hand side is
$\sum_{j,j'}\big(\frac1m\sum_k\zeta^{(j-j')k}\big)\Pi_jX\Pi_{j'}$ and the
inner average is $1$ for $j=j'$ and $0$ otherwise. In particular a pinching
does not increase the operator norm, a mixture of unitary conjugates of $X$
having norm at most $\|X\|$.
\end{digestclaim}

% claim foundations/def:renyi
\digestlabel{def:renyi}{1.2}\pdfbookmark[2]{Definition 1.2 Sandwiched Rényi divergence and conditional entropies}{digest-foundations/def:renyi}
\renewcommand*{\RestatedName}{\hyperref{contract.pdf}{}{claim-foundations/def:renyi}{Definition~1.2}}%
\begin{digestclaim}[Sandwiched Rényi divergence and conditional entropies]
For $\alpha\in(1,\infty)$, $\rho\in\Sleq{AB}$ and $\sigma\in\Pos{AB}$, on the
same system,
\begin{equation*}\tag{\getrefnumber{C-eq:sandwiched}}\digesteqlabel{eq:sandwiched}
  \Dsand(\rho\,\|\,\sigma)\coloneqq\frac{1}{\alpha-1}\log_2
  \tr\!\left[\left(\sigma^{\frac{1-\alpha}{2\alpha}}\,\rho\,
  \sigma^{\frac{1-\alpha}{2\alpha}}\right)^{\alpha}\right]
\end{equation*}
when $\operatorname{supp}\rho\subseteq\operatorname{supp}\sigma$, and
$\Dsand(\rho\|\sigma)\coloneqq+\infty$ otherwise. The exponent
$\frac{1-\alpha}{2\alpha}$ is negative, and the power of $\sigma$ is taken on
its support and is zero on its kernel, as for the generalised inverse; where
$\operatorname{supp}\rho\subseteq\operatorname{supp}\sigma$, $\rho=P\rho P$
for the projector $P$ onto that support, so the value reads the power on the
support alone. For
$\rho_{AB}\in\Seq{AB}$,
\begin{equation*}\tag{\getrefnumber{C-eq:Hdefs}}\digesteqlabel{eq:Hdefs}
  \Hsanddown(A|B)_{\rho}\coloneqq-\Dsand(\rho_{AB}\,\|\,\idon{A}\otimes\rho_B),
  \qquad
  \Hsandup(A|B)_{\rho}\coloneqq-\inf_{\sigma_B\in\Seq{B}}
  \Dsand(\rho_{AB}\,\|\,\idon{A}\otimes\sigma_B).
\end{equation*}
Both are in bits. These are the definitions of
\cite[Def.~5.2, Eqs.~(5.18)--(5.19)]{tomamichel2016finite}. Since the
infimum in \eqref{eq:Hdefs} includes $\sigma_B=\rho_B$,
\begin{equation*}\tag{\getrefnumber{C-eq:arrow-order}}\digesteqlabel{eq:arrow-order}
  \Hsanddown(A|B)_{\rho}\;\le\;\Hsandup(A|B)_{\rho}.
\end{equation*}
The substitution $\Hsandup\to\Hsanddown$ inside a lower bound is therefore
conservative, and the reverse is not.
\end{digestclaim}

% claim foundations/def:kl
\digestlabel{def:kl}{1.3}\pdfbookmark[2]{Definition 1.3 Classical divergences and the quantum relative entropy}{digest-foundations/def:kl}
\renewcommand*{\RestatedName}{\hyperref{contract.pdf}{}{claim-foundations/def:kl}{Definition~1.3}}%
\begin{digestclaim}[Classical divergences and the quantum relative entropy]
For probability vectors $\mathbf q,\mathbf p$ on a finite alphabet,
$\sfD(\mathbf q\|\mathbf p)\coloneqq\sum_c q(c)\log_2\frac{q(c)}{p(c)}$ with
the conventions $0\log_2\frac{0}{p}=0$ and $q\log_2\frac{q}{0}=+\infty$ for
$q>0$. For $\alpha\in(1,\infty)$ the classical R\'enyi divergence is
$D_{\alpha}(\mathbf q\|\mathbf p)\coloneqq
\frac{1}{\alpha-1}\log_2\sum_{c}q(c)^{\alpha}p(c)^{1-\alpha}$ (with
$+\infty$ if $\mathbf q\not\ll\mathbf p$), and
$D_{\infty}(\mathbf q\|\mathbf p)\coloneqq\log_2\max_c\frac{q(c)}{p(c)}$.
For $X,Y\in\Pos{}$ the quantum relative entropy is
$\Dtwo(X\|Y)=\tr[X(\log_2X-\log_2Y)]$ when
$\operatorname{supp}X\subseteq\operatorname{supp}Y$, else $+\infty$.

The Kullback--Leibler divergence of probability vectors is non-negative,
$\sfD(\mathbf q\|\mathbf p)\ge0$, with equality if and only if
$\mathbf q=\mathbf p$ (Gibbs' inequality). If $q(c)>0=p(c)$ for some $c$,
then $\sfD(\mathbf q\|\mathbf p)=+\infty$. Otherwise, with
$S\coloneqq\{c:q(c)>0\}$, the terms outside $S$ vanish by the first
convention and $\log_2$ is strictly concave, so Jensen's inequality with the
weights $q(c)$, $c\in S$, which sum to one, gives
$-\sfD(\mathbf q\|\mathbf p)=\sum_{c\in S}q(c)\log_2\frac{p(c)}{q(c)}
\le\log_2\sum_{c\in S}p(c)\le\log_21=0$. Equality in the first step
forces $p(c)/q(c)$ to be one constant $\gamma$ on $S$, and in the second
$\sum_{c\in S}p(c)=1$; then $\gamma=\gamma\sum_{c\in S}q(c)
=\sum_{c\in S}p(c)=1$, so $\mathbf p=\mathbf q$ on $S$, and $\mathbf p$
vanishes off $S$, its mass on $S$ being one. Conversely
$\sfD(\mathbf q\|\mathbf q)=0$ term by term.
\end{digestclaim}

% claim foundations/lem:kernel
\digestlabel{lem:kernel}{1.4}\pdfbookmark[2]{Lemma 1.4 Extended divergence kernel}{digest-foundations/lem:kernel}
\begin{mainclaim}
\renewcommand*{\RestatedName}{\hyperref{contract.pdf}{}{claim-foundations/lem:kernel}{Lemma~1.4}}%
\begin{digestclaim}[Extended divergence kernel]
Define $d:[0,\infty)\times[0,\infty)\to(-\infty,+\infty]$ by
$d(u,v)\coloneqq u\log_2\frac uv$ for $u,v>0$, $d(0,v)\coloneqq0$ for
$v\ge0$, and $d(u,0)\coloneqq+\infty$ for $u>0$. Then $d$ is jointly
convex and lower semicontinuous, and non-increasing in $v$.
\end{digestclaim}
\end{mainclaim}

% claim foundations/lem:pinchmono
\digestlabel{lem:pinchmono}{1.5}\pdfbookmark[2]{Lemma 1.5 Pinching lowers the trace power}{digest-foundations/lem:pinchmono}
\begin{mainclaim}
\renewcommand*{\RestatedName}{\hyperref{contract.pdf}{}{claim-foundations/lem:pinchmono}{Lemma~1.5}}%
\begin{digestclaim}[Pinching lowers the trace power]
Let $\tau\in\Pos{Q}$, let $\cZ$ be a pinching with projectors
$\{Z_j\}_{j=1}^{m}$, and let $\alpha\in(1,\infty)$. Then
\begin{equation*}\tag{\getrefnumber{C-eq:pinchmono}}\digesteqlabel{eq:pinchmono}
  \tr\!\left[\left(\cZ\big(\tau^{1/\alpha}\big)\right)^{\alpha}\right]
  \;\le\;\tr[\tau],
\end{equation*}
with equality if and only if $\cZ(\tau^{1/\alpha})=\tau^{1/\alpha}$,
equivalently iff $\tau$ is block-diagonal for $\{Z_j\}$, equivalently iff
$\cZ(\tau)=\tau$.
\end{digestclaim}
\end{mainclaim}

% claim foundations/fact:pinch-id
\digestlabel{fact:pinch-id}{1.6}\pdfbookmark[2]{Fact 1.6 Pinched trace-power form of the conditional Rényi entropy}{digest-foundations/fact:pinch-id}
\renewcommand*{\RestatedName}{\hyperref{contract.pdf}{}{claim-foundations/fact:pinch-id}{Fact~1.6}}%
\begin{digestclaim}[Pinched trace-power form of the conditional Rényi entropy]
Let $\rho_{QE}\in\Seq{QE}$ be pure, let $\{Z_j\}_{j=1}^{d}$ be orthogonal
projectors on $Q$ with $\sum_jZ_j=\idon{Q}$, let
$\VZ=\sum_j\ket{j}_{S}\otimes Z_j$ and
$\sigma_{SQE}\coloneqq\VZ\,\rho_{QE}\,\VZ^{\dagger}$, and let
$\cZ(\tau)\coloneqq\sum_jZ_j\tau Z_j$. Then, for $\alpha\in(1,\infty)$,
\begin{equation*}\tag{\getrefnumber{C-eq:pinch-id}}\digesteqlabel{eq:pinch-id}
  \Hsanddown(S|E)_{\sigma_{SE}}
  =\frac{1}{1-\alpha}\log_2
  \tr\!\left[\left(\cZ\big(\rho_Q^{1/\alpha}\big)\right)^{\alpha}\right].
\end{equation*}
\emph{Locator note.} The source states no $\alpha$ restriction in the
theorem; the domain $(1,\infty)$ used here lies inside the domain
$(0,1)\cup(1,\infty)$ of its entropy definitions and is the only range any
claim that uses this fact invokes.
\end{digestclaim}

% claim foundations/fact:concavity
\digestlabel{fact:concavity}{1.7}\pdfbookmark[2]{Fact 1.7 Concavity of the pinched trace-power functional}{digest-foundations/fact:concavity}
\renewcommand*{\RestatedName}{\hyperref{contract.pdf}{}{claim-foundations/fact:concavity}{Fact~1.7}}%
\begin{digestclaim}[Concavity of the pinched trace-power functional]
Let $\cZ$ be a pinching built from a complete family of mutually orthogonal
projectors (Def.~22 of the source) and for $\alpha'\in[0,2]$ let
$Q_{\alpha'}(\rho)\coloneqq
\tr[(\cZ(\rho^{\alpha'}))^{1/\alpha'}]$ (Eq.~(F6) of the source). Then
$\rho\mapsto Q_{\alpha'}(\rho)$ is concave on $\Sleq{Q}$ for
$\alpha'\in[0,1]$.
\end{digestclaim}

\section*{2\quad The round model}\pdfbookmark[1]{2 The round model}{digest-group-foundations/def:round-model}

% claim foundations/def:round-model
\digestlabel{def:round-model}{2.1}\pdfbookmark[2]{Definition 2.1 Round model, attack sequences and attack sets}{digest-foundations/def:round-model}
\renewcommand*{\RestatedName}{\hyperref{contract.pdf}{}{claim-foundations/def:round-model}{Definition~2.1}}%
\begin{digestclaim}[Round model, attack sequences and attack sets]
Fix a number $n\geq1$ of rounds, a per-round system $S$, Alice's system $A$,
and the instantiating model's source mode: in trusted-source mode Alice
supplies $A$ in every round, and in untrusted-source mode the adversary
does. Systems, operators and channels follow the conventions of
Definition~\ref{def:pinching}, except that $S$ may be separable of infinite
dimension where the instantiating model declares it so; a channel whose input
or output includes such a system is a normal CPTP map on trace-class
operators. The dimensions are the instantiating model's to declare. We write
$S_j$ and $A_j$ for the copies of $S$ and $A$ in round $j$, and drop the
round index where one round is in view. This is the \emph{round model} with
per-round system $S$.
\begin{enumerate}[label=(\roman*),leftmargin=2.4em]
\item An \emph{attack
  sequence} is a family $(\Aev_j)_{j=1}^{n}$ of channels whose output carries
  the systems the adversary supplies in the source mode:
  \begin{equation*}\tag{\getrefnumber{C-eq:round-attack}}\digesteqlabel{eq:round-attack}
    \Aev_j\in
    \begin{cases}
      \mathrm{CPTP}\big(E_{j-1},\ S_j\,E'_j\big),
        & \text{trusted-source mode},\\[0.3ex]
      \mathrm{CPTP}\big(E_{j-1},\ A_j\,S_j\,E'_j\big),
        & \text{untrusted-source mode},
    \end{cases}
  \end{equation*}
  where $E_0$ is Eve's initial register, $E'_j$ her outgoing memory, and
  $E_j$ her register after round $j$, which holds $E'_j$ together with what
  the instantiating model hands her in round $j$. In trusted-source mode $A_j$
  is \emph{not} an output of $\Aev_j$: it is supplied by Alice and carried
  into the round as a separate input.
\item Let $S$ be finite-dimensional, and let
  $\CList=\{(\hat O_k,\bowtie_k,c_k)\}_{k=1}^{m}$, $m\geq0$, be a
  finite list of Hermitian operators $\hat O_k$ on $S$, relations
  $\bowtie_k\in\{\preceq,=,\succeq\}$ between operators, and real numbers
  $c_k$. For registers $E',E''$, the \emph{attack
  set} of $\CList$ is
  \begin{equation*}\tag{\getrefnumber{C-eq:round-attset}}\digesteqlabel{eq:round-attset}
    \AttSet_{\CList}(E',\,S E'')\coloneqq
    \Big\{\Aev\in\mathrm{CPTP}(E',S E''):\
    \Aev^{\dagger}\big(\hat O_k\otimes\idon{E''}\big)
    \bowtie_k c_k\,\idon{E'}\ \ \forall k\Big\}
    \qquad\text{(trusted-source mode)},
  \end{equation*}
  and, with the constraint operator lifted to the larger supplied system,
  \begin{equation*}\tag{\getrefnumber{C-eq:round-attset-un}}\digesteqlabel{eq:round-attset-un}
    \AttSet_{\CList}(E',\,A S E'')\coloneqq
    \Big\{\Aev\in\mathrm{CPTP}(E',A S E''):\
    \Aev^{\dagger}\big(\idon{A}\otimes\hat O_k\otimes\idon{E''}\big)
    \bowtie_k c_k\,\idon{E'}\ \ \forall k\Big\}
  \end{equation*}
  in untrusted-source mode; the two are distinguished by the output system
  named in the argument. For $\CList=\varnothing$ the condition is empty and
  $\AttSet_{\varnothing}=\mathrm{CPTP}$.
\item Where $S$ is finite-dimensional and $S=\bigotimes_{f\in\mathbb{F}}S_f$
  is a product of factors indexed by a finite ordered set $\mathbb{F}$ (with a
  round index as well, a factor is written $S_{f,j}$), an entry of
  $\CList$ may be \emph{tagged} with a factor $f$; its operator $\hat O_k$ is
  then Hermitian on $S_f$, and the entry $(\hat O_k,\bowtie_k,c_k;f)$
  contributes to the attack set, in place of the condition of
  \eqref{eq:round-attset} or \eqref{eq:round-attset-un}, the condition
  \begin{equation*}\tag{\getrefnumber{C-eq:round-attset-pl}}\digesteqlabel{eq:round-attset-pl}
    \big(\idmap_{R'}\otimes\Aev\big)^{\dagger}
    \Big(X\otimes\big(\hat O_k-c_k\,\idon{S_f}\big)\otimes\idon{E''}\Big)
    \;\bowtie_k\;0
    \qquad\text{for every finite-dimensional }R'\text{ and every }
    X\in\Pos{R'R\,S_{f^c}} ,
  \end{equation*}
  as an operator on $R'E'$, where $S_{f^c}$ is the product of the factors
  other than $f$, the tensor factors of the argument are written in the order
  $R'R\,S_{f^c}$, $S_f$, $E''$, and $R$ is $A$ in untrusted-source mode and
  trivial in trusted-source mode. Untagged entries keep their conditions. The
  reference system $R'$ is what carries Alice's retained system in
  trusted-source mode, where $A$ is no output of $\Aev$. With $R'$ trivial and
  $X=\idon{}$, \eqref{eq:round-attset-pl} is the condition of
  \eqref{eq:round-attset} or \eqref{eq:round-attset-un} for the lifted operator
  $\idon{S_{f^c}}\otimes\hat O_k$, because $\Aev^{\dagger}$ is unital; so the
  attack set with tagged entries is contained in the attack set with the same
  entries untagged and their operators lifted.
\item An attack sequence is an \emph{admissible attack sequence} for
  $\CList$, or \emph{$\CList$-admissible}, if $\Aev_j$ lies in the attack set
  of $\CList$ of the source mode, \eqref{eq:round-attset} or
  \eqref{eq:round-attset-un} with the conditions of
  \eqref{eq:round-attset-pl} for tagged entries, with $E'=E_{j-1}$ and
  $E''=E'_j$, for every $j$.
\end{enumerate}
\end{digestclaim}

% claim foundations/def:round-homogeneity
\digestlabel{def:round-homogeneity}{2.2}\pdfbookmark[2]{Definition 2.2 Round homogeneity of a per-round family}{digest-foundations/def:round-homogeneity}
\renewcommand*{\RestatedName}{\hyperref{contract.pdf}{}{claim-foundations/def:round-homogeneity}{Definition~2.2}}%
\begin{digestclaim}[Round homogeneity of a per-round family]
Consider the $n$ rounds of a round model (Definition~\ref{def:round-model}),
and a \emph{per-round family}: a finite list of objects that a model of the
run declares for a round, such as operators, channels or a declaration. The
family is \emph{round-homogeneous}
if both of the following hold.
\begin{enumerate}[label=(\roman*),leftmargin=2.4em]
\item \emph{Sameness.} One list of objects, fixed before the run, describes
  every round: for every $j\in\{1,\dots,n\}$ round $j$ is described by the
  same objects. No object of the family depends on the round index or on the
  run.
\item \emph{No memory.} Every round behaves physically as the objects
  describe it, whatever the inputs and the outputs of the other rounds are.
  This excludes memory across rounds, such as detector memory, afterpulsing
  or correlations between a source's emissions in different rounds. It is a
  condition on the rounds, not on the objects, and does not follow from (i).
\end{enumerate}
An assumption may assert (i) alone of its family, and must then say so; an
assumption that calls its family round-homogeneous without saying so asserts
(i) and (ii). Rounds are not averaged: a family that describes the rounds
only on average over the rounds, or up to a per-round variation, does not
satisfy (i), and
per-round variation is not modelled. Round homogeneity constrains the family
and the rounds it describes, not the adversary: an attack sequence of the
round model remains arbitrary and may correlate the rounds.
\end{digestclaim}

% claim foundations/lem:register-compression
\digestlabel{lem:register-compression}{2.3}\pdfbookmark[2]{Lemma 2.3 Compressing a register onto a finite-dimensional subspace}{digest-foundations/lem:register-compression}
\renewcommand*{\RestatedName}{\hyperref{contract.pdf}{}{claim-foundations/lem:register-compression}{Lemma~2.3}}%
\begin{digestclaim}[Compressing a register onto a finite-dimensional subspace]
\emph{Setting.} In this lemma a Hilbert space may have any dimension, finite
or infinite, separable or not, in departure from the convention of
Definition~\ref{def:pinching}, and every system named is of this kind.
$\Tcl{Z}$ is the space of trace-class operators on $\Hs{Z}$ with the trace
norm $\|\cdot\|_1$, and $RM$ is the Hilbert-space tensor product of $R$ and
$M$. A \emph{channel} $\Tcl{Z}\to\Tcl{Z'}$ is a trace-preserving linear map
$\Phi$ that is completely positive: for every $k\ge1$ the map that applies
$\Phi$ to every entry of a $k\times k$ matrix of trace-class operators, which
is $\Tcl{\mathbb{C}^kZ}$, is positive. $\tr_R\colon\Tcl{RM}\to\Tcl{M}$ is the
partial trace, fixed by $\tr[X\tr_R(\sigma)]=\tr[(\idon{R}\otimes X)\sigma]$
for every bounded operator $X$ on $\Hs{M}$, and $\tr_M$ likewise.
\begin{enumerate}[label=(\roman*),leftmargin=2.6em]
\item \emph{(compression and inclusion)} Let $M^{\circ}$ be a non-zero
  finite-dimensional subspace of $\Hs{M}$, $V$ its inclusion isometry into
  $\Hs{M}$, $\Pi\coloneqq VV^{\dagger}$ the orthogonal projection onto it,
  $\Pi^{\perp}\coloneqq\idon{M}-\Pi$, and $\psi$ a state on $M^{\circ}$. For
  every system $R$ put, for $\sigma\in\Tcl{RM}$ and
  $\tau\in\Tcl{RM^{\circ}}$,
  \begin{equation*}\tag{\getrefnumber{C-eq:compression}}\digesteqlabel{eq:compression}
  \begin{gathered}
    \mathcal{P}_R(\sigma)\coloneqq
      (\idon{R}\otimes V)^{\dagger}\sigma(\idon{R}\otimes V)
      +\tr_M\big[(\idon{R}\otimes\Pi^{\perp})\,\sigma\,
        (\idon{R}\otimes\Pi^{\perp})\big]\otimes\psi ,\\
    \iota_R(\tau)\coloneqq(\idon{R}\otimes V)\,\tau\,(\idon{R}\otimes V)^{\dagger},
  \end{gathered}
  \end{equation*}
  and write $\mathcal{P}$ and $\iota$ for the case of a one-dimensional $R$.
  Then $\mathcal{P}_R\colon\Tcl{RM}\to\Tcl{RM^{\circ}}$ and
  $\iota_R\colon\Tcl{RM^{\circ}}\to\Tcl{RM}$ are channels,
  $\mathcal{P}_R\circ\iota_R$ is the identity, $\|\iota_R(\tau)\|_1=\|\tau\|_1$
  for every $\tau$, and $\mathcal{P}_R(x\otimes y)=x\otimes\mathcal{P}(y)$ and
  $\iota_R(x\otimes z)=x\otimes\iota(z)$ for $x\in\Tcl{R}$, $y\in\Tcl{M}$ and
  $z\in\Tcl{M^{\circ}}$; more generally, for $R=R_1R_2$,
  $\mathcal{P}_{R_1R_2}(x\otimes y)=x\otimes\mathcal{P}_{R_2}(y)$ and
  $\iota_{R_1R_2}(x\otimes z)=x\otimes\iota_{R_2}(z)$ for $x\in\Tcl{R_1}$,
  $y\in\Tcl{R_2M}$ and $z\in\Tcl{R_2M^{\circ}}$. The channel
  $\mathcal{C}_R\coloneqq\iota_R\circ\mathcal{P}_R$ on $\Tcl{RM}$ is the
  \emph{compression of $M$ onto $M^{\circ}$}.
\item \emph{(one compression)} For every state $\sigma$ on $RM$, with
  $\delta\coloneqq\tr\big[(\idon{R}\otimes\Pi^{\perp})\sigma\big]$,
  \begin{equation*}\tag{\getrefnumber{C-eq:compression-bound}}\digesteqlabel{eq:compression-bound}
    \big\|\mathcal{C}_R(\sigma)-\sigma\big\|_1\;\le\;2\sqrt{\delta}+\delta .
  \end{equation*}
  For every state $\sigma$ on $RM$ and every $\eta>0$ there is a non-zero
  finite-dimensional subspace $M^{\circ}$, spanned by finitely many
  eigenvectors of $\tr_R(\sigma)$, at which $\delta\le\eta$.
\item \emph{(a sequence of channels)} Let $Z_0,\dots,Z_m$ be systems, $\omega$
  a state on $Z_0$, $\Phi_k\colon\Tcl{Z_{k-1}}\to\Tcl{Z_k}$ channels for
  $k=1,\dots,m$, and $\mathcal{C}_k$ a channel on $\Tcl{Z_k}$ for
  $k=0,\dots,m$. Put $\sigma_0\coloneqq\omega$ and
  $\sigma_k\coloneqq\Phi_k(\sigma_{k-1})$. Then
  \begin{equation*}\tag{\getrefnumber{C-eq:compression-sequence}}\digesteqlabel{eq:compression-sequence}
    \big\|\mathcal{C}_m\Phi_m\mathcal{C}_{m-1}\cdots\Phi_1\mathcal{C}_0(\omega)
      -\sigma_m\big\|_1
    \;\le\;\sum_{k=0}^{m}\big\|\mathcal{C}_k(\sigma_k)-\sigma_k\big\|_1 .
  \end{equation*}
\item \emph{(maps on different factors)} Let $G\colon\Tcl{R}\to\Tcl{R'}$ and
  $H\colon\Tcl{M}\to\Tcl{M'}$ be linear maps, and let
  $G_M\colon\Tcl{RM}\to\Tcl{R'M}$, $G_{M'}\colon\Tcl{RM'}\to\Tcl{R'M'}$,
  $H_R\colon\Tcl{RM}\to\Tcl{RM'}$ and $H_{R'}\colon\Tcl{R'M}\to\Tcl{R'M'}$ be
  linear maps, bounded in trace norm, with $G_M(x\otimes y)=G(x)\otimes y$,
  $G_{M'}(x\otimes y')=G(x)\otimes y'$, $H_R(x\otimes y)=x\otimes H(y)$ and
  $H_{R'}(x'\otimes y)=x'\otimes H(y)$ for all $x\in\Tcl{R}$,
  $x'\in\Tcl{R'}$, $y\in\Tcl{M}$, $y'\in\Tcl{M'}$. Then
  $H_{R'}\circ G_M=G_{M'}\circ H_R$. In particular, with $H=\mathcal{P}$ or
  $H=\iota$, whose extensions are $\mathcal{P}_R$, $\mathcal{P}_{R'}$,
  $\iota_R$ and $\iota_{R'}$ by (i),
  \begin{equation*}\tag{\getrefnumber{C-eq:compression-commute}}\digesteqlabel{eq:compression-commute}
    \mathcal{P}_{R'}\circ G_M=G_{M^{\circ}}\circ\mathcal{P}_R,
    \qquad
    G_M\circ\iota_R=\iota_{R'}\circ G_{M^{\circ}} .
  \end{equation*}
\item \emph{(extension by an identity)} For every channel
  $\Phi\colon\Tcl{Z}\to\Tcl{Z'}$ and every system $K$ there is exactly one
  linear map $\Phi_K\colon\Tcl{ZK}\to\Tcl{Z'K}$, bounded in trace norm, with
  $\Phi_K(x\otimes y)=\Phi(x)\otimes y$ for all $x\in\Tcl{Z}$ and
  $y\in\Tcl{K}$; it is a channel, written $\Phi\otimes\idmap_K$.
\end{enumerate}
\end{digestclaim}

\section*{3\quad Binary64 numbers and exact sums}\pdfbookmark[1]{3 Binary64 numbers and exact sums}{digest-group-foundations/def:binary64}

% claim foundations/def:binary64
\digestlabel{def:binary64}{3.1}\pdfbookmark[2]{Definition 3.1 Binary64 numbers, declared values and dyadic lattices}{digest-foundations/def:binary64}
\renewcommand*{\RestatedName}{\hyperref{contract.pdf}{}{claim-foundations/def:binary64}{Definition~3.1}}%
\begin{digestclaim}[Binary64 numbers, declared values and dyadic lattices]
\begin{enumerate}[label=(\roman*),leftmargin=2.4em]
\item Let $\mathbb{B}$ be the set of real numbers $m\,2^{q}$ with
  $m,q\in\mathbb{Z}$, $|m|<2^{53}$ and $-1074\le q\le971$, and let
  $\mathbb{B}_{\mathbb{C}}\coloneqq\{x+iy:\ x,y\in\mathbb{B}\}$. The elements
  of $\mathbb{B}$ are the \emph{binary64 numbers}, the finite values of the
  IEEE~754 binary64 format. Every binary64 number is a dyadic rational, an
  element of $\mathbb{Z}[\tfrac12]$.
\item For a real number $x$ with $|x|<2^{1024}-2^{970}$, $\mathrm{fl}(x)$ is
  $x$ rounded to the nearest element of $\mathbb{B}$, a tie going to the
  element whose $m$ is even in its representation of (i) with the least $q$.
  For $|x|\ge2^{1024}-2^{970}$, $\mathrm{fl}(x)\coloneqq+\infty$ if $x>0$ and
  $-\infty$ if $x<0$; the threshold is the midpoint between the largest
  element $(2^{53}-1)\,2^{971}$ of $\mathbb{B}$ and $2^{1024}$, which is where
  IEEE~754 rounding to nearest overflows. Further,
  $\mathrm{fl}(\pm\infty+y)\coloneqq\pm\infty$ for $y\in\mathbb{B}$. The
  \emph{binary64 sum} of $a$
  and $b$ is $\mathrm{fl}(a+b)$. On
  $\mathbb{B}_{\mathbb{C}}$, and on matrices over it, $\mathrm{fl}$ acts on the
  real and on the imaginary part of every entry separately. Two properties of
  $\mathrm{fl}$ are all that is used of it: on real input its value lies in
  $\mathbb{B}\cup\{\pm\infty\}$, and $\mathrm{fl}(x)=x$ for $x\in\mathbb{B}$.
  The binary64 sum of $a$ and $b$ is \emph{exact} when it equals $a+b$; by
  those two properties it is exact if and only if the real and the imaginary
  part of every entry of $a+b$ lie in $\mathbb{B}$.
\item A \emph{declared value} is a binary64
  number stated as a datum. It
  denotes that dyadic rational and nothing else: there is no exact object
  behind it of which it is an approximation, and every predicate stated of
  declared values is a predicate of those rationals. A declared complex value
  is an element of $\mathbb{B}_{\mathbb{C}}$, and a declared matrix is a
  matrix of declared complex values.
\item A \emph{lattice exponent} is an integer
  $k$ with $1\le k\le53$. A real number $x$ is a \emph{lattice value} at $k$ when
  $2^{k}x\in\mathbb{Z}$ and
  $|x|\le1$; a complex number, or a matrix, is one when the real and the
  imaginary part of every entry are.
\end{enumerate}
\end{digestclaim}

% claim foundations/lem:binary64-lattice
\digestlabel{lem:binary64-lattice}{3.2}\pdfbookmark[2]{Lemma 3.2 Binary64 numbers on dyadic lattices}{digest-foundations/lem:binary64-lattice}
\renewcommand*{\RestatedName}{\hyperref{contract.pdf}{}{claim-foundations/lem:binary64-lattice}{Lemma~3.2}}%
\begin{digestclaim}[Binary64 numbers on dyadic lattices]
\begin{enumerate}[label=(\roman*),leftmargin=2.4em]
\item For every lattice exponent $k$, every real lattice value at $k$ is a
  binary64 number, and every complex lattice value at $k$ lies in
  $\mathbb{B}_{\mathbb{C}}$.
\item For every integer $e$ with $-1022\le e\le1023$,
  \begin{equation*}\tag{\getrefnumber{C-eq:binade}}\digesteqlabel{eq:binade}
    \mathbb{B}\cap\big[2^{e},2^{e+1}\big)
    \;=\;2^{e-52}\,\mathbb{Z}\cap\big[2^{e},2^{e+1}\big).
  \end{equation*}
  In particular the binary64 numbers in $[\tfrac12,1)$ are the multiples of
  $2^{-53}$ in that interval, and those in $[1,2)$ are the multiples of
  $2^{-52}$ in it.
\end{enumerate}
\end{digestclaim}

% claim foundations/def:declared-order
\digestlabel{def:declared-order}{3.3}\pdfbookmark[2]{Definition 3.3 Sums in a declared order and sums that round nowhere}{digest-foundations/def:declared-order}
\renewcommand*{\RestatedName}{\hyperref{contract.pdf}{}{claim-foundations/def:declared-order}{Definition~3.3}}%
\begin{digestclaim}[Sums in a declared order and sums that round nowhere]
Let $(x_1,\dots,x_n)$, $n\ge1$, be a finite sequence of elements of
$\mathbb{B}_{\mathbb{C}}$, or of matrices of one shape over
$\mathbb{B}_{\mathbb{C}}$ (Definition~\ref{def:binary64}\,(i)): declared values
and declared matrices (Definition~\ref{def:binary64}\,(iii)), or the computed
inner sums of a sum of sums below. Its
\emph{sequential sum} is
\begin{equation*}\tag{\getrefnumber{C-eq:seqsum}}\digesteqlabel{eq:seqsum}
  S_0\coloneqq0,\qquad S_i\coloneqq S_{i-1}+x_i\quad(1\le i\le n),
\end{equation*}
and its \emph{computed sequential sum} is $\hat S_0\coloneqq0$,
$\hat S_i\coloneqq\mathrm{fl}(\hat S_{i-1}+x_i)$ for $1\le i\le n$
(Definition~\ref{def:binary64}\,(ii)), formed left to right, entrywise, with no
re-ordering and no re-association. The accumulation \emph{rounds nowhere}
when
every one of its $n$ binary64 additions is exact: for every $1\le i\le n$, every
entry of $\hat S_{i-1}$ is finite and $\hat S_i=\hat S_{i-1}+x_i$.

A \emph{declared order} on a finite alphabet $J$ is a total order
$j_1\prec j_2\prec\dots\prec j_{|J|}$ on $J$, fixed together with $J$. For a
family $\{a_j\}_{j\in J}$ of declared values, or of declared matrices of one
shape, the \emph{sum in the declared order} is the sequential
sum of
$(a_{j_1},\dots,a_{j_{|J|}})$, and the family's sum rounds nowhere when that
accumulation does. The trace of a declared $d\times d$ matrix $X$ is
accumulated as the sequential sum of $(X_{11},X_{22},\dots,X_{dd})$, in
increasing row index. A sum of such sums over an alphabet with a declared order
rounds nowhere when every inner accumulation rounds nowhere and the sequential
sum, in that order, of the computed inner sums rounds nowhere; that sequential
sum is then the computed sum of sums. Each computed inner sum is then over
$\mathbb{B}_{\mathbb{C}}$: its last addition is exact, so it equals a sum of
finite numbers, and it is a value of $\mathrm{fl}$, which lies in
$\mathbb{B}\cup\{\pm\infty\}$ entry by entry and part by part
(Definition~\ref{def:binary64}\,(ii)). The outer sequential sum is therefore
one of the sequences above.
\end{digestclaim}

% claim foundations/lem:exact-accumulation
\digestlabel{lem:exact-accumulation}{3.4}\pdfbookmark[2]{Lemma 3.4 Exactness of a sum that rounds nowhere}{digest-foundations/lem:exact-accumulation}
\begin{mainclaim}
\renewcommand*{\RestatedName}{\hyperref{contract.pdf}{}{claim-foundations/lem:exact-accumulation}{Lemma~3.4}}%
\begin{digestclaim}[Exactness of a sum that rounds nowhere]
Let $(x_1,\dots,x_n)$ be a finite sequence of elements of
$\mathbb{B}_{\mathbb{C}}$, or of matrices of one shape over
$\mathbb{B}_{\mathbb{C}}$, such as declared values or the computed inner sums
of a sum of sums, with sequential sum $(S_i)_{i=0}^{n}$ and computed
sequential sum $(\hat S_i)_{i=0}^{n}$ (Definition~\ref{def:declared-order}).
The following are equivalent:
\begin{enumerate}[label=(\alph*),leftmargin=2.4em]
\item the accumulation rounds nowhere;
\item $\hat S_i=S_i$ for every $0\le i\le n$;
\item for every $1\le i\le n$, the real and the imaginary part of every entry
  of $S_i$ are binary64 numbers.
\end{enumerate}
If they hold, the computed sum $\hat S_n$ is the exact sum
$\sum_{i=1}^{n}x_i$.
\end{digestclaim}
\end{mainclaim}

\section*{4\quad Declared data and exact admissibility}\pdfbookmark[1]{4 Declared data and exact admissibility}{digest-group-foundations/def:declared-data}

% claim foundations/def:declared-data
\digestlabel{def:declared-data}{4.1}\pdfbookmark[2]{Definition 4.1 Declared data and the parties to a module boundary}{digest-foundations/def:declared-data}
\renewcommand*{\RestatedName}{\hyperref{contract.pdf}{}{claim-foundations/def:declared-data}{Definition~4.1}}%
\begin{digestclaim}[Declared data and the parties to a module boundary]
\begin{enumerate}[label=(\roman*),leftmargin=2.4em]
\item An \emph{instance} is the party that fixes one protocol's
  device data, protocol data and run configuration. The data it states are
  \emph{declared data}, and it is their
  \emph{declarer}.
\item A module is a \emph{supplier} when its own claims decide their
  conditions on the data it reads or emits in floating point, at slacks an
  instance declares to it, rather than exactly. A floating-point evaluation of
  a device model is an example. Its numerical output is not
  declared data until an instance states it as its own declared data, and the
  instance is then its declarer. Numbers of a supplier's output that reach
  another module in a run are declared by the instance that fixed the run:
  those it rounds and states, and those it passes on unchanged, whether in a
  declaration of its own or by handing the supplier's output on untouched.
\item A module other than a supplier that emits numerical data to another
  module as its own declaration, constructed from declared data it read, is a
  \emph{producer} of those data; they are declared data, and the
  producer is
  their declarer. Declared data that a module passes on unchanged keep their
  declarer, and the module is not their producer. A scalar bound that a
  producer emits and that a consumer has replaced by its rounding to a
  binary64 number, in a direction that the producer's own claims license as
  weakening it, is read, at its replaced value, as its producer's declaration
  of that bound, and keeps its producer as declarer; a set defined by such
  bounds is then declared data by (v).
\item A module other than a supplier that reads declared data is a
  \emph{consumer} of
  them. The passage of declared data from their declarer to a consumer is a
  \emph{module boundary}.
\item A set defined by declared data, such as a feasible set or an acceptance
  set, is declared data. An object that a module computes itself in the course
  of a computation, such as an iterate of an optimiser or the return of an
  oracle, is a \emph{computed object} and is not declared data.
\item Every entry of declared operator data is a declared value
  (Definition~\ref{def:binary64}\,(iii)).
\end{enumerate}
\end{digestclaim}

% claim foundations/def:sector-map
\digestlabel{def:sector-map}{4.2}\pdfbookmark[2]{Definition 4.2 Sector maps and block-list operators}{digest-foundations/def:sector-map}
\renewcommand*{\RestatedName}{\hyperref{contract.pdf}{}{claim-foundations/def:sector-map}{Definition~4.2}}%
\begin{digestclaim}[Sector maps and block-list operators]
A \emph{sector map} $\mathfrak{m}$ consists of two
components: a block kind,
which is a label naming what its blocks are, and a finite non-empty set
$\mathsf{S}$ of block identifiers listed in a fixed order, each identifier
$s\in\mathsf{S}$ with a dimension $d_s\ge1$. Two sector maps \emph{agree} when
they have the same block kind and the same identifiers, in the same order, with
the same dimensions. A system of dimension $d$ without a declared sector
decomposition carries the \emph{one-block map} of dimension
$d$: one
identifier, with dimension $d$, and a block kind and an identifier that are
the same for every such system, so that two one-block maps agree if and only
if their dimensions are equal.

A \emph{block list} against
$\mathfrak{m}$ is a list
$X=(X^{(s)})_{s\in\mathsf{S}}$ of complex $d_s\times d_s$ matrices. It denotes
the operator $\bigoplus_{s\in\mathsf{S}}X^{(s)}$ on
$\bigoplus_{s\in\mathsf{S}}\mathbb{C}^{d_s}$, and we identify the two: the list
is the operator, and no entry outside the blocks exists to be written, so a
block list is block diagonal by construction. A family of block lists against
$\mathfrak{m}$ is a family
$\{X_j\}_{j\in J}$ of such lists, $X^{(s)}_j$ being the block $s$ of $X_j$.

Because the direct sum acts block by block, $X=X^{\dagger}$ if and only if
$X^{(s)}=X^{(s)\dagger}$ for every $s$; $X\succeq0$ if and only if
$X^{(s)}\succeq0$ for every $s$, since
$\langle v,Xv\rangle=\sum_s\langle v_s,X^{(s)}v_s\rangle$ for $v=(v_s)_s$; and
$\sum_{j\in J}X_j=c\,\idon{}$ for a scalar $c$ if and only if
$\sum_{j\in J}X^{(s)}_j=c\,\idon{d_s}$ for every $s$.
\end{digestclaim}

% claim foundations/def:exact-admissibility
\digestlabel{def:exact-admissibility}{4.3}\pdfbookmark[2]{Definition 4.3 Exact admissibility of a declared operator family}{digest-foundations/def:exact-admissibility}
\begin{mainclaim}
\renewcommand*{\RestatedName}{\hyperref{contract.pdf}{}{claim-foundations/def:exact-admissibility}{Definition~4.3}}%
\begin{digestclaim}[Exact admissibility of a declared operator family]
Let $\mathfrak{m}$ be a sector map with blocks $s\in\mathsf{S}$ of
dimensions $d_s$ (Definition~\ref{def:sector-map}), let $J$ be a finite
non-empty alphabet with a declared order, let $c$ be a declared value,
and let
$\{a_j\}_{j\in J}$ be a family of block lists against a sector map
$\mathfrak{m}'$ whose entries are declared values
(Definition~\ref{def:binary64}\,(iii)). The family is
\emph{exactly admissible} for $\mathfrak{m}$ with constant $c$ when
\begin{itemize}[leftmargin=2.6em]
\item[(T)] \emph{(sector map)} $\mathfrak{m}'$ agrees with $\mathfrak{m}$;
\item[(H)] \emph{(Hermiticity)} $a^{(s)}_j=a^{(s)\dagger}_j$ for every
  $j\in J$ and $s\in\mathsf{S}$;
\item[(P)] \emph{(positivity)} $a^{(s)}_j\succeq0$ for every $j\in J$ and
  $s\in\mathsf{S}$;
\item[(C)] \emph{(completeness)} $\sum_{j\in J}a^{(s)}_j=c\,\idon{d_s}$ for
  every $s\in\mathsf{S}$, as an identity of the declared numbers;
\item[(S)] \emph{(the sum rounds nowhere)} for every
  $s\in\mathsf{S}$, the sum of $\{a^{(s)}_j\}_{j\in J}$ in the declared order
  of $J$ rounds nowhere (Definition~\ref{def:declared-order}).
\end{itemize}
The family is exactly admissible for $\mathfrak{m}$ when it is so with
constant $c=1$; it is then a positive operator-valued measure on
$\bigoplus_{s}\mathbb{C}^{d_s}$, by the block-by-block reading of
Definition~\ref{def:sector-map}. Since the block $s$ of a block list against
$\mathfrak{m}'$ is a square matrix of the dimension $\mathfrak{m}'$ gives $s$,
(T) holds if and only if $\mathfrak{m}'$ has the block kind and the
identifiers of $\mathfrak{m}$, in the same order, and every block $a^{(s)}_j$
is a $d_s\times d_s$ matrix, $J$ being non-empty. A family on a system without a declared sector
decomposition is read against the one-block map of that system's dimension,
and (T) then reduces to the agreement of that dimension with the size of the
family's matrices.

Each condition is a property of the declared values themselves and involves no
threshold and no norm: (T) compares declared labels and dimensions, (H) and (C)
are identities between exactly represented numbers, (P) is a property of
matrices with dyadic entries, and (S) is a property of the binary64 additions
that form the declared-order sum. Under (S) the computed
sum in the declared order is the exact sum
(Lemma~\ref{lem:exact-accumulation}), so (C) holds if and only if that
computed sum equals $c\,\idon{d_s}$ on every block.
\end{digestclaim}
\end{mainclaim}

\section*{5\quad The declared tolerance and the snap}\pdfbookmark[1]{5 The declared tolerance and the snap}{digest-group-foundations/def:tolerance}

% claim foundations/def:tolerance
\digestlabel{def:tolerance}{5.1}\pdfbookmark[2]{Definition 5.1 The declared tolerance and the defects it bounds}{digest-foundations/def:tolerance}
\renewcommand*{\RestatedName}{\hyperref{contract.pdf}{}{claim-foundations/def:tolerance}{Definition~5.1}}%
\begin{digestclaim}[The declared tolerance and the defects it bounds]
The \emph{declared tolerance} is a binary64 number
$\tnum$ with
$0\le\tnum<1$, stated by an instance in its run configuration
(Definition~\ref{def:declared-data}). For a complex matrix $Y$ we write
$\mxn{Y}\coloneqq\max_{i,l}|Y_{il}|$ for its \emph{entrywise maximum modulus}.
The \emph{defects} of data are the entrywise maximum moduli of the
differences
their predicates name:
\begin{enumerate}[label=(\roman*),leftmargin=2.4em]
\item the \emph{Hermiticity defect} $\mxn{X-X^{\dagger}}$ of a square matrix
  $X$;
\item the \emph{completeness residue}
  $\max_{s\in\mathsf{S}}\mxn{\sum_{j\in J}a^{(s)}_j-c\,\idon{d_s}}$ of a family of
  block lists $\{a_j\}_{j\in J}$ against a
  sector map with blocks $s\in\mathsf{S}$
  (Definition~\ref{def:sector-map}), read with a declared constant $c$;
\item the \emph{column residue} $\max_{l}\big|\sum_{i}P_{il}-1\big|$ of a real
  matrix $P$ read as column-stochastic;
\item the \emph{trace residue} $|\tr X-1|$ of a square matrix $X$ read as
  having unit trace, and the \emph{sum residue} $\big|\sum_{x}p_x-1\big|$ of a
  vector $p$ read as a probability vector.
\end{enumerate}
The sums here are exact sums. A square matrix $X$ passes the
\emph{positivity floor} at $\tnum$
when
$\lmin\big(\tfrac12(X+X^{\dagger})\big)\ge-\tnum$. Data lie
\emph{within the declared tolerance} when each of their defects is at most
$\tnum$ and each of their blocks read as positive semidefinite passes the
positivity floor at $\tnum$.
\end{digestclaim}

% claim foundations/lem:floor-stricter
\digestlabel{lem:floor-stricter}{5.2}\pdfbookmark[2]{Lemma 5.2 The positivity floor against entrywise closeness to the cone}{digest-foundations/lem:floor-stricter}
\renewcommand*{\RestatedName}{\hyperref{contract.pdf}{}{claim-foundations/lem:floor-stricter}{Lemma~5.2}}%
\begin{digestclaim}[The positivity floor against entrywise closeness to the cone]
Let $d\ge1$ and $\delta\ge0$, and write $\opn{\cdot}$ for the operator norm and
$\mxn{\cdot}$ for the entrywise maximum modulus of complex $d\times d$
matrices.
\begin{enumerate}[label=(\roman*),leftmargin=2.4em]
\item $\mxn{Y}\le\opn{Y}\le d\,\mxn{Y}$ for every complex $d\times d$
  matrix $Y$.
\item If $X$ is Hermitian and $\lmin(X)\ge-\delta$, then
  $X+\delta\,\idon{d}\succeq0$ and
  $\mxn{(X+\delta\,\idon{d})-X}=\delta$.
\item If $X$ is Hermitian and $\mxn{X-P}\le\delta$ for some $P\succeq0$, then
  $\lmin(X)\ge-d\delta$. The bound is attained: for the $d\times d$ matrix $E$
  with every entry equal to $1$, the matrices $X\coloneqq-\delta E$ and
  $P\coloneqq0$ satisfy $\mxn{X-P}=\delta$ and $\lmin(X)=-d\delta$.
\end{enumerate}
\end{digestclaim}

% claim foundations/def:exact-reference
\digestlabel{def:exact-reference}{5.3}\pdfbookmark[2]{Definition 5.3 Reference of a measurement family, exact or enclosed}{digest-foundations/def:exact-reference}
\renewcommand*{\RestatedName}{\hyperref{contract.pdf}{}{claim-foundations/def:exact-reference}{Definition~5.3}}%
\begin{digestclaim}[Reference of a measurement family, exact or enclosed]
Let $\mathfrak{m}$ be a sector map with blocks $s\in\mathsf{S}$ of
dimensions $d_s$ (Definition~\ref{def:sector-map}), let $J$ be a finite
non-empty alphabet with a declared order, and let $c>0$ be a
declared value (Definition~\ref{def:binary64}\,(iii)). A
\emph{entrywise enclosure} of a complex $d\times d$ matrix
$X$ is a pair $(X_c,r)$ of a Hermitian $d\times d$ matrix $X_c$ with entries
in $\mathbb{Q}(i)=\{x+iy:x,y\in\mathbb{Q}\}$, its \emph{centre}, and
rationals $r_{il}=r_{li}\ge0$, its \emph{radii}, with
\begin{equation*}\tag{\getrefnumber{C-eq:entrywise-enclosure}}\digesteqlabel{eq:entrywise-enclosure}
  \big|X_{il}-(X_c)_{il}\big|\;\le\;r_{il}\qquad\text{for all }i,l .
\end{equation*}
A \emph{reference} against $\mathfrak{m}$ over $J$
with constant $c$ is a family $G=\{G_j\}_{j\in J}$ of block lists against
$\mathfrak{m}$ whose blocks are Hermitian and positive semidefinite, with
\begin{equation*}\tag{\getrefnumber{C-eq:reference-complete}}\digesteqlabel{eq:reference-complete}
  \sum_{j\in J}G^{(s)}_j\;=\;c\,\idon{d_s}\qquad\text{for every }s\in\mathsf{S},
\end{equation*}
read together with an entrywise enclosure $\big(C^{(s)}_j,r^{(s)}_j\big)$ of
every block $G^{(s)}_j$. It is an
\emph{exact reference} when every radius is $0$; then
$G^{(s)}_j=C^{(s)}_j$ has entries in $\mathbb{Q}(i)$, and its Hermiticity,
positive semidefiniteness and completeness are properties of finitely many
Gaussian rationals.

For a Hermitian $d_s\times d_s$ matrix $Y$ with entries in $\mathbb{Q}(i)$,
the pair $\big(C^{(s)}_j-Y,\,r^{(s)}_j\big)$ is an entrywise enclosure of
$G^{(s)}_j-Y$, since the two differences in \eqref{eq:entrywise-enclosure}
are equal. A reference is an object of its own beside a family declared from
it: the entries of the declared family are declared values, and each denotes
exactly the dyadic rational it is (Definition~\ref{def:binary64}\,(iii)).
\end{digestclaim}

% claim foundations/def:factor-enclosure
\digestlabel{def:factor-enclosure}{5.4}\pdfbookmark[2]{Definition 5.4 Factor enclosure of a block}{digest-foundations/def:factor-enclosure}
\renewcommand*{\RestatedName}{\hyperref{contract.pdf}{}{claim-foundations/def:factor-enclosure}{Definition~5.4}}%
\begin{digestclaim}[Factor enclosure of a block]
Let $d\ge1$ and $k\ge0$ be integers. A
\emph{matrix box} of complex
$d\times k$ matrices is a pair $(L_c,t)$ of a $d\times k$ matrix $L_c$ with
entries in $\mathbb{Q}(i)$, its \emph{centre}, and rationals $t_{im}\ge0$,
its \emph{radii}; it \emph{contains} a complex $d\times k$ matrix $L$ when
$|L_{im}-(L_c)_{im}|\le t_{im}$ for all $i,m$. For $k=0$ it contains the one
$d\times0$ matrix.

A
\emph{factor enclosure} of a complex $d\times d$ matrix $X$ is a triple $(L_c,t,W)$ of a
matrix box $(L_c,t)$ of complex $d\times k$ matrices and a diagonal
$k\times k$ matrix $W$ with positive rational diagonal entries, its
\emph{core}, such that
\begin{equation*}\tag{\getrefnumber{C-eq:factor-enclosure}}\digesteqlabel{eq:factor-enclosure}
  X\;=\;LWL^{\dagger}\qquad\text{for some }L\text{ that }(L_c,t)\text{ contains}.
\end{equation*}
The number $k$ is its \emph{number of columns}, a matrix $L$ that the box
contains is a \emph{factor} of the triple, and the matrices $LWL^{\dagger}$
with $L$ a factor are its \emph{members}, $X$ among them; for $k=0$ the one
member is the zero $d\times d$ matrix. Every member is Hermitian and positive
semidefinite, since $x^{\dagger}LWL^{\dagger}x=\|W^{1/2}L^{\dagger}x\|^{2}$
for every $x\in\mathbb{C}^{d}$.

A reference $G$ (Definition~\ref{def:exact-reference}) \emph{carries} a
factor enclosure on a block $G^{(s)}_j$ when it is read, besides the
entrywise enclosure $\big(C^{(s)}_j,r^{(s)}_j\big)$ of that block, with
a factor enclosure of $G^{(s)}_j$; the block then satisfies both
\eqref{eq:entrywise-enclosure} and \eqref{eq:factor-enclosure}.
\end{digestclaim}

% claim foundations/def:snap
\digestlabel{def:snap}{5.5}\pdfbookmark[2]{Definition 5.5 Snapping declared data onto a lattice}{digest-foundations/def:snap}
\renewcommand*{\RestatedName}{\hyperref{contract.pdf}{}{claim-foundations/def:snap}{Definition~5.5}}%
\begin{digestclaim}[Snapping declared data onto a lattice]
Let $D$ be data a declarer is about to declare, a producer's exact
construction of what it emits (Definition~\ref{def:declared-data}), or a
reference of an operator family (Definition~\ref{def:exact-reference})
--- a family of block lists against a sector map
(Definition~\ref{def:sector-map}) over an alphabet with a declared order, a
matrix or a vector --- read with the conditions its declaration states of it,
and let $k$ be a lattice exponent
(Definition~\ref{def:binary64}\,(iv)). For a family read as a measurement,
those conditions are exact admissibility with its declared constant
(Definition~\ref{def:exact-admissibility}); every sum they read is read in its
declared order (Definition~\ref{def:declared-order}). A \emph{snap} of $D$ at
$k$ is data $D'$ of the same shape, over the same alphabet in the same order,
every entry of which is a lattice value at $k$, and which satisfy those
conditions exactly. Its \emph{displacement} is the largest entrywise modulus
of $D'-D$, taken over every block of every element
(Definition~\ref{def:tolerance}), with $D$ replaced by its centres where $D$
is a reference.

An entry of $D$ is
\emph{structurally zero} when
it is known to be exactly zero, which is the following and nothing else: where
$D$ is given as numbers, an entry that is the number $0$; where $D$ is a
reference, an entry whose centre and radius are both $0$. An entry of a
reference whose centre is $0$ and whose radius is not is not structurally
zero. The
\emph{structural support} of a block
$D^{(s)}_j$ is the set of indices $i$ such that row $i$ or column $i$ of
$D^{(s)}_j$ holds an entry that is not structurally zero. A row, a column or
an element all of whose entries are structurally zero is structurally zero,
and an index lies outside the structural support exactly when both its row
and its column are. For a reference, whose centres are Hermitian and whose
radii are symmetric (Definition~\ref{def:exact-reference}), row $i$ of a block
is structurally zero exactly when column $i$ is.

Where $D$ is an operator family read as a measurement against a sector map
with blocks $s\in\mathsf{S}$, over an alphabet $J$, a block $s$ of a snap
$F=D'$ of $D$ is
\emph{kept} when
$F^{(s)}_j=D^{(s)}_j$ is decided exactly for every $j\in J$, and
\emph{moved} otherwise. Here $D$ on block $s$ is the data the snap is made
from: the centres of an exact reference of the block where the snap
runs on them, a producer's exact construction, or the data as they stand
where the declarer states them with no model behind them. Where the
declarer computes the data of the block from a model of it and the snap
runs on that evaluation, the data the snap is made from are the model, of
which the evaluation is no exact presentation, so no exact decision equates
the snap with them and the block is moved, even where the snap leaves the
evaluation unchanged. A snap with no moved block changes nothing.

The definition fixes what a snap produces and not how. The displacement is a
record of the move, and it is the whole of what relates $D'$ to $D$ here.
\end{digestclaim}

% claim foundations/def:evaluated-off-snap
\digestlabel{def:evaluated-off-snap}{5.7}\pdfbookmark[2]{Definition 5.7 A reference handed beside a snap; a block evaluated off its snap}{digest-foundations/def:evaluated-off-snap}
\renewcommand*{\RestatedName}{\hyperref{contract.pdf}{}{claim-foundations/def:evaluated-off-snap}{Definition~5.7}}%
\begin{digestclaim}[A reference handed beside a snap; a block evaluated off its snap]
Let $F$ be a snap (Definition~\ref{def:snap}) of data $D$ against a
sector map with blocks $s\in\mathsf{S}$
(Definition~\ref{def:sector-map}), over an alphabet $J$, where $D$ is an
operator family read as a measurement, the data an instance states of it or a
reference of it, or a producer's exact construction of what it emits
(Definition~\ref{def:declared-data}).
\begin{enumerate}[label=(\roman*),leftmargin=2.4em]
\item Where $D$ is an operator family read as a measurement and the declarer
  holds a reference $G$ of it (Definition~\ref{def:exact-reference}), the
  declarer \emph{hands the reference beside the snap} on a block $s$ when it
  passes to every consumer of $F$, together with $F$ and altered by no step,
  the centre $C^{(s)}_j$, with entries in $\mathbb{Q}(i)$, and the radii
  $r^{(s)}_j$, rationals, of an entrywise enclosure of $G^{(s)}_j$ for every
  $j\in J$, with the identity of the reference and the claims that establish
  it. The reference so passed is the
  \emph{handed reference},
  and $s$ is a \emph{handed block}.
\item A block $s$ of $F$ is
  \emph{evaluated off its snap}
  in a run when the record of the run shows that no number from which the run
  computes its key length reads $F^{(s)}_j$ for any $j\in J$, save in two
  places: the points at which those numbers are evaluated, which any
  computation may have chosen, a computation on $F$ among them; and every
  value the run fixes before its first round other than the output length,
  the acceptance set and the transcript length $\lEC$ among them, whatever
  they were computed from.
\end{enumerate}
\end{digestclaim}

\section*{6\quad Domination and its limits}\pdfbookmark[1]{6 Domination and its limits}{digest-group-foundations/def:domination}

% claim foundations/def:domination
\digestlabel{def:domination}{6.1}\pdfbookmark[2]{Definition 6.1 Domination of one operator family by another at block weights}{digest-foundations/def:domination}
\renewcommand*{\RestatedName}{\hyperref{contract.pdf}{}{claim-foundations/def:domination}{Definition~6.1}}%
\begin{digestclaim}[Domination of one operator family by another at block weights]
Let $\mathfrak{m}$ be a sector map with blocks $s\in\mathsf{S}$ of
dimensions $d_s$ (Definition~\ref{def:sector-map}), let $J$ be a finite non-empty
alphabet, and let $G=\{G_j\}_{j\in J}$ and $F=\{F_j\}_{j\in J}$ be families of
block lists against $\mathfrak{m}$ whose blocks are
Hermitian and positive semidefinite; their entries are arbitrary complex
numbers. A family is \emph{complete at} a real number $c$ on block $s$ when
$\sum_{j\in J}G^{(s)}_j=c\,\idon{d_s}$. For weights $q=(q_s)_{s\in\mathsf{S}}$
with $q_s\in[0,1)$, the family $G$
\term[said of an operator family G and an operator family F over the same
alphabet and sector map, at block weights q in [0, 1): G − (1−q)F is positive
semidefinite for every outcome on every block, with the weight of that
block]{dominates} $F$ at $q$ when
\begin{equation*}\tag{\getrefnumber{C-eq:domination}}\digesteqlabel{eq:domination}
  G^{(s)}_j-(1-q_s)\,F^{(s)}_j\;\succeq\;0
  \qquad\text{for every }j\in J\text{ and }s\in\mathsf{S}.
\end{equation*}
The weight $q_s$ belongs to the block $s$ and is the same for every outcome.
On a block $s$, the
\emph{residual family} of the domination is
$R^{(s)}_j\coloneqq q_s^{-1}\big(G^{(s)}_j-(1-q_s)F^{(s)}_j\big)$ for
$j\in J$ where $q_s>0$, and $R^{(s)}_j\coloneqq G^{(s)}_j$ where $q_s=0$.
For a Hermitian matrix $X$ we write $\mathrm{supp}\,X$ for its range, the
orthogonal complement of its kernel.
\end{digestclaim}

% claim foundations/lem:domination
\digestlabel{lem:domination}{6.2}\pdfbookmark[2]{Lemma 6.2 Elementary properties of domination}{digest-foundations/lem:domination}
\renewcommand*{\RestatedName}{\hyperref{contract.pdf}{}{claim-foundations/lem:domination}{Lemma~6.2}}%
\begin{digestclaim}[Elementary properties of domination]
Let $\mathfrak{m}$, $J$ and the families $G=\{G_j\}_{j\in J}$,
$F=\{F_j\}_{j\in J}$ and $E=\{E_j\}_{j\in J}$ of block lists against
$\mathfrak{m}$ with Hermitian positive semidefinite blocks be as in
Definition~\ref{def:domination}, and let $s\in\mathsf{S}$ be a block.
Say that $G$ dominates $F$ at $q_s$ on block $s$ when
\eqref{eq:domination} holds at that block.
\begin{enumerate}[label=(\roman*),leftmargin=2.4em]
\item \emph{(Upward closure.)} If $G$ dominates $F$ at $q_s$ on block $s$ and
  $q_s\le q'_s<1$, then $G$ dominates $F$ at $q'_s$ on block $s$.
\item \emph{(Support.)} If $G$ dominates $F$ at some $q_s\in[0,1)$ on block
  $s$, then $\mathrm{supp}\,F^{(s)}_j\subseteq\mathrm{supp}\,G^{(s)}_j$ for
  every $j\in J$. Conversely, if
  $\mathrm{supp}\,F^{(s)}_j\subseteq\mathrm{supp}\,G^{(s)}_j$ for every
  $j\in J$, then $G$ dominates $F$ on block $s$ at some weight in $[0,1)$.
\item \emph{(Weight zero.)} If $G$ and $F$ are both complete at the same
  $c$ on block $s$, then $G$ dominates $F$ at $0$ on block $s$ if and only
  if $F^{(s)}_j=G^{(s)}_j$ for every $j\in J$.
\item \emph{(Composition.)} If $G$ dominates $F$ at $q_s$ and $F$ dominates
  $E$ at $q'_s$ on block $s$, then $G$ dominates $E$ on block $s$ at
  $1-(1-q_s)(1-q'_s)$, which lies in $[0,1)$ and is at most $q_s+q'_s$.
\item \emph{(Classical post-processing.)} Let $W$ be a finite non-empty
  alphabet and let $P_{w|j}\ge0$ be real numbers for $w\in W$ and $j\in J$.
  Put $G'_w\coloneqq\sum_{j\in J}P_{w|j}G_j$ and
  $F'_w\coloneqq\sum_{j\in J}P_{w|j}F_j$, block by block. If $G$ dominates
  $F$ at $q_s$ on block $s$, then $G'$ dominates $F'$ at the same $q_s$ on
  block $s$. If moreover $\sum_{w\in W}P_{w|j}=1$ for every $j$ and $G$ and
  $F$ are complete at $c$ on block $s$, then so are $G'$ and $F'$. In
  particular, multiplying each $G_j$ and $F_j$ by the same number $p_j\ge0$
  preserves domination at the same weight.
\item \emph{(Residual family.)} If $G$ and $F$ are both complete at $c$ on
  block $s$ and $G$ dominates $F$ at $q_s$ there, then the residual family
  $R^{(s)}$ of Definition~\ref{def:domination} has positive semidefinite
  elements, is complete at $c$, and satisfies
  \begin{equation*}\tag{\getrefnumber{C-eq:residual-mixture}}\digesteqlabel{eq:residual-mixture}
    G^{(s)}_j=(1-q_s)\,F^{(s)}_j+q_s\,R^{(s)}_j
    \qquad\text{for every }j\in J.
  \end{equation*}
\end{enumerate}
\end{digestclaim}

% claim foundations/lem:absorbing-domination
\digestlabel{lem:absorbing-domination}{6.3}\pdfbookmark[2]{Lemma 6.3 Domination of a family that absorbs its residue in one element}{digest-foundations/lem:absorbing-domination}
\renewcommand*{\RestatedName}{\hyperref{contract.pdf}{}{claim-foundations/lem:absorbing-domination}{Lemma~6.3}}%
\begin{digestclaim}[Domination of a family that absorbs its residue in one element]
Let $s$ be a block of dimension $d_s$, $J$ a finite non-empty alphabet,
$c>0$, and $j^{\star}\in J$. Let $\{G_j\}_{j\in J}$ be Hermitian positive
semidefinite $d_s\times d_s$ matrices with $\sum_{j\in J}G_j=c\,\idon{d_s}$,
and let $\{F_j\}_{j\in J}$ be Hermitian matrices with
\begin{equation*}\tag{\getrefnumber{C-eq:absorbing}}\digesteqlabel{eq:absorbing}
  0\preceq F_j\preceq G_j\quad(j\neq j^{\star}),
  \qquad
  F_{j^{\star}}=c\,\idon{d_s}-\sum_{j\neq j^{\star}}F_j .
\end{equation*}
Put $\Delta\coloneqq\sum_{j\neq j^{\star}}(G_j-F_j)$, the \emph{absorbed
residue}.
\begin{enumerate}[label=(\roman*),leftmargin=2.4em]
\item $\Delta\succeq0$, $F_{j^{\star}}=G_{j^{\star}}+\Delta\succeq G_{j^{\star}}$,
  every $F_j$ is positive semidefinite, and $\sum_jF_j=c\,\idon{d_s}$.
\item For $q\in[0,1)$, $\{G_j\}$ dominates $\{F_j\}$ at $q$ on this block
  (Definition~\ref{def:domination}) if and only if
  \begin{equation*}\tag{\getrefnumber{C-eq:absorbing-condition}}\digesteqlabel{eq:absorbing-condition}
    q\,G_{j^{\star}}-(1-q)\,\Delta\;\succeq\;0 .
  \end{equation*}
  For every $j\neq j^{\star}$ and every $q\in[0,1)$,
  $G_j-(1-q)F_j\succeq0$ holds.
\item Some $q\in[0,1)$ satisfies \eqref{eq:absorbing-condition} if and only if
  $\mathrm{supp}\,\Delta\subseteq\mathrm{supp}\,G_{j^{\star}}$; this holds in
  particular when $G_{j^{\star}}$ is positive definite.
\item Suppose $\mathrm{supp}\,\Delta\subseteq\mathrm{supp}\,G_{j^{\star}}$,
  let $\lambda>0$ be the least non-zero eigenvalue of $G_{j^{\star}}$ (its
  least eigenvalue when it is positive definite) and let
  $\delta\ge\opn{\Delta}$. If $G_{j^{\star}}\neq0$, then
  \eqref{eq:absorbing-condition} holds for every $q$ with
  $\delta/(\lambda+\delta)\le q<1$. If $\Delta=0$, it holds at $q=0$.
\end{enumerate}
\end{digestclaim}

% claim foundations/lem:rational-obstruction
\digestlabel{lem:rational-obstruction}{6.4}\pdfbookmark[2]{Lemma 6.4 Rationality obstruction to the domination of a lattice family}{digest-foundations/lem:rational-obstruction}
\renewcommand*{\RestatedName}{\hyperref{contract.pdf}{}{claim-foundations/lem:rational-obstruction}{Lemma~6.4}}%
\begin{digestclaim}[Rationality obstruction to the domination of a lattice family]
Let $d\ge1$, and read every vector of $\mathbb{C}^{d}$ and every
$d\times d$ matrix in the standard basis. Call a subspace $U$ of
$\mathbb{C}^{d}$ \emph{$\mathbb{Q}(i)$-rational} when it is spanned over
$\mathbb{C}$ by vectors of $\mathbb{Q}(i)^{d}$ (the subspace $\{0\}$ is
spanned by the empty set), and for a subspace $V$ of $\mathbb{C}^{d}$ put
\begin{equation*}\tag{\getrefnumber{C-eq:rational-part}}\digesteqlabel{eq:rational-part}
  \mathcal{R}(V)\;\coloneqq\;\mathrm{span}_{\mathbb{C}}\big(V\cap\mathbb{Q}(i)^{d}\big),
\end{equation*}
the \emph{largest $\mathbb{Q}(i)$-rational subspace of $V$}.
\begin{enumerate}[label=(\roman*),leftmargin=2.4em]
\item $\mathcal{R}(V)$ is a $\mathbb{Q}(i)$-rational subspace contained in
  $V$, and it contains every $\mathbb{Q}(i)$-rational subspace contained in
  $V$.
\item For every $d\times d'$ matrix $A$ with entries in $\mathbb{Q}(i)$, the
  range of $A$ is $\mathbb{Q}(i)$-rational. In particular the support of a
  Hermitian $d\times d$ matrix with entries in $\mathbb{Q}(i)$ is.
\item Let $\{G_j\}_{j\in J}$ and $\{F_j\}_{j\in J}$ be Hermitian positive
  semidefinite $d\times d$ matrices over a finite non-empty alphabet $J$,
  every $F_j$ with entries in $\mathbb{Q}(i)$, such that
  $G_j-(1-q)F_j\succeq0$ for every $j$, for some $q\in[0,1)$. Then
  $\mathrm{supp}\,F_j\subseteq\mathcal{R}_j\coloneqq
  \mathcal{R}(\mathrm{supp}\,G_j)$ for every $j$.
\item If, in (iii), moreover $\sum_jG_j=\sum_jF_j=c\,\idon{d}$ with $c>0$,
  then for every $j\in J$
  \begin{equation*}\tag{\getrefnumber{C-eq:rational-bound}}\digesteqlabel{eq:rational-bound}
    q\;\ge\;\frac{\lmax\big(P_jG_jP_j\big)}{c},
  \end{equation*}
  where $P_j$ is the orthogonal projector onto $\mathcal{R}_j^{\perp}$.
\item In particular, if $\mathrm{supp}\,G_j$ contains no non-zero vector of
  $\mathbb{Q}(i)^{d}$, then $q\ge\lmax(G_j)/c$, which equals $\tr G_j/c$
  when $G_j$ has rank one.
\end{enumerate}
\end{digestclaim}

% claim foundations/lem:enclosure-psd
\digestlabel{lem:enclosure-psd}{6.5}\pdfbookmark[2]{Lemma 6.5 Positive semidefiniteness from an entrywise enclosure}{digest-foundations/lem:enclosure-psd}
\renewcommand*{\RestatedName}{\hyperref{contract.pdf}{}{claim-foundations/lem:enclosure-psd}{Lemma~6.5}}%
\begin{digestclaim}[Positive semidefiniteness from an entrywise enclosure]
\begin{enumerate}[label=(\roman*),leftmargin=2.4em]
\item Let $X$ be a Hermitian complex $d\times d$ matrix, $X_c$ a Hermitian
  $d\times d$ matrix with entries in $\mathbb{Q}(i)$, and $r_{il}\ge0$
  rationals with $|X_{il}-(X_c)_{il}|\le r_{il}$ for all $i,l$. Let $\rho$ be
  a rational with
  \begin{equation*}\tag{\getrefnumber{C-eq:enclosure-radius}}\digesteqlabel{eq:enclosure-radius}
    \rho\;\ge\;\max_{i}\sum_{l}r_{il}
    \qquad\text{or}\qquad
    \rho\ge0\ \text{ and }\ \rho^{2}\;\ge\;\sum_{i,l}r_{il}^{2},
  \end{equation*}
  and $\mu\ge0$ a rational. If $X_c-(\rho+\mu)\,\idon{d}\succeq0$, then
  $X\succeq\mu\,\idon{d}$; in particular $X\succeq0$, and $X$ is positive
  definite when $\mu>0$. If $X_c-\rho\,\idon{d}$ is positive definite, then
  $X$ is positive definite. Every hypothesis on $X_c$, $r$, $\rho$ and $\mu$
  is a property of finitely many rationals and Gaussian rationals.
\item Let $n_1,\dots,n_d$ be positive rationals and
  $D\coloneqq\mathrm{diag}(\sqrt{n_1},\dots,\sqrt{n_d})$, and let $a_i,e_i$
  be rationals with $0\le e_i\le a_i$ and
  \begin{equation*}\tag{\getrefnumber{C-eq:sqrt-enclosure}}\digesteqlabel{eq:sqrt-enclosure}
    (a_i-e_i)^{2}\,n_i\;\le\;1\;\le\;(a_i+e_i)^{2}\,n_i .
  \end{equation*}
  Then $|a_i-n_i^{-1/2}|\le e_i$. Let $S$ and $Y$ be Hermitian $d\times d$
  matrices with entries in $\mathbb{Q}(i)$, put $X\coloneqq D^{-1}SD^{-1}-Y$,
  $A\coloneqq\mathrm{diag}(a_1,\dots,a_d)$, $X_c\coloneqq ASA-Y$ and
  \begin{equation*}\tag{\getrefnumber{C-eq:scaled-radii}}\digesteqlabel{eq:scaled-radii}
    r_{il}\;\coloneqq\;\sigma_{il}\,\big(a_ie_l+e_ia_l+e_ie_l\big),
    \qquad
    \sigma_{il}\coloneqq|\mathrm{Re}\,S_{il}|+|\mathrm{Im}\,S_{il}| .
  \end{equation*}
  Then $X$, $X_c$ and $r$ satisfy the hypotheses of (i) on $X$, $X_c$ and
  $r$, and the radii do not depend on $Y$.
\item Let $X$, $X_c$ and $r$ be as in (i), let $F$ be a Hermitian positive
  semidefinite $d\times d$ matrix with entries in $\mathbb{Q}(i)$, and let
  $q\in[0,1]$ be rational. Then $X-(1-q)F$, $X_c-(1-q)F$ and $r$ satisfy the
  hypotheses of (i) on $X$, $X_c$ and $r$. If, for some $\rho$ and $\mu$ as
  in (i), $X_c-(1-q)F-(\rho+\mu)\idon{d}\succeq0$, then
  $X_c-(1-q')F-(\rho+\mu)\idon{d}\succeq0$ for every rational
  $q'\in[q,1]$.
\end{enumerate}
\end{digestclaim}

\section*{7\quad The exact island, enclosures and the tiers}\pdfbookmark[1]{7 The exact island, enclosures and the tiers}{digest-group-foundations/def:exact-island}

% claim foundations/def:exact-island
\digestlabel{def:exact-island}{7.1}\pdfbookmark[2]{Definition 7.1 The exact island}{digest-foundations/def:exact-island}
\renewcommand*{\RestatedName}{\hyperref{contract.pdf}{}{claim-foundations/def:exact-island}{Definition~7.1}}%
\begin{digestclaim}[The exact island]
The \emph{exact island} is
the following set of computations of the engine,
the declarations of its instances included.
\begin{enumerate}[label=(\roman*),leftmargin=2.4em]
\item \emph{The admission boundary}: every decision that data satisfy a
  condition their declaration states of them --- among them the conditions of
  Definition~\ref{def:exact-admissibility}, a trace, a column sum, an identity
  between declared operators and a lattice condition --- taken by a consumer
  of declared data it reads, or by a declarer, an instance or a producer
  (Definition~\ref{def:declared-data}), of data it is about to declare; the
  decisions by which a snap establishes the conditions of its output
  (Definition~\ref{def:snap}) are among the latter. So are, for a
  reference of an operator family that a declarer uses or hands beside
  its snap (Definition~\ref{def:exact-reference}), the computation of the
  centres and radii of its entrywise enclosures,
  whose enclosure of the reference, \eqref{eq:entrywise-enclosure}, the claims
  of the module that constructs the reference establish, the decisions on
  those centres and radii and, for an exact reference, on its
  positivity and completeness; and, for any snap of an operator family read
  as a measurement, the decisions by which the snap keeps a block, and the
  decisions on which entries are structurally zero. It
  contains, further, a
  producer's construction, block by block, of the operator data it emits, and,
  for each scalar bound that a consumer replaces as in
  Definition~\ref{def:declared-data}\,(iii), the decision that the
  replacement lies on the weakening side of the bound emitted.
\item \emph{The acceptance side of the optimisation}, the honest model and
  the completeness budget being declared as exact rational data: the
  construction of an acceptance set from them; the exact establishment that
  the optimiser's acceptance-side iterate lies in the acceptance set it reads,
  for its initial point and for each update as a convex combination, with
  exact rational weights, of points so established; and the exact
  certification of the minimum of a linear functional over that set, or of a
  lower bound on that minimum. The acceptance set the optimiser reads is the
  set so declared: the constructed set as its producer emits it, with each
  scalar bound that a consumer has replaced as in
  Definition~\ref{def:declared-data}\,(iii) at its replaced value.
\item \emph{The feasible side of the optimisation}: the certification of a
  lower bound on the minimum of a linear functional over the set the
  optimiser's linear oracle is asked to minimise over, formed from the
  functional, from the data of that set and from multipliers, each given to
  it as values, together with the bounds on least eigenvalues that the
  certification computes on the way.
\end{enumerate}
The decision whether data an instance is about to declare lie
within the declared tolerance (Definition~\ref{def:tolerance}) is not on
the exact island. The optimiser's iterate apart from
its acceptance side, the optimiser's oracle, the objective and its gradient lie
outside the exact island, and so does a binary64 computation that only
proposes a candidate, such as a point or a multiplier, for a decision or a
certification of the island. A certification of (iii) holds of the
functional and the data as they are given to it, and not of the quantities
that a binary64 computation of them approximates.
\end{digestclaim}

% claim foundations/lem:least-eigenvalue-bound
\digestlabel{lem:least-eigenvalue-bound}{7.2}\pdfbookmark[2]{Lemma 7.2 Least-eigenvalue lower bound from an approximate eigenbasis}{digest-foundations/lem:least-eigenvalue-bound}
\renewcommand*{\RestatedName}{\hyperref{contract.pdf}{}{claim-foundations/lem:least-eigenvalue-bound}{Lemma~7.2}}%
\begin{digestclaim}[Least-eigenvalue lower bound from an approximate eigenbasis]
Let $n\ge1$, let $R$ be a Hermitian complex $n\times n$ matrix, and let
\begin{equation*}\tag{\getrefnumber{C-eq:real-embedding}}\digesteqlabel{eq:real-embedding}
  \widetilde R\;\coloneqq\;
  \begin{pmatrix}\mathrm{Re}\,R&-\mathrm{Im}\,R\\
  \mathrm{Im}\,R&\mathrm{Re}\,R\end{pmatrix}\in\mathbb{R}^{2n\times2n},
\end{equation*}
the real and imaginary parts taken entrywise. Let $V$ be any real
$2n\times2n$ matrix, put
\begin{equation*}
  B\coloneqq V^{\transpose}\widetilde R\,V,\qquad
  E\coloneqq V^{\transpose}V-\idon{2n},\qquad
  \gamma\coloneqq\min_{1\le i\le2n}\Big(B_{ii}-\sum_{l\ne i}|B_{il}|\Big),
\end{equation*}
and let $\theta$ be a real number with $0\le\theta<1$ and
$\sum_{i,l}E_{il}^{2}\le\theta^{2}$.
\begin{enumerate}[label=(\roman*),leftmargin=2.4em]
\item $\widetilde R$ is real symmetric, and $\lmin(\widetilde R)=\lmin(R)$.
\item $V$ is invertible, and
\begin{equation*}\tag{\getrefnumber{C-eq:least-eigenvalue-bound}}\digesteqlabel{eq:least-eigenvalue-bound}
  \lmin(R)\;\ge\;
  \begin{cases}
    \gamma/(1+\theta), & \gamma\ge0,\\
    \gamma/(1-\theta), & \gamma<0.
  \end{cases}
\end{equation*}
\item If the entries of $R$ lie in $\mathbb{Q}(i)$, the entries of $V$ are
  rational and $\theta$ is rational, then $\widetilde R$, $B$ and $E$ have
  rational entries, $\gamma$ is rational, the hypothesis on $\theta$ is a
  comparison of rationals, and the right-hand side of
  \eqref{eq:least-eigenvalue-bound} is a rational number.
\item Applied to $-R$, (ii) bounds the largest eigenvalue of $R$, which is
  $-\lmin(-R)$, from above.
\end{enumerate}
\end{digestclaim}

% claim foundations/def:tiers
\digestlabel{def:tiers}{7.6}\pdfbookmark[2]{Definition 7.6 The tiers of a reported key rate}{digest-foundations/def:tiers}
\renewcommand*{\RestatedName}{\hyperref{contract.pdf}{}{claim-foundations/def:tiers}{Definition~7.6}}%
\begin{digestclaim}[The tiers of a reported key rate]
A \emph{tier} of a reported key-rate number is one
of the following three labels.
\begin{enumerate}[label=(\roman*),leftmargin=2.4em]
\item \emph{estimate}: the
  number is the value of the objective at the point
  an optimiser returned. It is reported as an estimate of the minimum of the
  objective and is claimed as no bound.
\item \emph{first-order}: the number is a lower bound on the minimum of the
  objective, established by a claim under hypotheses that claim names.
\item \emph{certified}: the number is a lower bound on the minimum
  of the objective of the declaration it is formed on, established by a claim
  under hypotheses that claim names, and every arithmetic step between that
  declaration's data and the number is exact or
  \emph{rigorously enclosed}: a
  computation is rigorously enclosed when each of its operations returns a
  set of numbers, such as an interval or a ball, that contains the exact
  result of the operation at every value the sets of its inputs contain, so
  that the set it ends with contains the exact value, at the exact data, of
  the quantity it computes. None of those hypotheses is an assumption about
  how a number was computed, that is, none is an assumption of the category
  numerical-trust.
\end{enumerate}
The tiers are ordered, estimate below first-order below certified, and a
tier is \emph{stronger} than the tiers below it. A certified number is a
first-order number as well. There is no other tier.
\end{digestclaim}

\section*{8\quad The snapped-measurement assumption}\pdfbookmark[1]{8 The snapped-measurement assumption}{digest-group-foundations/ass:snapped-measurement}

% claim foundations/ass:snapped-measurement
\digestlabel{ass:snapped-measurement}{8.1}\pdfbookmark[2]{Assumption 8.1 The device measures the snapped family where the run reads the snap}{digest-foundations/ass:snapped-measurement}
\SetShortId{A-SNAP}
\renewcommand*{\RestatedName}{\hyperref{contract.pdf}{}{claim-foundations/ass:snapped-measurement}{Assumption~8.1}}%
\begin{digestassumption}[The device measures the snapped family where the run reads the snap]
Let $F$ be a snap (Definition~\ref{def:snap}) of an operator family $D$ read
as a measurement, against a sector map with blocks $s\in\mathsf{S}$
(Definition~\ref{def:sector-map}), over an alphabet $J$, with $D$ the data the
snap is made from: the data as an instance states them, or, where the
instance computes those data from a model of the family, that model, a
reference of the family (Definition~\ref{def:exact-reference}), or a
producer's exact construction of what it emits
(Definition~\ref{def:declared-data}). Let $B\subseteq\mathsf{S}$ be the set of
blocks that the record of the run names as resting on this assumption: the
moved blocks of $F$ (Definition~\ref{def:snap}), those on which the
record does not show, by an exact decision, that $F^{(s)}_j=D^{(s)}_j$ for
every $j\in J$, that are not evaluated off its snap
(Definition~\ref{def:evaluated-off-snap}). A kept block and a block
evaluated off its snap do not belong to $B$. Every other moved block belongs
to $B$, whatever its cause; among them are the blocks on which the snap gives
an element support outside the structural support of $D^{(s)}_j$
(Definition~\ref{def:snap}).
\begin{enumerate}[label=(\roman*),leftmargin=2.4em]
\item Where $D$ is the data an instance states as the measurement of a
  device, the model it computes them from, or a reference of it, then on every block $s\in B$ the block on $s$
  of the measurement the device performs, read against the sector map, is
  the family $\{F^{(s)}_j\}_{j\in J}$ with the entries recorded for the snap.
  On such a block, any statement of another assumption of the run that a
  model describes that measurement is not asserted, and this clause holds in
  its place.
\item Where $D$ is a producer's exact construction of what it emits, every
  conclusion about the device that the claims of the run draw from the
  emitted family, under their hypothesis that it is the producer's exact
  construction, holds when the emitted family is $F^{(s)}$ on every block
  $s\in B$.
\end{enumerate}
\end{digestassumption}

\part*{Engine requirements}\pdfbookmark[0]{Engine requirements}{digest-part-2}

\section*{9\quad Admission}\pdfbookmark[1]{9 Admission}{digest-group-foundations/req:boundary-admission}

% claim foundations/req:boundary-admission
\digestlabel{req:boundary-admission}{9.1}\pdfbookmark[2]{Requirement 9.1 Exact admissibility at every module boundary}{digest-foundations/req:boundary-admission}
\begin{mainclaim}
\renewcommand*{\RestatedName}{\hyperref{contract.pdf}{}{claim-foundations/req:boundary-admission}{Requirement~9.1}}%
\begin{digestrequirement}[Exact admissibility at every module boundary]
Declared operator data are exactly admissible at every
module boundary (Definition~\ref{def:declared-data}). Precisely:
\begin{enumerate}[label=(\roman*),leftmargin=2.4em]
\item every declared operator family that a consumer reads as a measurement,
  that is, as a positive operator-valued measure or as a family with a
  declared completeness constant, is exactly admissible with that constant for
  the sector map the consumer holds
  (Definition~\ref{def:exact-admissibility});
\item every other declared operator, and every declared matrix or vector
  whose entries are declared values, satisfies exactly,
  on those values, each condition its declaration states of it, such as
  Hermiticity, positivity, a trace, a column sum or an identity with another
  declared operator, and every sum such a condition reads rounds nowhere in
  its declared order (Definition~\ref{def:declared-order}).
\end{enumerate}
This is established before the data cross the boundary. An instance
establishes it by constructing or snapping what it declares
(Requirement~\ref{req:who-snaps}), or by passing on unchanged, from a
supplier's output (Definition~\ref{def:declared-data}\,(ii)), data whose
conditions the consumer decides exactly on entry, that decision then
establishing it. A producer establishes it by constructing exactly what it
emits and by deciding exactly, before emitting, each condition that its
construction does not already establish, a sum that rounds nowhere among them;
for a measurement family that decision is the one of
Requirement~\ref{req:exact-emission}. Data that fail it do not cross, because
a consumer that finds one of its admission predicates violated refuses them
(Requirement~\ref{req:consumer-decides-exactly}).
\end{digestrequirement}
\end{mainclaim}

% claim foundations/req:consumer-decides-exactly
\digestlabel{req:consumer-decides-exactly}{9.2}\pdfbookmark[2]{Requirement 9.2 Exact decision by every consumer}{digest-foundations/req:consumer-decides-exactly}
\renewcommand*{\RestatedName}{\hyperref{contract.pdf}{}{claim-foundations/req:consumer-decides-exactly}{Requirement~9.2}}%
\begin{digestrequirement}[Exact decision by every consumer]
Every consumer (Definition~\ref{def:declared-data}) decides every admission
predicate it states of declared data exactly, on the declared values,
and none
more loosely than exactly: no consumer has an input through which a predicate
on declared data would be decided at a nonzero defect, and none admits data on
the ground that they lie within a tolerance of data it would admit. A consumer
that finds such a predicate violated refuses, with a machine-readable reason
naming the predicate that failed. It repairs nothing it was handed, and it
neither snaps nor otherwise alters declared data of another declarer, with one
exception: it may replace a scalar bound that a producer emits by that bound
rounded to a binary64 number in a direction that the producer's own claims
license as weakening it. The direction in which each bound so replaced was
rounded is recorded where the consumer's reading of the emission is
specified, nothing else of the emission is altered, and every predicate on
such a bound is decided exactly on its replaced value. The
declared tolerance therefore widens nothing that any consumer admits;
it bounds only how far an instance may move its own declaration
(Requirement~\ref{req:declared-tolerance}).
\end{digestrequirement}

% claim foundations/req:declared-order
\digestlabel{req:declared-order}{9.3}\pdfbookmark[2]{Requirement 9.3 Summation order declared with the data}{digest-foundations/req:declared-order}
\renewcommand*{\RestatedName}{\hyperref{contract.pdf}{}{claim-foundations/req:declared-order}{Requirement~9.3}}%
\begin{digestrequirement}[Summation order declared with the data]
Where a predicate on declared data is a sum, the order in which the sum is
accumulated is part of the declared data: it is declared with the alphabet
that indexes the sum, as that alphabet's declared order
(Definition~\ref{def:declared-order}), or with an alphabet containing it,
whose declared order it then carries by restriction; and it travels with the
artefact that carries the data. Where a declared order arrives both with the
alphabet and with an alphabet containing it, the two are one order: the first
is the restriction of the second. A consumer that receives two such orders
that differ takes the order of the containing alphabet, decides the predicate
in its restriction, and refuses nothing for the difference. Every consumer that decides the predicate,
and every declarer that decides it of data it is about to declare
(Definition~\ref{def:declared-data}), decides it in that order, and no
consumer substitutes an order of its own; a supplier deciding its own
conditions at slacks declared to it is not bound by this. Only where no
declared order arrives with the alphabet or with an alphabet containing it is
the declaration's own listing of the alphabet taken as its declared order.
Where such a predicate reads the trace of a declared matrix, the trace is
accumulated in increasing row index.
\end{digestrequirement}

\section*{10\quad The tolerance, norms and snapping}\pdfbookmark[1]{10 The tolerance, norms and snapping}{digest-group-foundations/req:declared-tolerance}

% claim foundations/req:declared-tolerance
\digestlabel{req:declared-tolerance}{10.1}\pdfbookmark[2]{Requirement 10.1 One declared tolerance per instance}{digest-foundations/req:declared-tolerance}
\renewcommand*{\RestatedName}{\hyperref{contract.pdf}{}{claim-foundations/req:declared-tolerance}{Requirement~10.1}}%
\begin{digestrequirement}[One declared tolerance per instance]
Each instance declares exactly one declared tolerance $\tnum$
(Definition~\ref{def:tolerance}) in its run configuration, and $\tnum$ has no
default: a run whose configuration does not declare it, or declares a value
that is not a binary64 number in $[0,1)$, is refused before anything is
constructed. $\tnum$ is the ceiling on the defect an instance may repair in
what it declares. It binds the data an instance is about to declare other than
numbers of a supplier's output that it passes on unchanged
(Definition~\ref{def:declared-data}\,(ii)); those it neither measures nor
repairs, and their consumer decides their conditions exactly on entry
(Requirement~\ref{req:boundary-admission}). Of the data it binds that have no
exact reference on whose centres the snap runs (for those, see the
next paragraph), the instance
measures in binary64 the defects (Definition~\ref{def:tolerance}) and, for each
block read as positive semidefinite, the least eigenvalue of its Hermitian
part, each on the data as they stand before any step of the snap, and records
those measured numbers. Data with a measured defect above $\tnum$ or a
measured least eigenvalue below $-\tnum$ are refused as a modelling error; the
others may be snapped (Definition~\ref{def:snap}), the snap being recorded
with the measured defects it corrected (Requirement~\ref{req:who-snaps}). Near
$\tnum$ this decision on measured numbers may differ from the predicate
\emph{within the declared tolerance} of Definition~\ref{def:tolerance},
whose sums are exact, and no bound on that difference is kept. This decision
is not on the exact island; the exact decisions are those by which the
snap establishes the conditions of its output
(Definition~\ref{def:exact-island}\,(i)).

An operator family that an instance snaps on the centres of an
exact reference (Definition~\ref{def:exact-reference}), as
Requirement~\ref{req:snap-reference}\,(ii) states, stands before the snap as
that reference, which is Hermitian, positive semidefinite and complete by the
exact decision on its entries, and never by a binary64 evaluation. Its defects are
therefore zero and no block of it falls below the positivity floor:
$\tnum$ decides nothing for it, and the instance records the defects as zero
on that ground rather than on a binary64 measurement. $\tnum$ keeps its role
for the data it binds on which the snap runs as they stand, an evaluation of
a model whose reference has a radius that is not zero included, since the
snap does not run on such a reference.

A probability vector that an instance prepares from ordinary numerical inputs
before declaring it is decided instead by the rule of the input preparation.
It is refused unless its normalisation defect, computed exactly on the inputs
as given, is at most $\tnum$, or the sum of the rounding cells of its inputs
contains $1$: a binary64 input contributes the rounding cell it denotes under
the preparation's documented representation policy, and an input given
exactly, as a decimal or a fraction, contributes the one point that is its
value. A value below $0$ or above $1$ is refused. The preparation records the
defect, the cells and which of the two grounds admitted the vector. Its snap
then runs at a \emph{guard}, the larger of $\tnum$ and the completeness
residue the snap measures on the binary64 conversion of the inputs, which may
exceed $\tnum$; the guard is recorded, and it admits no vector the rule has
refused.
\end{digestrequirement}

% claim foundations/req:tolerance-scope
\digestlabel{req:tolerance-scope}{10.2}\pdfbookmark[2]{Requirement 10.2 Scope of the declared tolerance}{digest-foundations/req:tolerance-scope}
\renewcommand*{\RestatedName}{\hyperref{contract.pdf}{}{claim-foundations/req:tolerance-scope}{Requirement~10.2}}%
\begin{digestrequirement}[Scope of the declared tolerance]
The declared tolerance governs the repair of an instance's own
declared data (Requirement~\ref{req:declared-tolerance}) and nothing
else; in particular no consumer reads it
(Requirement~\ref{req:consumer-decides-exactly}).
\begin{enumerate}[label=(\roman*),leftmargin=2.4em]
\item A set defined by declared data is the exact set those data define, and
  the declared tolerance relaxes none of its defining conditions. A consumer
  that certifies a supplied point as a member of such a set certifies
  membership in the exact set, deciding it exactly on the declared values of
  the point. A check of a supplied point at a window of the module's own
  certifies no membership in such a set, and its window is a number of that
  module, as in (ii).
\item The declared tolerance does not govern a module's arithmetic on
  computed objects (Definition~\ref{def:declared-data}):
  a window at which a module checks a computed object, or within which it
  repairs one, is a number that module declares and records in its own right,
  and it is never the declared tolerance.
\item The declared tolerance is a numerical tolerance and not a
  device-uncertainty budget. A declared number that bounds a physical
  quantity, such as a tail budget, is declared separately and is never derived
  from the declared tolerance. A slack that an instance declares to a supplier
  (Definition~\ref{def:declared-data}\,(ii)) for the supplier's own conditions
  is a declared number of its own. One that bounds the error of the
  supplier's floating-point evaluation of a quantity, such as a slack on a
  computed eigenvalue (Requirement~\ref{req:exact-island}), is never derived
  from the declared tolerance; any other is derived from it, if at all, only
  as Requirement~\ref{req:norm-conversion} says.
\end{enumerate}
\end{digestrequirement}

% claim foundations/req:norm-conversion
\digestlabel{req:norm-conversion}{10.3}\pdfbookmark[2]{Requirement 10.3 Explicit conversion between norms}{digest-foundations/req:norm-conversion}
\renewcommand*{\RestatedName}{\hyperref{contract.pdf}{}{claim-foundations/req:norm-conversion}{Requirement~10.3}}%
\begin{digestrequirement}[Explicit conversion between norms]
No bound stated in one norm is read as a bound in another. The defects of
Definition~\ref{def:tolerance} are entrywise maximum moduli and the
positivity floor is a bound on a least eigenvalue, and neither is
converted, implicitly, into the other or into a bound in any other norm.
A declared number that is
derived from the declared tolerance, or from any bound stated in one norm, for
use as a bound in another norm is derived explicitly by the party that
declares it, with the conversion factor between the two norms concerned
written down (for the operator norm and the entrywise maximum modulus, that of
Lemma~\ref{lem:floor-stricter}\,(i)), and it is recorded together with that
factor.
\end{digestrequirement}

% claim foundations/req:who-snaps
\digestlabel{req:who-snaps}{10.4}\pdfbookmark[2]{Requirement 10.4 Snapping by the declarer only}{digest-foundations/req:who-snaps}
\renewcommand*{\RestatedName}{\hyperref{contract.pdf}{}{claim-foundations/req:who-snaps}{Requirement~10.4}}%
\begin{digestrequirement}[Snapping by the declarer only]
Declared data are snapped (Definition~\ref{def:snap}) only by their declarer,
at most once, and before any consumer reads them: an instance snaps what it
declares about its device and its protocol, on the admission that
Requirement~\ref{req:declared-tolerance} states (for an operator family
snapped on the centres of an exact reference, the decisions on the
reference that Requirement~\ref{req:snap-reference}\,(i) lists; for a
probability vector prepared from numerical inputs, the rule of the input
preparation; otherwise within the declared tolerance), and a producer snaps
only its own emission. No party snaps data it was handed, and no party
otherwise alters them, with the one exception that
Requirement~\ref{req:consumer-decides-exactly} states and limits: a consumer
may replace a scalar bound that a producer emits by that bound rounded in a
direction that the producer's own claims license as weakening it, recording
the direction. Every snap is recorded: an instance's in the record of the run,
with the defects it measured, or decided to be zero, before the snap
(Definition~\ref{def:tolerance}),
the displacement and the rule it applied, and a producer's in its emission,
with the disclosures of Requirement~\ref{req:exact-emission}. The record of a
snap of an operator family read as a measurement carries, besides, the
identity of every reference the declarer uses or hands beside the snap
(the model, the exact parameters from which it is formed, and its
entrywise enclosures where any radius is not
zero, with the claims that establish them;
Definition~\ref{def:exact-reference}); the \emph{reference kind} of the snap
per block (\emph{exact}, where the snap runs on the centres of an exact
reference of the block, or \emph{none}, where it runs on the data as they
stand); whether the run declares that enclosed references are not used; the
status of every block, kept or moved
(Definition~\ref{def:snap}); the steps of the procedure that made each moved
block (Requirement~\ref{req:support-safe-snap}); and the flag
\iface{support_not_preserved} of every block snapped by the last resort
(Requirement~\ref{req:support-safe-snap}\,(iv)). Where the declarer hands the
reference beside the snap (Definition~\ref{def:evaluated-off-snap}\,(i)), the
record names the handed blocks, with the kind of the reference handed on each
(\emph{exact}, every radius zero, or \emph{enclosed}, with the largest radius
of the block and the largest radius of the enclosures of irrational inputs
(square roots, trigonometric values) from which the constructing module
computed the centres exactly, as its claims state it), and, for every block
on which it hands none, the cause of
Requirement~\ref{req:snap-reference}\,(iii)(a). The record of the run names
every moved block of the snap that is evaluated off its snap
(Definition~\ref{def:evaluated-off-snap}\,(ii)), and every moved block that
rests on Assumption~\ref{ass:snapped-measurement}, with its cause, group and
fix (Requirement~\ref{req:snap-reference}\,(iii)); a kept block is named as
kept and in neither list. Handing the reference alters neither the snap nor
the reference. A producer that declares a lattice exponent
re-establishes its lattice at each emission whose exact construction lies off
it, by a snap of that emission or by refusing it
(Requirements~\ref{req:declared-lattice} and~\ref{req:exact-emission}); a
producer that declares no lattice exponent snaps nothing.
\end{digestrequirement}

% claim foundations/req:snap-reference
\digestlabel{req:snap-reference}{10.5}\pdfbookmark[2]{Requirement 10.5 Snaps of a measurement family: kept, or moved and disclosed}{digest-foundations/req:snap-reference}
\renewcommand*{\RestatedName}{\hyperref{contract.pdf}{}{claim-foundations/req:snap-reference}{Requirement~10.5}}%
\begin{digestrequirement}[Snaps of a measurement family: kept, or moved and disclosed]
\begin{enumerate}[label=(\roman*),leftmargin=2.4em]
\item \emph{The reference.} Before it snaps an operator family read as a
  measurement, an instance of what it declares or a producer of its own
  emission (Definition~\ref{def:declared-data}), the declarer decides
  exactly, of every reference of the family it uses or hands beside
  the snap (Definition~\ref{def:exact-reference}), that every centre is
  Hermitian with entries in $\mathbb{Q}(i)$ and that the radii are
  non-negative rationals with $r_{il}=r_{li}$, of every factor
  enclosure it carries (Definition~\ref{def:factor-enclosure}) that the
  centre of its matrix box has entries in $\mathbb{Q}(i)$, that its
  radii are non-negative rationals and that its core is diagonal with
  positive rational entries, and, for an exact reference, that it is
  positive semidefinite and complete at its constant. For a reference with a
  radius that is not zero, three properties are not decided by the declarer:
  that its enclosures enclose it, \eqref{eq:entrywise-enclosure}, and that it
  is positive semidefinite and complete at its constant. They are
  established by the claims of the module that constructs the reference,
  which the record names, and its centres and radii are computed in exact
  rational arithmetic on the exact island
  (Definition~\ref{def:exact-island}\,(i)); a binary64 evaluation
  establishes none of them. A reference of which any of them is not so
  established is neither used nor handed. Nor is it decided by the
  declarer, of any reference, that a block on which it carries a factor
  enclosure is a member of it, \eqref{eq:factor-enclosure}: that too is
  established by the claims of the module that constructs the reference,
  which the record names, and the centres and radii of the factor enclosure are computed in exact
  rational arithmetic or by a rigorously enclosed computation
  (Definition~\ref{def:tiers}). A factor enclosure of which either is not so
  established is not carried, and the reference is read without it. A producer's exact
  construction, where it is exactly Hermitian, positive semidefinite and
  complete, is an exact reference of its emission. An instance that holds a
  reference so established hands it beside its snap
  (Definition~\ref{def:evaluated-off-snap}\,(i)) on every block, save that,
  in a run whose configuration declares that enclosed references are not
  used, by the option \iface{reference_enclosure}, it hands only the blocks
  every radius of which is zero.
\item \emph{The snap.} Every such snap is made block by block. A block is
  kept, $F^{(s)}=G^{(s)}$, where the declarer establishes an exact
  reference $G$ of that block, against the sector map of the block
  alone, decided as (i) states, every element of which is a
  lattice value at the lattice exponent $k$
  (Definition~\ref{def:binary64}\,(iv)), and where exact admissibility at the
  constant $c$ (Definition~\ref{def:exact-admissibility}) and the lattice
  condition at $k$, decided last on the block, hold. For a producer the
  block kept is the one its own claims keep. Every other block is made by a
  procedure that meets Requirement~\ref{req:support-safe-snap} and aims at
  fidelity to the data it starts from rather than at any other property: an
  instance's is one procedure for every such block, the support-safe steps
  specified with the shared numerical code, run on the centres of an exact
  reference of the block where the declarer establishes one, and otherwise
  on the data as they stand, admitted within the declared tolerance; a
  producer's is the procedure its own claims state. Where that procedure
  yields on a block no data that satisfy the conditions of
  Definition~\ref{def:snap}, the block is snapped by the last resort of
  Requirement~\ref{req:support-safe-snap}\,(iv), which the record flags;
  where the last resort too yields no such data, the declarer refuses the
  data, naming the block and the condition that failed. Every block made by
  this procedure or the last resort is moved
  (Definition~\ref{def:snap}) unless the snap is decided exactly to equal on
  it the data it is made from.
\item \emph{Resting.} A moved block \emph{rests on
  Assumption~\ref{ass:snapped-measurement}} when it is not evaluated
  off its snap (Definition~\ref{def:evaluated-off-snap}); a kept block and
  a block evaluated off its snap do not. The record of the run names every
  block that rests on the assumption with its cause, its group and a fix.
  The group is \emph{last resort} for a block snapped by the last resort,
  whose fix is a model that includes the noise the device has, so that
  every element of the block is positive definite, together with the fix of
  its cause; otherwise it is the group of its cause:
  \begin{enumerate}[label=(\alph*),leftmargin=2.0em]
  \item \emph{no usable reference}: no reference is handed on the block,
    because the declarer holds no reference of the family, because a
    reference fails the decisions of (i) or a property that (i) leaves to
    the constructing module is not established, or because the run declares
    that enclosed references are not used and a radius of the block is not
    zero; for a class of a producer's snap of its own emission, because the
    producer forms no companion that carries its exact construction past
    the snap, a \emph{bypassing companion}. The fix is a model of the family
    with a reference that the run uses, and for the last of the instance's
    causes an exact reference, or a run that uses enclosed references; for a
    producer's class, declared numbers on which its exact construction stays
    on the lattice, so that the class is kept;
  \item \emph{reference bracket withdrawn}: a reference is handed on the
    block, or, for a producer's class, a bypassing companion is formed, and
    a companion carrying the reference on the block, or bypassing the class,
    is evaluated, but the key length the run reports is formed from a number
    on the snapped declaration, for one of three causes. Either no evaluated
    point contributes a bound on the companion, so the run's first-order
    number is formed on the snapped declaration
    (\xref{meat-core/frank-wolfe/req:bracket-on-reference}); or the certified
    evaluation finds the two enclosures of a block the companion carries
    inconsistent, the bracket on the companion being withdrawn and the run's
    first-order number formed on the snapped declaration, the cause
    \iface{factor_enclosure_inconsistent}
    (\xref{meat-core/frank-wolfe/req:certified-evaluation}\,(iv)); or a point
    contributes, but the certified evaluation on the companion certifies no
    point and the key length is formed from a number certified on the snapped
    declaration in its place, the cause \iface{companion_not_certified}
    (\xref{meat-core/frank-wolfe/req:certified-evaluation}\,(i)). The fix is
    the one the reasons recorded at the evaluated points of the companion
    name: the withdrawal reasons for the first cause and the reasons of the
    certified evaluation for the third; for the second, no change of the
    declaration, an inconsistency being a defect of the engine in the
    modules that carry the enclosures, so the fix says that and that the
    defect be reported;
  \item \emph{reference bracket not evaluated}: a reference is handed on the
    block, or, for a producer's class, a bypassing companion is formed, but
    no companion carrying the reference on the block, or bypassing the class,
    is evaluated, because none is formed, compiled or admitted. The fix is
    the one the reason recorded for the companion names.
  \end{enumerate}
\item \emph{The strict option.} A run whose configuration declares the
  \emph{strict option} \iface{strict_snaps} refuses wherever a moved
  block would rest on Assumption~\ref{ass:snapped-measurement}, naming the
  block, its cause, its group and the fix, and reports no key length that
  rests on the assumption. It refuses as soon as the resting is decided: a
  declarer when it snaps a moved block on which it hands no reference, a
  producer when it snaps a moved class that no bypassing companion it forms
  carries past the snap, and the run, after its search and before any key
  length is reported, where a moved block or class is not evaluated off its
  snap. A strict run forms no certified number on the snapped declaration in
  place of a companion's
  (\xref{meat-core/frank-wolfe/req:certified-evaluation}\,(i)), so the
  cause \iface{companion_not_certified} of (iii)(b) does not arise in it,
  and it refuses on the other two. A strict run none of whose
  moved blocks rests on the assumption proceeds as one that does not declare
  the option. In a run that does not declare it, the run proceeds; its record
  carries, for every snap, the status of each block, kept or moved, the flag
  of every block snapped by the last resort, the blocks that rest on the
  assumption with their causes, groups and fixes, and
  Assumption~\ref{ass:snapped-measurement} among the assumptions the run
  rests on, naming those blocks, so that a report on the key rate can flag
  that the rate rests on a snap. A run none of whose blocks rests on the
  assumption does not rest on it.
\end{enumerate}
\end{digestrequirement}

% claim foundations/req:support-safe-snap
\digestlabel{req:support-safe-snap}{10.6}\pdfbookmark[2]{Requirement 10.6 Support-safe snaps of a measurement family}{digest-foundations/req:support-safe-snap}
\renewcommand*{\RestatedName}{\hyperref{contract.pdf}{}{claim-foundations/req:support-safe-snap}{Requirement~10.6}}%
\begin{digestrequirement}[Support-safe snaps of a measurement family]
Every snap $F$ at a lattice exponent $k$ of an operator family $D$ read as a
measurement (Definition~\ref{def:snap}), against a sector map with blocks
$s\in\mathsf{S}$ (Definition~\ref{def:sector-map}) and with constant $c$, by
an instance of what it declares or by a producer of its own emission, with
$D$ the data as they stand or a reference of the family
(Definition~\ref{def:exact-reference}), satisfies (i) to (iii) on every block
$s$ and for every outcome $j$, save on a block snapped by the last resort of
(iv).
\begin{enumerate}[label=(\roman*),leftmargin=2.4em]
\item \emph{No revival.} The structural support of $F^{(s)}_j$ is
  contained in that of $D^{(s)}_j$. In particular, if $D^{(s)}_j$ is
  structurally zero, then $F^{(s)}_j=0$; and for every index $i$ at
  which both row $i$ and column $i$ of $D^{(s)}_j$ are structurally zero,
  row $i$ and column $i$ of $F^{(s)}_j$ are zero, which for a reference is
  every index whose row is structurally zero.
\item \emph{Lift per coordinate.} A step of the snap that changes the
  elements in order to make them positive semidefinite does one of two
  things. Either it adds to $F^{(s)}_j$ only a matrix $\sum_it_ie_ie_i^{*}$
  with $t_i\ge0$, $e_i$ the $i$-th coordinate vector of the block, and
  $t_i>0$ only for indices $i$ in the structural support of $D^{(s)}_j$; or
  it replaces the family on the block by $(1-w)E+wT$. Here
  $w\in(0,1)$ is rational; $E$ is the family on the block as the
  preceding steps of the snap left it or, where every entry of $D^{(s)}$ is a Gaussian rational, $D^{(s)}$
  itself; and $T$ is a family complete at $c$ on the block whose element
  $T^{(s)}_j$ is, for every $j$, a diagonal matrix with non-negative entries
  that are non-zero only at indices in the structural support of
  $D^{(s)}_j$.
\item \emph{Residue.} A step of the snap that adds to an element what the sum
  of the elements lacks of $c\,\idon{d_s}$, or a part of it, adds it to
  $F^{(s)}_j$ only where the structural support of $D^{(s)}_j$ contains every
  index $i$ such that row $i$ or column $i$ of what it adds has a non-zero
  entry.
\item \emph{Last resort.} Where a procedure meeting (i) to (iii) yields on a
  block $s$ no family that satisfies the conditions of
  Definition~\ref{def:snap}, the snap may, on that block alone, replace the
  family by $(1-w)E+wU$ and then restore the lattice condition and
  completeness; (i) to (iii) do not bind on that block. Here $w\in(0,1)$ is
  rational; $E$ is the family that procedure produced on the block or, where
  every entry of $D^{(s)}$ is a Gaussian rational, $D^{(s)}$ itself; and $U$
  is a family complete at $c$ on the block whose element $U^{(s)}_j$ is, for
  every $j$, a positive rational multiple of $\idon{d_s}$, or zero where
  $D^{(s)}_j$ is structurally zero. The record of the snap flags the block
  \iface{support_not_preserved}, since an element may then have support
  where $D^{(s)}_j$ is structurally zero.
\end{enumerate}
\end{digestrequirement}

\section*{11\quad Emission}\pdfbookmark[1]{11 Emission}{digest-group-foundations/req:declared-lattice}

% claim foundations/req:declared-lattice
\digestlabel{req:declared-lattice}{11.1}\pdfbookmark[2]{Requirement 11.1 The declared lattice condition}{digest-foundations/req:declared-lattice}
\renewcommand*{\RestatedName}{\hyperref{contract.pdf}{}{claim-foundations/req:declared-lattice}{Requirement~11.1}}%
\begin{digestrequirement}[The declared lattice condition]
(D-LAT) Let $k$ be the one lattice exponent
(Definition~\ref{def:binary64}\,(iv)) at which a declarer declares operator
data. Every entry of every operator of those data is a lattice value at
$k$: the real and the imaginary part of the entry are integer multiples of
$2^{-k}$ and have modulus at most $1$. Each entry is then exactly a binary64
number, by Lemma~\ref{lem:binary64-lattice}\,(i).

Operator data are declared at a lattice exponent $k$, and (D-LAT) binds them,
in these cases and no other: the data an instance snaps at $k$
(Definition~\ref{def:snap}); the data an instance declares to a consumer that
admits declared data only at a declared lattice exponent $k$; and every
operator that a producer which declares a lattice exponent $k$ emits
(Definition~\ref{def:declared-data}). Where data of the second case are
numbers of a supplier's output that the instance passes on unchanged
(Definition~\ref{def:declared-data}\,(ii)), the instance does not establish
(D-LAT) of them: the consumer's exact decision on entry establishes it, and
the consumer refuses data that fail it
(Requirement~\ref{req:boundary-admission}). A producer that declares no
lattice exponent has no obligation under (D-LAT).
\end{digestrequirement}

% claim foundations/req:rounding-direction
\digestlabel{req:rounding-direction}{11.2}\pdfbookmark[2]{Requirement 11.2 Rounding direction of emitted bounds}{digest-foundations/req:rounding-direction}
\renewcommand*{\RestatedName}{\hyperref{contract.pdf}{}{claim-foundations/req:rounding-direction}{Requirement~11.2}}%
\begin{digestrequirement}[Rounding direction of emitted bounds]
(D-DIR) Every emitted scalar bound is rounded onto the lattice in the direction
that weakens the claim it carries, never to nearest: a lower bound \emph{down},
an upper bound \emph{up}, a constraint constant in the direction that enlarges
the feasible set for its declared relation. The direction is recorded per
bound, and where a family's completeness depends on a rounded weight the
complementary quantity is \emph{recomputed} from the rounded weights, never
rounded independently.

It binds every producer that declares a lattice exponent $k$
(Definition~\ref{def:declared-data}), for each scalar bound it emits, the
lattice being that of the lattice values at $k$
(Definition~\ref{def:binary64}\,(iv)); the claim that states the emission names
the direction of each such bound. A producer that declares no lattice exponent,
and that constructs an emitted scalar bound exactly and then rounds it to a
binary64 number, rounds it in the same direction, never to nearest.
\end{digestrequirement}

% claim foundations/req:exact-emission
\digestlabel{req:exact-emission}{11.3}\pdfbookmark[2]{Requirement 11.3 Exact verification before emission}{digest-foundations/req:exact-emission}
\renewcommand*{\RestatedName}{\hyperref{contract.pdf}{}{claim-foundations/req:exact-emission}{Requirement~11.3}}%
\begin{digestrequirement}[Exact verification before emission]
(D-EMIT) Every producer (Definition~\ref{def:declared-data}) emits an operator
family as a measurement only after establishing, of the
numbers it will emit and by exact arithmetic on them: (i) the
sector map is the
consumer's, decided and never measured; (ii) exact Hermiticity; (iii) exact
positive semidefiniteness, by a rational $LDL^{*}$ with the zero-pivot rule or
an equivalent exact decision, never by a floating-point eigenvalue; (iv) exact
completeness per sector block, as an identity of the emitted numbers; (v) the
completeness sum in the declared order rounds nowhere in
binary64, and, where
it emits Alice's measurement $\{M_x\}$ and her marginal $\tau_A$ of a protocol
declaration, so do the completeness sum of the one and the trace of the other;
(vi) where it declares a lattice exponent $k$, Requirement~\ref{req:declared-lattice}
of every entry at $k$. Conditions (i)--(v) of a family are its exact
admissibility with constant $1$ for the sector map the consumer holds
(Definition~\ref{def:exact-admissibility}), and the clause of (v) on
$\{M_x\}$ and $\tau_A$ binds only a producer that emits them.

A producer that cannot establish the conditions that bind it \emph{refuses,
naming the predicate}, save in the one case below; it repairs nothing it was
handed. A producer that declares a lattice exponent $k$ may, where the family it
constructed fails any of (ii)--(iv) or (vi), instead \emph{declare} a nearby
family established under all six: a snap at $k$ of its own emission
(Definition~\ref{def:snap}), made as Requirement~\ref{req:snap-reference}
states. Its moved blocks are related to its exact construction by no
claim; a run that consumes the emission evaluates the numbers its key length
rests on with the exact construction in their place where it can, and
otherwise rests, on the blocks that Requirement~\ref{req:snap-reference}\,(iii)
names, on Assumption~\ref{ass:snapped-measurement}\,(ii): what it concludes
from the emitted family holds of the snapped family. A run that declares the
strict option of Requirement~\ref{req:snap-reference}\,(iv) refuses there
instead. The producer must disclose that the emitted numbers are a
declaration and not the constructed ones, the measured per-block
displacement, the rounding direction of every weight that it rounds in
constructing the emission, and, separately, the status of every block, kept
or moved, and the flag of every block snapped by the last resort of
Requirement~\ref{req:support-safe-snap}\,(iv). A producer that
declares no lattice exponent declares no nearby family.
\end{digestrequirement}

\section*{12\quad Arithmetic and tiers}\pdfbookmark[1]{12 Arithmetic and tiers}{digest-group-foundations/req:exact-island}

% claim foundations/req:exact-island
\digestlabel{req:exact-island}{12.1}\pdfbookmark[2]{Requirement 12.1 Exact arithmetic on the island, unbudgeted rounding elsewhere}{digest-foundations/req:exact-island}
\renewcommand*{\RestatedName}{\hyperref{contract.pdf}{}{claim-foundations/req:exact-island}{Requirement~12.1}}%
\begin{digestrequirement}[Exact arithmetic on the island, unbudgeted rounding elsewhere]
Every decision, certification and construction of the exact island
(Definition~\ref{def:exact-island}) is computed in exact rational arithmetic on
the values it is given, so that what it decides, certifies or constructs holds
of those values exactly. A candidate proposed by binary64 arithmetic, such as
a point or a multiplier, enters the exact island only as a value given to such
a decision or certification, and nothing is established of it by the
arithmetic that proposed it. Outside the exact island the engine may compute
in binary64, and it neither bounds nor budgets the rounding of those
computations: no module other than a supplier tracks a floating-point error
budget, and no result of the engine rests on one except as follows. A
supplier (Definition~\ref{def:declared-data}\,(ii)) may allow for the error of
one of its own floating-point evaluations by a slack that an instance declares
to it, under an assumption of its own claims, which it does not verify, that
the slack bounds that error on the side that would make its output unsound.
A result that rests on such an evaluation rests on that assumption, and the
supplier carries the assumption, with its slack, in its output. A computation
outside the exact island that is rigorously enclosed
(Definition~\ref{def:tiers}) is neither: each of its operations bounds its
own rounding, so what it establishes holds of the exact quantity and no
allowance is assumed to bound its error. The engine forms every number it
reports at the certified tier from such computations and from the
exact island alone.
\end{digestrequirement}

% claim foundations/req:tiers
\digestlabel{req:tiers}{12.2}\pdfbookmark[2]{Requirement 12.2 Reporting of the tiers}{digest-foundations/req:tiers}
\renewcommand*{\RestatedName}{\hyperref{contract.pdf}{}{claim-foundations/req:tiers}{Requirement~12.2}}%
\begin{digestrequirement}[Reporting of the tiers]
Every run that reports a key rate reports it at every tier of
Definition~\ref{def:tiers} that it reaches: the estimate always, the
first-order number wherever the hypotheses of the claim that establishes it
held, and the certified number wherever the requirement of the solver that
forms it reports it at that tier, which may withhold the tier from a number
that meets Definition~\ref{def:tiers}\,(iii) but rests on an upstream
premise; together with the field
\iface{reported_tier}, which names the strongest tier reached, the tier that
carries the run's claim. Where the first-order number is reported, the number
at that strongest tier is the headline key rate, and the key length the run
reports is formed from it; a number at a weaker tier is recorded beside it.
No module reports a number under a tier stronger than the checks that number
has passed support, and a request naming a tier that
Definition~\ref{def:tiers} does not list is refused as an unknown tier.
\end{digestrequirement}

\bibliographystyle{plain}
\bibliography{../../../refs/references}

\end{document}
